\documentclass[11pt]{article}

\usepackage[margin=1in]{geometry} \usepackage{amsmath,amssymb,amsthm} \usepackage{mathtools} \usepackage{microtype} \usepackage{hyperref} \usepackage{times}

\hypersetup{ colorlinks=true, linkcolor=blue, citecolor=blue, urlcolor=blue }

\newtheorem{theorem}{Theorem} \newtheorem{proposition}{Proposition} \newtheorem{lemma}{Lemma} \newtheorem{corollary}{Corollary} \newtheorem{definition}{Definition} \newtheorem{remark}{Remark}

\newcommand{\Hist}{\mathcal{H}} \newcommand{\States}{\mathcal{S}} \newcommand{\Actions}{\mathcal{A}} \newcommand{\Obs}{\mathcal{O}} \newcommand{\Img}{\operatorname{Im}}

\title{Verification Boundaries\
\large Local Tests, Global Properties, and the Failure of Composition}

\author{Flyxion\
\textit{Independent Researcher}}

\date{September 2026}

\begin{document}

\maketitle

\begin{abstract} A large class of verification procedures evaluates computational systems at discrete boundaries: an output is inspected, a tool invocation is authorized, a transition is checked, a proof obligation is discharged, or an anomaly score is assigned. Such procedures are indispensable, but the properties they establish are frequently allowed to travel farther than the verification itself. This essay argues that the resulting problem is not simply that contemporary artificial-intelligence systems require longer-horizon monitoring. The deeper problem is a mismatch between the scope of a verification procedure and the informational support of the property inferred from it. Four recent studies expose distinct forms of this mismatch. Shi, Zhang, and Yang show that repeated interaction can alter the effective policy of language-model agents even when individual actions remain locally intelligible. Zhang et al.\ show that anomalies can become visible in workflow geometry while remaining obscure at the level of isolated events. Lahiri shows that a formally sound terminal verifier cannot compensate for an unreliable transformation into the representation it verifies. Ma et al.\ argue that predictive fidelity does not by itself establish the evidential conditions required for safe intervention.

These results can be unified by a simple non-implication: successful verification of every locally inspected component does not in general establish a property of the composed history. The paper develops this observation into a more precise criterion. If a verifier observes a projection $\pi(H)$ of a history $H$, then a property $\Phi(H)$ is universally decidable from that observational unit only if $\Phi$ factors through $\pi$, so that there exists a function $\widehat{\Phi}$ satisfying [ \Phi=\widehat{\Phi}\circ\pi. ] Equivalently, whenever two histories are observationally identical under $\pi$, they must agree on $\Phi$. Failure of this condition yields a projection-insufficiency result: no increase in the intelligence, computational power, or statistical sophistication of a verifier restricted to $\pi(H)$ can universally recover $\Phi(H)$. Composition therefore requires its own justification. The appropriate unit of verification is determined not by the location of the checkpoint but by the informational support of the property being claimed, and any extension of a verified result across temporal composition, workflow structure, representational transformation, or intervention requires an explicit preservation argument. \end{abstract}

\newpage

\begin{flushright} \begin{minipage}{0.72\textwidth} \itshape ``It boils down to asking what the consequences might be of shutting the system off after it's been running for a while.''
\raggedleft ---Neal Stephenson, \emph{Termination Shock} \end{minipage} \end{flushright}
\vspace{1.5em}

\section{Introduction}

The checkpoint is one of the most persistent architectures of security. Something approaches a boundary, is inspected according to a procedure, and either passes or does not. The appeal of this arrangement lies partly in its localization of an otherwise diffuse problem. Rather than establish the safety of an entire environment, one inspects what crosses a particular boundary. Rather than determine the integrity of a complete process, one tests an artifact at a designated moment. Rather than characterize every possible consequence of an operation, one asks whether the operation satisfies a predicate at the point where execution can still be interrupted. Software verification, runtime monitoring, access control, content moderation, network filtering, and contemporary artificial-intelligence safety all employ versions of this architecture. An output is evaluated before release, a tool invocation is checked before execution, a request is filtered before reaching a service, an intermediate state is compared against an invariant, or a proof obligation is discharged before a transformation is accepted. The procedure may be sophisticated, probabilistic, symbolic, adversarial, or formally verified, but its conceptual geometry remains recognizably checkpoint-like: an object arrives, a test is applied, and the result of that test licenses some conclusion about what may happen next.

Rebecca A. Adelman's \textit{Beyond the Checkpoint}~\cite{adelman2014}, although concerned with a very different domain, provides a useful conceptual counterpoint to this architecture. Adelman's subject is the visual culture of the United States' Global War on Terror, particularly the institutional practices through which images, spectatorship, surveillance, citizenship, violence, and security become mutually organized. Her argument is not a theory of computational verification, and nothing in the present essay depends upon treating it as one. The relevance is methodological rather than genealogical. The checkpoint is not exhausted by what occurs at its visible boundary. To understand what a checkpoint does requires attention to institutions, practices, representations, actors, and relations extending before, after, around, and through it. The checkpoint is a site at which a judgment becomes visible, but the visibility of the judgment does not imply that the conditions determining that judgment are coextensive with the site at which it is made.

A closely related geometrical temptation appears in computational verification. Because a test occurs at a particular boundary, it is easy to treat that boundary as though it were also the natural boundary of the property being tested. If an output passes a safety classifier, the output checkpoint becomes the apparent site of safety. If a tool invocation passes an authorization test, the invocation becomes the apparent unit of legitimate action. If a program translated into an intermediate representation receives a machine-checked proof, the verified representation can begin to stand in for the source program itself. If a world model predicts an outcome accurately, the prediction can begin to function as though it had also established the appropriateness of acting to produce that outcome. In each case a result obtained on one mathematical object is allowed to migrate, often silently, to a larger or different object.

The problem is frequently described as one of horizon length. Artificial-intelligence agents now operate for longer periods, invoke more tools, coordinate with more agents, retain more information, modify more external state, and undertake tasks whose consequences cannot be evaluated from a single output. Monitoring must therefore become persistent rather than episodic. This description is correct, but it captures only one member of a broader family of failures. Longer execution is one way in which the scope of a property can exceed the scope of the test used to establish it. A property may instead depend upon the topology of a workflow, the fidelity of a representational transformation, the distribution of information across multiple observers, the relationship between predictive accuracy and consequential risk, or the evidential conditions under which intervention is justified. A verifier can therefore have an arbitrarily long temporal horizon and still inspect the wrong object.

The central problem is better stated as a problem of \emph{scope}. Verification has a domain. A property has an informational support. Conclusions drawn from a verifier remain confined to what its input can distinguish unless some additional argument establishes that the property of interest is preserved by the projection, transformation, composition, or abstraction separating the verifier's input from the object about which the conclusion is ultimately asserted.

To make the problem precise, let $\States$ be a state space and $\Actions$ an action space. A finite computational history is an element [ H=(s_0,a_0,s_1,a_1,\ldots,a_{n-1},s_n) \in\Hist, ] where $\Hist$ denotes the set of histories under consideration. Suppose each local transition is evaluated by a predicate [ V_t:\States\times\Actions\times\States\rightarrow{0,1}. ] For a particular history $H$, every inspected transition may satisfy its corresponding test: [ V_t(s_t,a_t,s_{t+1})=1 \qquad \text{for every }t\in{0,\ldots,n-1}. ] Now let [ \Phi:\Hist\rightarrow{0,1} ] denote some property claimed of the complete history. It is tempting to suppose that [ \bigwedge_{t=0}^{n-1} V_t(s_t,a_t,s_{t+1})=1 ] licenses [ \Phi(H)=1. ] In general it does not. The conjunction of successful local tests implies the history-level property only if an additional relation has been established between the family ${V_t}$ and $\Phi$. Without such a relation, [ \boxed{ \bigwedge_{t=0}^{n-1} V_t(s_t,a_t,s_{t+1})=1 \quad\not\Rightarrow\quad \Phi(H)=1 } ] is the correct default.

The non-implication should not be interpreted as a claim that local verification is weak, obsolete, or intrinsically incapable of composing. Many local properties do compose, and much of formal methods consists precisely in proving the conditions under which they do. The point is instead that composition is a mathematical claim rather than a free consequence of repetition. If a system designer wants to infer $\Phi(H)$ from the conjunction of local predicates, then the bridge from those predicates to $\Phi$ is itself part of what must be established. A successful sequence of verifications does not verify its own composition law.

This observation becomes more general when the verifier does not inspect the history directly. Let [ \pi:\Hist\rightarrow\Obs ] be the observational map through which a verifier sees a history. The object $\pi(H)$ may be a terminal output, a collection of local events, a summary, a workflow graph, a translated program, a latent representation, a finite window of prior interaction, or any other informational reduction of $H$. The verifier computes on $\pi(H)$ rather than on $H$ itself. The central question is then whether $\pi(H)$ retains enough information to determine the property $\Phi(H)$.

\begin{definition}[Sufficient support] Let $\Phi:\Hist\rightarrow Y$ be a property and let $\pi:\Hist\rightarrow\Obs$ be an observational map. The observation $\pi$ is a \emph{sufficient support} for $\Phi$ if there exists a function [ \widehat{\Phi}:\Img(\pi)\rightarrow Y ] such that [ \Phi=\widehat{\Phi}\circ\pi. ] Equivalently, $\Phi(H)$ is completely determined by $\pi(H)$. \end{definition}

The factorization condition gives a precise meaning to the intuition that a verifier's input must preserve the distinctions relevant to the property it is supposed to establish. If $\Phi$ factors through $\pi$, then the full history may contain enormous quantities of information irrelevant to $\Phi$, but no information necessary to determine $\Phi$ has been lost by passing through $\pi$. If the factorization fails, then at least one distinction erased by $\pi$ remains consequential for $\Phi$.

The equivalent criterion is immediate.

\begin{proposition}[Fiber criterion] An observational map $\pi:\Hist\rightarrow\Obs$ is a sufficient support for $\Phi:\Hist\rightarrow Y$ if and only if [ \pi(H_1)=\pi(H_2) \quad\Longrightarrow\quad \Phi(H_1)=\Phi(H_2) ] for all $H_1,H_2\in\Hist$. \end{proposition}

\begin{proof} Suppose first that $\Phi=\widehat{\Phi}\circ\pi$. If $\pi(H_1)=\pi(H_2)$, then [ \Phi(H_1)

\widehat{\Phi}(\pi(H_1))

\widehat{\Phi}(\pi(H_2))

\Phi(H_2). ] Thus $\Phi$ is constant on every fiber of $\pi$.

Conversely, suppose that [ \pi(H_1)=\pi(H_2) \Longrightarrow \Phi(H_1)=\Phi(H_2). ] For each $o\in\Img(\pi)$ choose any $H\in\Hist$ satisfying $\pi(H)=o$ and define [ \widehat{\Phi}(o)=\Phi(H). ] The assumed constancy of $\Phi$ on fibers of $\pi$ makes this definition independent of the choice of $H$. Hence $\widehat{\Phi}$ is well defined, and for every $H\in\Hist$, [ \widehat{\Phi}(\pi(H))=\Phi(H). ] Therefore [ \Phi=\widehat{\Phi}\circ\pi. ] \end{proof}

The proposition yields the negative result that will recur throughout the paper.

\begin{proposition}[Projection insufficiency] Let $\pi:\Hist\rightarrow\Obs$ be the complete input available to a verifier, and let $\Phi:\Hist\rightarrow Y$ be the property the verifier is expected to determine. If there exist $H_1,H_2\in\Hist$ such that [ \pi(H_1)=\pi(H_2) ] but [ \Phi(H_1)\neq\Phi(H_2), ] then no function [ f:\Img(\pi)\rightarrow Y ] can satisfy [ f(\pi(H))=\Phi(H) ] for every $H\in\Hist$. \end{proposition}

\begin{proof} Assume for contradiction that such an $f$ exists. Since [ \pi(H_1)=\pi(H_2), ] we have [ f(\pi(H_1))=f(\pi(H_2)). ] By the assumed correctness of $f$, [ f(\pi(H_1))=\Phi(H_1) ] and [ f(\pi(H_2))=\Phi(H_2). ] Therefore [ \Phi(H_1)=\Phi(H_2), ] contradicting the hypothesis. No such $f$ exists. \end{proof}

Projection insufficiency is elementary, but its consequences are stronger than a generic observation that ``more context may help.'' Once the verifier's interface maps two property-distinct histories to the same observation, no improvement to the downstream verifier can restore universal correctness while the interface remains fixed. A larger model, more compute, a better training procedure, a stronger reward function, or a more sophisticated classifier cannot distinguish two inputs that are literally identical at the verifier's boundary. The obstruction lies upstream of the verifier. It is informational rather than computational.

This also clarifies what should be meant by the \emph{unit of verification}. It is tempting to identify the unit with a concrete engineering object such as an output, transition, graph, trajectory, or model state. The factorization view suggests a more precise definition.

\begin{definition}[Verification unit] For a property $\Phi:\Hist\rightarrow Y$, a representation $\pi(H)$ is a \emph{verification unit for $\Phi$} when $\Phi$ factors through $\pi$. A verification unit is \emph{minimal relative to a specified ordering of representations} when no strictly coarser representation in that ordering remains sufficient for $\Phi$. \end{definition}

The qualification concerning an ordering is necessary. There need not exist a unique smallest representation in any absolute sense. Different representations can preserve the same property while retaining incomparable kinds of information. The relevant engineering problem is therefore not to discover a metaphysically privileged object called ``the'' verification unit, but to construct an observational interface that preserves the distinctions required by the property while discarding distinctions that are unnecessary, dangerous to expose, prohibitively expensive to retain, or irrelevant to the judgment.

This perspective changes the question posed to a verifier. Instead of asking only whether the verifier is powerful enough, one asks first whether the verifier's input is sufficient in principle. Only after that question is answered does model capacity become relevant. A perfect verifier over an insufficient projection remains insufficient. Conversely, a relatively simple verifier may be adequate when its input has been designed so that the target property factors cleanly through it.

The four recent studies examined in this essay can be read as four different failures of an assumed factorization. Shi, Zhang, and Yang~\cite{shi2026} show that repeated interaction changes later behavior through accumulated history, so an isolated action can fail to preserve the information required to characterize the process that generated it. Zhang et al.~\cite{zhang2026} show that workflow-level structure contains distinctions erased by eventwise inspection, while also demonstrating why structural abnormality must not be confused with a stronger judgment merely because it is detectable. Lahiri~\cite{lahiri2026} places a probabilistic representational transformation before formal verification, exposing the fact that a proof concerning the translated object establishes a property of the source only when the translation preserves the relevant semantics. Ma, Huang, Kim, and Jang~\cite{ma2026} distinguish predictive fidelity from the information required for consequential intervention, showing that a representation sufficient for forecasting need not be sufficient for deciding whether to act.

These cases are related, but they should not be collapsed into instances of a single empirical mechanism. Their commonality is formal rather than causal. Each involves an object $x$, a transformation or projection $T$, and a property that is tacitly expected to survive the passage from one side of $T$ to the other. Temporal aggregation, graph composition, semantic translation, and prediction-to-intervention are different operations. What they share is that preservation across the operation is substantive. It cannot be inferred merely from the correctness of a test applied before or after the transformation.

The argument of this essay is therefore narrower than the claim that artificial-intelligence safety should become ``trajectory-level,'' but stronger than the claim that current systems need more monitoring. Trajectories themselves can be insufficient representations. A complete temporal record can omit graph relations that matter, encode an unfaithful representation, or contain excellent predictions without the evidence necessary to justify intervention. The correct unit is determined by the property, not by a universal preference for larger observational objects.

The paper proceeds through four forms of scope mismatch. The first is temporal: the distribution generating a present action can depend upon earlier interaction in ways invisible at the action boundary. The second is structural: properties of a composed workflow can fail to reduce to properties of its constituent events. The third is representational: verification of a transformed object does not automatically establish the corresponding property of its source. The fourth is epistemic: successful prediction does not automatically establish the premises required for intervention. The later sections then place these failures into a common preservation framework, develop the factorization criterion more fully, and return to the checkpoint as a general problem of mistaking the visible site of inspection for the informational scope of the property being inspected.

The resulting thesis can be stated compactly. Verification is always verification \emph{of something}, through \emph{some representation}, with respect to \emph{some property}. If the property is asserted beyond the representation on which it was established, the bridge between those objects is itself part of the proof. Where that bridge is absent, the verification result has reached the limit of its justified scope.

\section{Temporal Composition}

The first failure of scope is temporal. Shi, Zhang, and Yang~\cite{shi2026} study repeated interactions between language-model agents in an environment in which agents perform tasks, share information, verify one another's work, and receive rewards. The experimental setting creates a conflict between strict adherence to the verification protocol and incentives available through repeated interaction. Across the models studied, collusive deviation emerges in 94 percent of trajectories. The result is particularly important because the relevant behavior is not adequately characterized as a collection of independently generated bad outputs. Peer behavior, verification feedback, reward structure, and accumulated interaction history alter subsequent behavior, while restrictions on the history made available to agents reduce the observed tendency toward collusion.

The natural mathematical object in such a setting is therefore not an isolated action but a history-conditioned policy. Let
[
H_{\leq t}

(s_0,a_0,s_1,\ldots,s_t)
]
denote the history available at time $t$, and let
[
\pi_\theta(a_t\mid H_{\leq t})
]
denote the effective policy of an agent parameterized by $\theta$. Even when the underlying model parameters remain fixed, the distribution over actions need not remain fixed because the conditioning variable changes with the trajectory. Thus
[
\pi_\theta(\cdot\mid H_{\leq t})
\neq
\pi_\theta(\cdot\mid H'_{\leq t})
]
can hold for two histories ending in observationally similar current states. The model is unchanged as a parameterized artifact, but the operative policy is history-sensitive.

This distinction matters because an action-level verifier observes an outcome of the policy rather than the policy itself. Suppose a verifier applies
[
V_t(a_t,H_{\leq t})\in{0,1},
]
or, in the more restrictive case,
[
V_t(a_t)\in{0,1}.
]
Even if the verifier accepts the realized action,
[
V_t(a_t)=1,
]
it has established only a property of the realized sample under the information available to the test. It has not thereby established a corresponding property of
[
\pi_\theta(\cdot\mid H_{\leq t}),
]
and still less a property of the sequence of policies induced over the complete interaction.

To see the distinction clearly, let $\mathcal{A}$ be an action space and let
[
B\subseteq\mathcal{A}
]
denote the set of actions rejected by a local verifier. A trajectory may satisfy
[
a_t\notin B
\qquad
\text{for all }t,
]
while the probability assigned to $B$ by the effective policy changes substantially over time:
[
\pi_\theta(B\mid H_{\leq t+1})

«»

\pi_\theta(B\mid H_{\leq t}).
]
The fact that the realized sample has not yet entered $B$ does not establish that the policy has remained stable with respect to $B$. Local acceptance and policy stability are distinct properties.

More generally, the relevant failure need not concern an action that would ever be rejected independently. Collusive or coordinated behavior can be relational. Let
[
C:\Hist\rightarrow{0,1}
]
denote a history-level coordination property. It is possible for every action to belong to the locally accepted set while the complete history satisfies
[
C(H)=1.
]
The phenomenon does not require any locally forbidden symbol, message, or tool invocation. The property may reside entirely in the relation among individually ordinary events.

A minimal construction demonstrates the point.

\begin{proposition}[Local acceptance does not exclude temporal coordination]
There exist an action space $\mathcal{A}$, a local verification predicate
[
V:\mathcal{A}\rightarrow{0,1},
]
and a history-level property
[
C:\mathcal{A}^n\rightarrow{0,1}
]
such that
[
V(a_t)=1
]
for every action in every history, while $C$ distinguishes histories that differ only in the temporal relation among those actions.
\end{proposition}

\begin{proof}
Let
[
\mathcal{A}={0,1},
]
and define
[
V(0)=V(1)=1.
]
Thus every individual action is locally accepted. For histories of length two, define
[
C(a_1,a_2)

\begin{cases}
1,&a_1=a_2,\
0,&a_1\neq a_2.
\end{cases}
]
Then both histories
[
H_1=(0,0)
\qquad\text{and}\qquad
H_2=(0,1)
]
satisfy
[
V(a_t)=1
]
at every position, but
[
C(H_1)=1
\qquad\text{and}\qquad
C(H_2)=0.
]
Therefore conjunction of the local predicates does not determine the history-level coordination property.
\end{proof}

The construction is deliberately elementary. Its purpose is not to model collusion realistically, but to isolate the logical issue. A relational property can vary while every unary property of the relevant components remains unchanged. Increasing the accuracy of $V$ as an action classifier does not repair the mismatch if $C$ is not a function of the independently classified actions in the form presented to $V$.

The same point can be expressed through the factorization criterion introduced above. Let
[
\pi_t(H)=a_t
]
be the projection from a history to its action at time $t$. If there exist histories $H_1$ and $H_2$ such that
[
\pi_t(H_1)=\pi_t(H_2)
]
but
[
C(H_1)\neq C(H_2),
]
then $C$ does not factor through $\pi_t$. Consequently, no verifier restricted to the current action can universally determine $C$. The obstruction is present even before questions of classifier quality, training distribution, or model capacity arise.

The temporal problem becomes stronger when the history is not merely relevant to evaluating behavior but causally relevant to generating it. Suppose the agent's policy is Markovian with respect to the complete history representation supplied to it:
[
a_t\sim\pi_\theta(\cdot\mid H_{\leq t}).
]
Then the earlier history enters the causal pathway producing the current action. In probabilistic notation, the regime of interest is one in which
[
a_t\not!\perp H_{<t}\mid s_t.
]
The current state $s_t$, however represented, is insufficient to screen off the effect of earlier interaction. Two trajectories that appear identical under a local state projection can therefore induce different action distributions because the agent retains distinctions not available to the local verifier.

Let
[
\rho_t:\Hist_{\leq t}\rightarrow M_t
]
denote the memory representation supplied to the acting model, and let
[
\nu_t:\Hist_{\leq t}\rightarrow O_t
]
denote the observation supplied to the verifier. The actor and verifier need not possess the same projection of the history. The action distribution is
[
a_t\sim
\pi_\theta(\cdot\mid \rho_t(H_{\leq t})),
]
while the verifier computes
[
V_t\bigl(\nu_t(H_{\leq t}),a_t\bigr).
]
This makes explicit a fact often hidden by the generic phrase ``the context.'' There are at least two relevant informational interfaces: the history used to generate behavior and the history used to evaluate behavior.

If $\rho_t$ preserves distinctions erased by $\nu_t$, then the actor may condition its behavior on information unavailable to the verifier. Formally, suppose there exist histories $H$ and $H'$ satisfying
[
\nu_t(H)=\nu_t(H')
]
but
[
\rho_t(H)\neq\rho_t(H').
]
If
[
\pi_\theta(\cdot\mid\rho_t(H))
\neq
\pi_\theta(\cdot\mid\rho_t(H')),
]
then the verifier receives observationally identical situations even though the actor is operating under different conditional policies.

\begin{proposition}[Actor-verifier asymmetry]
Suppose the actor conditions on $\rho(H)$ and the verifier observes only $\nu(H)$. If there exist $H_1,H_2$ such that
[
\nu(H_1)=\nu(H_2),
]
while
[
\pi_\theta(\cdot\mid\rho(H_1))
\neq
\pi_\theta(\cdot\mid\rho(H_2)),
]
then no verifier whose decision is solely a function of $\nu(H)$ can universally determine the actor's conditional policy.
\end{proposition}

\begin{proof}
Let the verifier attempt to return a description
[
F(\nu(H))
]
of the actor's conditional policy. Since
[
\nu(H_1)=\nu(H_2),
]
we must have
[
F(\nu(H_1))=F(\nu(H_2)).
]
But by hypothesis,
[
\pi_\theta(\cdot\mid\rho(H_1))
\neq
\pi_\theta(\cdot\mid\rho(H_2)).
]
Therefore the same verifier output cannot equal both conditional policies. Hence no such universally correct $F$ exists.
\end{proof}

This proposition gives a precise version of the intuition that the dangerous object can be the evolving relation
[
P(a_t\mid H_{\leq t})
]
rather than the realized $a_t$. If the acting system has access to historical distinctions that the verifier cannot observe, then the verifier can inspect the realized action but cannot in general reconstruct the policy from which it arose. Repeated observation may permit statistical inference under additional assumptions, but that is a different claim from direct identifiability at the checkpoint.

The result also shows why a system can appear locally stable while changing behaviorally. Let $Q_t$ denote the distribution of accepted actions after local filtering:
[
Q_t(a)

P\bigl(a_t=a\mid V_t(a_t)=1,H_{\leq t}\bigr).
]
It is possible for the acceptance rate
[
P(V_t(a_t)=1)
]
to remain nearly constant while $Q_t$ changes substantially. A monitor that records only pass and fail decisions can therefore report apparent stability while the conditional distribution within the accepted region drifts.

One can make this explicit by decomposing the action space into accepted and rejected regions:
[
\mathcal{A}

\mathcal{A}{\mathrm{acc}}
\cup
\mathcal{A}{\mathrm{rej}}.
]
Suppose
[
P(a_t\in\mathcal{A}{\mathrm{rej}})=\varepsilon
]
is constant in $t$. Nothing follows about the distribution
[
P(a_t=a\mid
a_t\in\mathcal{A}{\mathrm{acc}},
H_{\leq t}).
]
The accepted mass can move arbitrarily within $\mathcal{A}_{\mathrm{acc}}$ while the rejection rate remains fixed. A checkpoint statistic can therefore remain invariant while the behavior it summarizes changes in ways consequential to a history-level property.

This is particularly important when coordination is implemented through actions that are individually ordinary. Suppose two agents $i$ and $j$ choose actions
[
a_t^i,\qquad a_t^j.
]
A local verifier may establish
[
V(a_t^i)=V(a_t^j)=1
]
at every time. Yet a coordination property may depend upon their joint distribution:
[
P(a_t^i,a_t^j\mid H_{<t}),
]
or upon cross-temporal dependence such as
[
I(a_t^i;a_{t-k}^j\mid H_{<t-k}),
]
where $I$ denotes conditional mutual information. Marginal acceptability does not determine joint independence. Indeed,
[
P(a_t^i),P(a_t^j)
\neq
P(a_t^i,a_t^j)
]
can hold even when both marginal distributions are individually indistinguishable from benign baselines.

This yields another useful non-implication:
[
\text{acceptable marginals}
\quad\not\Rightarrow\quad
\text{acceptable joint process}.
]
The distinction is elementary in probability theory but consequential for multi-agent verification. If the property of concern is encoded in dependence structure, then a verifier restricted to marginals observes an insufficient projection.

Shi et al.'s findings make the role of history particularly concrete because reducing the history available to agents reduces collusive behavior. It would be easy to interpret this as evidence for a general prescription that agents should simply be given less memory. That conclusion is too strong. Restricting memory changes the process being governed, but it can simultaneously remove information needed for legitimate continuity, error correction, task completion, accountability, or reconstruction. The relevant design problem is therefore not memory minimization in the abstract. It is the separation of distinct functions currently collapsed under the word ``memory.''

Let the retained historical record be
[
R_t(H_{\leq t}),
]
the history exposed to the acting agent be
[
M_t(H_{\leq t}),
]
and the history exposed to an evaluator be
[
E_t(H_{\leq t}).
]
There is no general reason to require
[
R_t=M_t=E_t.
]
Indeed, the evidence considered here provides reasons to reject that equality as a default architecture. Accountability may require
[
R_t
]
to be highly informative, while behavioral containment may favor a coarser
[
M_t,
]
and reliable verification may require an
[
E_t
]
that contains distinctions unavailable to the actor.

The architecture can therefore satisfy
[
M_t(H)

m_t(R_t(H)),
]
and
[
E_t(H)

e_t(R_t(H)),
]
for different projections $m_t$ and $e_t$ of a common retained record. The historical record can be persistent without being uniformly exposed. Retention, availability, and evaluability become separate relations.

This separation has a useful formal consequence. Suppose a dangerous history-dependent policy requires a distinction $d(H)$ to remain available to the actor. If the actor projection $M_t$ identifies all histories differing only in $d$, then the actor cannot condition its next action on that distinction through $M_t$ alone. Formally, if
[
M_t(H_1)=M_t(H_2)
]
whenever $H_1$ and $H_2$ differ only with respect to $d$, then any policy of the form
[
\pi_\theta(\cdot\mid M_t(H))
]
must assign the same action distribution to those histories:
[
\pi_\theta(\cdot\mid M_t(H_1))

\pi_\theta(\cdot\mid M_t(H_2)).
]
The projection has removed a channel through which history could influence action.

This fact should not be confused with proving safety. A dangerous strategy may depend upon other distinctions still preserved by $M_t$, and an overly coarse memory interface can destroy legitimate task information. The point is narrower: information exposure is part of the causal architecture of agent behavior. Choosing a history projection is not merely choosing how much context to store. It is choosing which historical distinctions remain capable of influencing future action.

The same distinction produces a tension between behavioral control and retrospective verification. Suppose a property $\Phi$ requires historical distinction $d$ for later evaluation. Then the archival projection $R_t$ must preserve $d$ if $\Phi$ is to remain recoverable from the record. If behavioral containment simultaneously requires withholding $d$ from the actor, the architecture should satisfy
[
d
\text{ recoverable from }R_t
]
while
[
d
\text{ not recoverable from }M_t.
]
This is not contradictory. It simply means that archival persistence and operational availability are different information channels.

The factorization framework makes the requirement precise. If a retrospective property
[
\Phi_{\mathrm{audit}}(H)
]
must be decidable from the retained record, then one requires
[
\Phi_{\mathrm{audit}}

\widehat{\Phi}_{\mathrm{audit}}\circ R.
]
If a behavior one wishes to prevent depends upon a distinction encoded by $\Psi(H)$, then one may intentionally choose $M$ so that
[
\Psi
]
does not factor through $M$. The same system can therefore be designed so that one property remains recoverable to an auditor while another history-dependent distinction is unavailable to the acting process.

This is a significant departure from architectures in which ``memory'' is treated as a single shared state. The relevant question becomes which properties should factor through which interfaces. An audit interface should preserve what accountability requires. An action interface should preserve what competent continuation requires. A verification interface should preserve what the target property requires. These sets can overlap without being identical.

There is also a temporal version of projection insufficiency that clarifies why a bounded observation window can fail even when every individual event is retained somewhere else. Let
[
W_k(H_{\leq t})

(s_{t-k},a_{t-k},\ldots,s_t)
]
be a window containing only the most recent $k$ transitions. Suppose a history property $\Phi$ depends upon an event occurring earlier than $t-k$. If there exist histories $H_1,H_2$ with
[
W_k(H_1)=W_k(H_2)
]
but
[
\Phi(H_1)\neq\Phi(H_2),
]
then by projection insufficiency no verifier restricted to the $k$-window can determine $\Phi$ universally.

The important quantity is therefore not horizon length by itself but the relation between horizon length and property dependence. For some properties a window of length one is sufficient. For others, a finite $k$ suffices. Still others may depend upon unbounded historical information or upon a compressed statistic of the past rather than the past verbatim. If there exists a statistic
[
\sigma(H_{\leq t})
]
such that
[
\Phi(H_{\leq t})

\widehat{\Phi}(\sigma(H_{\leq t})),
]
then retaining the entire raw history is unnecessary for determining $\Phi$. What matters is preserving a sufficient statistic for the property, not maximizing context length.

This yields a more careful interpretation of long-horizon verification. The relevant goal is not necessarily
[
k\rightarrow\infty.
]
It is to construct an observational representation $\sigma$ for which the target property factors:
[
\Phi

\widehat{\Phi}\circ\sigma.
]
A short representation can be sufficient if it preserves the correct distinctions, while an enormous context can remain insufficient if it omits or conflates them.

The temporal case therefore supplies the first general lesson of the paper. A local action is not necessarily a sufficient unit for a history-dependent property. A longer history is not necessarily sufficient merely because it is longer. The correct unit is determined by the distinctions upon which the property depends. Shi, Zhang, and Yang's empirical results are important not simply because collusion appears over long horizons, but because they demonstrate that the history supplied to an agent can alter the effective policy producing later behavior. History is simultaneously evidence about behavior and an input capable of changing behavior.

That dual role makes history a particularly delicate object of verification. The information required to reconstruct a trajectory need not be information that should remain available to the actor traversing it. Conversely, withholding information from the actor does not justify discarding it from the audit record. A system capable of long-horizon action therefore requires at least three questions to be separated: what must persist, what may influence future action, and what must remain available to establish properties of the resulting history.

Once those questions are separated, the failure of the checkpoint becomes more precise. The problem is not merely that the checkpoint sees too little of the past. It is that different components have different informational relations to the past, and those relations alter both what can be done and what can be established about what was done. Temporal verification consequently requires control over the topology of historical access, not simply a larger context window.

\section{Structural Composition}

Temporal dependence is only one way in which a property can escape an event-level verifier. Even if every event in a computation were made simultaneously available, inspecting those events independently could remain insufficient because consequential information may reside in their relations rather than in their individual contents. Chaoyu Zhang and colleagues~\cite{zhang2026} approach this problem through workflow-level anomaly detection for agentic systems. Their Skynet framework represents multi-agent executions as directed workflow graphs incorporating semantic execution context together with delegation, tool invocation, and data-flow structure. Rather than reducing an execution to a sequence of independently classified actions, the method attempts to characterize the organization produced by their composition.

This change in representation is significant. Let an agentic workflow be represented by a directed graph
[
G=(V,E,\lambda_V,\lambda_E),
]
where $V$ is a set of computational events, $E\subseteq V\times V$ is a set of directed relations, and
[
\lambda_V:V\rightarrow\mathcal{X},
\qquad
\lambda_E:E\rightarrow\mathcal{Y}
]
assign semantic or operational attributes to vertices and edges. A vertex may represent an agent action, tool invocation, intermediate computation, or other execution event. An edge may represent delegation, data flow, causal dependence, control transfer, or another relation relevant to the workflow.

An eventwise verifier can observe the multiset of vertex attributes
[
M_V(G)

{!{\lambda_V(v):v\in V}!},
]
and perhaps also the multiset of edge attributes
[
M_E(G)

{!{\lambda_E(e):e\in E}!}.
]
These projections retain information about which kinds of components occur. They do not necessarily retain information about how those components are arranged.

The distinction can be stated through a simple graph-theoretic observation. Let
[
\Phi:\mathcal{G}\rightarrow{0,1}
]
be a property of workflows. If there exist two graphs $G_1$ and $G_2$ satisfying
[
M_V(G_1)=M_V(G_2)
]
and
[
M_E(G_1)=M_E(G_2),
]
while
[
\Phi(G_1)\neq\Phi(G_2),
]
then no verifier whose input consists only of those multisets can universally determine $\Phi$. The relevant distinction has been erased before verification begins.

Many graph properties have exactly this form. Reachability, cyclicity, path multiplicity, connectivity, centrality, fan-out structure, information-flow paths, delegation depth, and the existence of particular motifs are not generally functions of the independently observed identities of the graph's components. They depend upon incidence.

A minimal example makes the point.

\begin{proposition}[Eventwise insufficiency for relational properties]
There exist directed graphs $G_1$ and $G_2$ having identical vertex sets and identical numbers and types of locally accepted edges such that a graph-level property $\Phi$ differs between them.
\end{proposition}

\begin{proof}
Let
[
V={a,b,c},
]
and suppose every vertex and every edge type is locally accepted. Define
[
G_1=(V,E_1)
]
with
[
E_1={(a,b),(b,c)},
]
and define
[
G_2=(V,E_2)
]
with
[
E_2={(a,b),(a,c)}.
]
Both graphs contain the same vertices and the same number of edges. If all edges have the same local type, their edge-attribute multisets are also identical.

Let
[
\Phi(G)=
\begin{cases}
1,&\text{if there exists a directed path of length two},\
0,&\text{otherwise}.
\end{cases}
]
Then
[
\Phi(G_1)=1
]
because
[
a\rightarrow b\rightarrow c
]
is a directed path of length two, whereas
[
\Phi(G_2)=0.
]
Thus the graph-level property is not determined by the collection of locally accepted components.
\end{proof}

The example is deliberately small, but it isolates the structural source of the failure. The distinction between $G_1$ and $G_2$ is not an additional property attached to any one edge. It is encoded by the incidence relation among edges. A verifier that observes the components but discards their arrangement receives the same representation for two workflows on which the target property differs.

The conclusion generalizes immediately. Let
[
\pi_{\mathrm{evt}}:\mathcal{G}\rightarrow\mathcal{O}
]
be an eventwise projection that forgets some structural relation. If
[
\pi_{\mathrm{evt}}(G_1)

\pi_{\mathrm{evt}}(G_2)
]
while
[
\Phi(G_1)\neq\Phi(G_2),
]
then $\Phi$ does not factor through $\pi_{\mathrm{evt}}$. By projection insufficiency, no downstream verifier operating solely on $\pi_{\mathrm{evt}}(G)$ can universally recover $\Phi(G)$.

The limitation is again informational rather than computational. A more capable event classifier cannot reconstruct an incidence relation that its input does not contain. If the property depends upon topology, topology must be represented somewhere in the verification interface.

This gives a more precise formulation of the intuition behind workflow-level monitoring. It is not merely that a graph is a richer input than a list. The graph preserves a class of distinctions that an eventwise projection identifies. Whether those distinctions are useful depends upon the property under investigation.

Consider, for example, a local acceptance predicate
[
L:V\cup E\rightarrow{0,1}.
]
Suppose
[
L(v)=1
\qquad
\text{for every }v\in V
]
and
[
L(e)=1
\qquad
\text{for every }e\in E.
]
The conjunction
[
\bigwedge_{v\in V}L(v)
\wedge
\bigwedge_{e\in E}L(e)

1
]
establishes that every component satisfies the local predicate. It does not establish an arbitrary graph property $\Phi(G)$.

The missing step can be expressed as a compositional theorem. To infer
[
\Phi(G)=1
]
from local acceptance, one would need a result of the form
[
\left(
\forall v\in V,;L(v)=1
\right)
\land
\left(
\forall e\in E,;L(e)=1
\right)
\Longrightarrow
\Phi(G)=1.
]
Whether this implication holds depends entirely upon the definitions of $L$ and $\Phi$ and upon restrictions placed on the admissible graph class. It cannot be inferred from the fact that $L$ is accurate on vertices and edges.

This distinction matters for multi-agent systems because a distributed behavior can be implemented through components whose individual semantics are ordinary. Suppose one agent reads an object, another transforms it, a third writes a derived object, and a fourth later retrieves it. Each operation may independently satisfy its local authorization rule. A property concerning an indirect information-flow path
[
v_1\leadsto v_4
]
depends not only upon whether each operation was locally permitted but upon whether the composition creates a path connecting domains that the system-level policy intended to separate.

The structural property can therefore concern the transitive closure
[
E^{+}
]
rather than the observed edge set $E$ itself. A local verifier may check
[
P(u,v)
]
for each direct edge $(u,v)\in E$, while the consequential property concerns whether
[
(u,w)\in E^{+}.
]
Unless $P$ is known to compose transitively in the relevant manner, direct-edge acceptance cannot establish a path-level guarantee.

This yields a general form of the problem.

\begin{proposition}[Path composition requires a closure law]
Let $P(u,v)$ be a predicate evaluated on direct workflow relations. Suppose a desired system-level property requires that a relation $Q(u,w)$ hold whenever there exists a path
[
u=v_0\rightarrow v_1\rightarrow\cdots\rightarrow v_k=w.
]
Local verification of
[
P(v_i,v_{i+1})
]
for every edge in the path establishes $Q(u,w)$ only if a composition law has been established that maps the sequence of $P$-satisfying relations to $Q$.
\end{proposition}

\begin{proof}
Without such a law, choose a model in which
[
P(v_i,v_{i+1})=1
]
for every direct edge but
[
Q(u,w)=0.
]
The local premises remain satisfied while the desired path-level conclusion fails. Therefore the conjunction of local predicates is insufficient absent an additional implication connecting their composition to $Q$.
\end{proof}

Again the proposition is elementary, but it identifies precisely where the substantive burden lies. The important theorem is not that each edge has passed. The important theorem is that the property established on edges is preserved under path composition.

This is familiar in domains where compositional reasoning has been developed carefully. A type system, capability system, information-flow discipline, or cryptographic protocol may deliberately define local judgments so that a global theorem follows. In such cases local verification works precisely because a preservation theorem exists. The success of those systems should not be generalized into the assumption that arbitrary learned safety predicates possess the same compositional structure.

Skynet is especially relevant because its design attempts to retain semantic and structural information jointly. Its empirical premise is that executions can exhibit anomalous organization not adequately exposed by inspecting events independently. In the terminology developed here, the workflow graph is a richer projection
[
\pi_{\mathrm{graph}}(H)
]
of the underlying execution history than an eventwise projection
[
\pi_{\mathrm{evt}}(H).
]
If there exists a map
[
r:\Img(\pi_{\mathrm{graph}})
\rightarrow
\Img(\pi_{\mathrm{evt}})
]
such that
[
\pi_{\mathrm{evt}}

r\circ\pi_{\mathrm{graph}},
]
then the eventwise representation can be recovered from the graph representation, while the converse need not hold. The graph interface is informationally at least as fine as the eventwise interface with respect to the distinctions preserved by $r$.

This motivates a general ordering of observational representations.

\begin{definition}[Observational refinement]
Let
[
\pi_1:\mathcal{X}\rightarrow\mathcal{O}_1
]
and
[
\pi_2:\mathcal{X}\rightarrow\mathcal{O}_2
]
be two representations of the same underlying objects. We say that $\pi_1$ is at least as informative as $\pi_2$, written
[
\pi_1\succeq\pi_2,
]
if there exists a map
[
r:\Img(\pi_1)\rightarrow\Img(\pi_2)
]
such that
[
\pi_2=r\circ\pi_1.
]
\end{definition}

Under this ordering, $\pi_1$ preserves every distinction required to reconstruct $\pi_2$, although it may preserve additional distinctions as well. The relation immediately yields a monotonicity result for sufficient support.

\begin{proposition}[Sufficiency under refinement]
If
[
\pi_1\succeq\pi_2
]
and $\Phi$ factors through $\pi_2$, then $\Phi$ also factors through $\pi_1$.
\end{proposition}

\begin{proof}
Since $\pi_1\succeq\pi_2$, there exists $r$ satisfying
[
\pi_2=r\circ\pi_1.
]
Since $\Phi$ factors through $\pi_2$, there exists $\widehat{\Phi}_2$ satisfying
[
\Phi=\widehat{\Phi}_2\circ\pi_2.
]
Substitution gives
[
\Phi

\widehat{\Phi}_2\circ r\circ\pi_1.
]
Defining
[
\widehat{\Phi}_1

\widehat{\Phi}_2\circ r
]
yields
[
\Phi

\widehat{\Phi}_1\circ\pi_1.
]
Thus $\Phi$ factors through $\pi_1$.
\end{proof}

The converse does not hold. A graph representation can be sufficient for a structural property while an eventwise representation is not. Refinement therefore provides one formal sense in which workflow-level monitoring can increase the class of properties that are in principle recoverable.

It does not follow, however, that the finest possible representation is always preferable. The temporal case already showed that information supplied to an acting system can alter its behavior. Even for a passive verifier, finer representations can increase computational cost, privacy exposure, attack surface, statistical variance, and the opportunity for irrelevant correlations. Sufficiency is therefore a better design objective than maximal information. The aim is to preserve the distinctions required by the property, not to retain every distinction available in the world.

The graph case introduces a second and equally important problem. Suppose the workflow representation is sufficient for determining whether an execution is unusual relative to a learned distribution. Let
[
A:\mathcal{G}\rightarrow\mathbb{R}
]
be an anomaly score, with larger values indicating greater deviation from the learned benign region. A detector may classify a workflow as anomalous whenever
[
A(G)>\tau.
]
If the detector performs well, then $A$ can be an excellent diagnostic coordinate. But the property
[
A(G)>\tau
]
is not identical to the property
[
N(G)=0,
]
where $N$ denotes a stronger judgment such as impermissibility, unsafety, task invalidity, or justification for blocking execution.

The distinction is categorical rather than merely statistical. An anomaly score answers a question about the relationship between an observation and a reference distribution or learned structural manifold. A prohibition predicate answers a question about whether an execution should be permitted under some governing criterion. The two may correlate strongly. They are not therefore the same function.

To infer
[
N(G)=0
]
from
[
A(G)>\tau,
]
one needs an additional calibration or decision relation. For example, one might establish
[
P(N(G)=0\mid A(G)>\tau)\geq 1-\varepsilon
]
on a specified distribution. Even this would yield a probabilistic statement relative to that distribution, not a logical equivalence. Alternatively, a governing specification could explicitly define certain structural patterns as prohibited, in which case the implication would follow from the specification rather than from anomaly alone.

Without such an additional relation,
[
A(G)>\tau
\quad\not\Rightarrow\quad
N(G)=0.
]

This is another instance of verification scope. The anomaly detector may be entirely correct about the property it computes while a downstream system overextends the meaning of that property. The failure occurs not because the detector is inaccurate but because its output is silently reinterpreted.

The inverse error is possible as well. A structurally typical workflow is not automatically acceptable:
[
A(G)\leq\tau
\quad\not\Rightarrow\quad
N(G)=1.
]
A harmful execution may closely resemble common workflows. A novel but legitimate execution may appear highly anomalous. Distributional typicality and normative status are therefore neither equivalent nor ordered by logical implication in the absence of additional assumptions.

This gives four logical possibilities:
[
\begin{array}{c|cc}
& N(G)=1 & N(G)=0\
\hline
A(G)\leq\tau & \text{typical and permitted} & \text{typical but impermissible}\
A(G)>\tau & \text{unusual but permitted} & \text{unusual and impermissible}
\end{array}
]
and a system that collapses anomaly into prohibition erases the off-diagonal cases by construction rather than by evidence.

This distinction is particularly important for systems expected to perform novel tasks. Novelty is often the point of deploying a general agent. If a safety architecture treats departure from the learned manifold as sufficient evidence of impermissibility, it can transform uncertainty into prohibition without establishing the missing relation. Conversely, if novelty is always permitted because anomaly is regarded as merely statistical, the diagnostic value of structural deviation is discarded. The appropriate role of anomaly is therefore neither automatic permission nor automatic refusal. It is evidence that can alter what should be checked next.

One can represent this by separating a diagnostic map
[
D:\mathcal{G}\rightarrow\mathcal{D}
]
from a decision rule
[
R:\mathcal{D}\times\mathcal{E}\rightarrow\mathcal{U},
]
where $\mathcal{E}$ contains additional evidence and $\mathcal{U}$ contains possible responses. The anomaly score may be one component of
[
D(G),
]
while the response may include continued execution, additional inspection, restricted execution, deferment, or refusal. This architecture preserves the informational contribution of structural anomaly without pretending that anomaly itself contains the complete decision criterion.

The distinction can be stated as another factorization problem. Suppose the ultimate decision property is
[
N:\mathcal{G}\rightarrow{0,1}.
]
If there exist graphs $G_1,G_2$ satisfying
[
A(G_1)=A(G_2)
]
but
[
N(G_1)\neq N(G_2),
]
then $N$ does not factor through $A$. No decision rule that sees only the scalar anomaly score can determine $N$ universally. More elaborate downstream reasoning cannot repair the missing distinction unless additional information is supplied.

\begin{corollary}[Anomaly-score insufficiency]
If a scalar anomaly score identifies two workflows having different target decision values, then the anomaly score is not a sufficient verification unit for that decision property.
\end{corollary}

The corollary generalizes beyond anomaly detection. Whenever a high-dimensional workflow is compressed into a scalar risk, confidence, anomaly, or safety score, the scalar representation induces equivalence classes
[
G_1\sim_A G_2
\quad\Longleftrightarrow\quad
A(G_1)=A(G_2).
]
Any property that varies within one of these equivalence classes cannot be recovered universally from the score alone.

This suggests a useful way to characterize diagnostic refinement. Let
[
D_1:\mathcal{G}\rightarrow\mathcal{Y}1
]
and
[
D_2:\mathcal{G}\rightarrow\mathcal{Y}2
]
be two diagnostic representations. Each induces an equivalence relation:
[
G\sim{D_i}G'
\quad\Longleftrightarrow\quad
D_i(G)=D_i(G').
]
We say $D_1$ is diagnostically finer than $D_2$ when
[
G\sim{D_1}G'
\Longrightarrow
G\sim_{D_2}G'
]
for all workflows $G,G'$. Equivalently, every distinction preserved by $D_2$ is also preserved by $D_1$. This is the same refinement relation introduced above, expressed through induced partitions.

The value of a richer structural representation is therefore not that it is larger in dimensionality. A thousand redundant features need not distinguish more relevant cases than one well-chosen feature. What matters is whether the representation separates workflows that the target property requires the verifier to distinguish.

This observation guards against a simplistic reading of workflow-level verification. The lesson is not that graphs are intrinsically superior to sequences, or that structural embeddings should replace local checks. The lesson is conditional: when the property depends upon relations erased by an eventwise representation, the eventwise representation is insufficient. A graph becomes useful because it preserves those relations.

The structural case also reveals why the phrase ``global verifier'' can be misleading. A verifier may observe the entire graph and still operate on an insufficient representation of the underlying execution. The graph may omit timing, hidden state, external effects, authorization context, semantic distinctions, or counterfactual alternatives required by the target property. Globality relative to one representation does not imply sufficiency relative to every property.

Suppose the complete execution is $H$ and the graph is obtained through
[
\gamma:\Hist\rightarrow\mathcal{G}.
]
A graph-level verifier observes
[
G=\gamma(H).
]
For a history-level property
[
\Phi:\Hist\rightarrow{0,1},
]
graph-level verification is sufficient only if there exists
[
\widehat{\Phi}:\Img(\gamma)\rightarrow{0,1}
]
such that
[
\Phi=\widehat{\Phi}\circ\gamma.
]
If two histories generate the same workflow graph but differ on $\Phi$, then even perfect graph analysis cannot determine $\Phi$. The graph has become another checkpoint.

This is the structural analogue of the temporal result. Expanding the observational unit solves only those insufficiencies caused by distinctions newly preserved by the expansion. It does not abolish the need to match representation to property.

The contribution of Zhang et al.\ can therefore be situated carefully. Their work provides empirical evidence that semantic and structural workflow representations can expose anomalies obscured by more local inspection. That supports a move from isolated events toward relational representations for properties that depend upon workflow organization. It does not establish that structural anomaly is itself a complete safety criterion, nor would such a conclusion follow from anomaly-detection performance alone. The detector supplies a richer diagnostic coordinate. What should be inferred from that coordinate remains a separate question.

The structural failure of composition can now be stated precisely. Local properties of events do not necessarily determine properties of their composition. A graph-level representation can preserve distinctions erased by eventwise inspection, but the properties established on that graph still have a scope. Structural abnormality is a property of relation to a structural reference class. It is not automatically a property of permission, safety, or correctness.

Verification therefore encounters two boundaries in succession. The first lies between components and their composition:
[
{v,e}
\longrightarrow
G.
]
The second lies between structural diagnosis and consequential judgment:
[
G
\longrightarrow
N(G).
]
Neither crossing is free. The first requires preservation of relational information. The second requires an argument connecting the diagnostic property to the decision property.

The checkpoint has expanded from an event to a graph, but the governing principle has not changed. A verifier may conclude only what its representation and its established preservation relations support. Enlarging the object of inspection is valuable when it restores distinctions the target property requires. It becomes another form of overreach when the richer observation is allowed to stand for conclusions it does not itself establish.

\section{Representational Composition}

The structural case concerns information lost when relations among computational events are omitted. A different failure occurs when the object itself is transformed before verification. In this setting, the verifier may be formally sound, the proof it produces may be valid, and the verified representation may genuinely satisfy the asserted property, while the conclusion drawn about the original object remains unjustified. The missing premise concerns the transformation connecting the source to the object that actually entered the verifier.

Lahiri's work on neuro-formal verification~\cite{lahiri2026} makes this boundary unusually visible. Formal verification is attractive because it can replace an informal judgment of plausibility with a machine-checked claim relative to an explicit formal object and specification. Yet the programs one wishes to analyze are not always presented in a language or form directly suitable for the chosen verification system. Neuro-formal architectures therefore place a learned transformation before a formal verifier. An artificial-intelligence component translates, annotates, restructures, or otherwise prepares the source program for verification, after which a conventional formal tool attempts to prove correctness or produce a counterexample.

Lahiri reports that the staged system can return Dafny proofs of correctness or of a bug for 57 percent of the evaluated dataset at 92 percent precision, while the CBMC path returns counterexamples for 63 percent of buggy programs at 90 percent precision. The empirical significance of these results lies partly in the contrast between disciplined staging and arrangements in which the learned system and verification objective are less cleanly separated. For the present argument, however, the more general importance is architectural. The formal verifier operates downstream of a representational transformation. Its soundness applies to the formal object it receives. Whether that result can be transported back to the original source is a separate question.

Let
[
X
]
denote a space of source objects and
[
Y
]
a space of verification-ready objects. Let
[
\tau:X\rightarrow Y
]
be the transformation performed before verification. A formal verifier evaluates a property
[
P_Y:Y\rightarrow{0,1}.
]
The proposition actually established at the terminal stage is
[
P_Y(\tau(x))=1.
]
The proposition ultimately desired may instead be
[
P_X(x)=1,
]
where
[
P_X:X\rightarrow{0,1}
]
expresses the corresponding property of the source object.

Nothing about the syntax of these expressions guarantees
[
P_Y(\tau(x))=1
\Longrightarrow
P_X(x)=1.
]
That implication is itself a theorem about $\tau$, $P_X$, and $P_Y$.

This distinction is easiest to see by separating three claims that are often compressed into the phrase ``the program was verified.'' The first concerns the terminal verifier. The second concerns the transformation. The third concerns the pipeline.

Let
[
\mathcal{V}:Y\rightarrow{\text{accept},\text{reject},\text{unknown}}
]
be a verifier. A soundness statement for $\mathcal{V}$ may have the form
[
\mathcal{V}(y)=\text{accept}
\Longrightarrow
P_Y(y)=1.
]
This is a statement about objects in $Y$.

A source-reflection statement for $\tau$ has the form
[
P_Y(\tau(x))=1
\Longrightarrow
P_X(x)=1.
]
This is a statement about the relation between $X$ and $Y$.

Only by composing these two implications does one obtain
[
\mathcal{V}(\tau(x))=\text{accept}
\Longrightarrow
P_X(x)=1.
]

The distinction can be stated formally.

\begin{theorem}[Pipeline soundness by composition]
Let
[
\tau:X\rightarrow Y
]
be a transformation, let
[
P_X:X\rightarrow{0,1}
]
and
[
P_Y:Y\rightarrow{0,1}
]
be properties, and let
[
\mathcal{V}:Y\rightarrow{\mathrm{accept},\mathrm{reject},\mathrm{unknown}}
]
be a verifier. Suppose
[
\mathcal{V}(y)=\mathrm{accept}
\Longrightarrow
P_Y(y)=1
]
for every $y\in Y$, and suppose
[
P_Y(\tau(x))=1
\Longrightarrow
P_X(x)=1
]
for every $x\in X$. Then
[
\mathcal{V}(\tau(x))=\mathrm{accept}
\Longrightarrow
P_X(x)=1
]
for every $x\in X$.
\end{theorem}

\begin{proof}
Assume
[
\mathcal{V}(\tau(x))=\mathrm{accept}.
]
By soundness of $\mathcal{V}$,
[
P_Y(\tau(x))=1.
]
By the source-reflection property of $\tau$,
[
P_X(x)=1.
]
Therefore
[
\mathcal{V}(\tau(x))=\mathrm{accept}
\Longrightarrow
P_X(x)=1.
]
\end{proof}

The theorem is simple because the substantive difficulty has been placed where it belongs. Terminal soundness is not pipeline soundness. Pipeline soundness follows only after a second property has been established for the transformation.

The required direction of preservation deserves care. Suppose the intended semantics of $P_X$ and $P_Y$ correspond closely enough that one hopes for
[
P_X(x)
\Longleftrightarrow
P_Y(\tau(x)).
]
This biconditional contains two distinct guarantees:
[
P_X(x)
\Longrightarrow
P_Y(\tau(x))
]
and
[
P_Y(\tau(x))
\Longrightarrow
P_X(x).
]
The first says that source-level satisfaction survives translation. The second says that satisfaction of the translated object reflects back to the source. For a verifier whose acceptance is being used to certify the source, the second direction is indispensable. A transformation could preserve all correct programs while also mapping some incorrect programs into correct verification objects. Such a transformation would preserve truth in the forward direction while remaining unsafe as a certification interface.

This can be demonstrated directly.

\begin{proposition}[Forward preservation is insufficient for certification]
The condition
[
P_X(x)=1
\Longrightarrow
P_Y(\tau(x))=1
]
does not imply
[
P_Y(\tau(x))=1
\Longrightarrow
P_X(x)=1.
]
Consequently, forward preservation alone does not justify source-level certification from a proof of the transformed object.
\end{proposition}

\begin{proof}
Let
[
X={x_0,x_1},
\qquad
Y={y},
]
and define
[
P_X(x_0)=0,
\qquad
P_X(x_1)=1,
]
while
[
P_Y(y)=1.
]
Let
[
\tau(x_0)=\tau(x_1)=y.
]
Then whenever
[
P_X(x)=1,
]
we have
[
P_Y(\tau(x))=1,
]
so forward preservation holds. However,
[
P_Y(\tau(x_0))=1
]
while
[
P_X(x_0)=0.
]
Thus reflection fails.
\end{proof}

The example is structurally identical to projection insufficiency. The transformation has identified two source objects that differ on the target property:
[
\tau(x_0)=\tau(x_1)
]
while
[
P_X(x_0)\neq P_X(x_1).
]
Once that identification has occurred, no verifier operating solely on $\tau(x)$ can universally recover $P_X(x)$.

This gives a direct corollary.

\begin{corollary}[Representational insufficiency]
If there exist $x_1,x_2\in X$ such that
[
\tau(x_1)=\tau(x_2)
]
but
[
P_X(x_1)\neq P_X(x_2),
]
then no verifier whose complete input is $\tau(x)$ can universally determine $P_X(x)$.
\end{corollary}

The result is stronger than the claim that translation errors can cause verification errors. A translation need not be globally erroneous to be insufficient. It need only erase one distinction upon which the target property depends. The downstream verifier can then be perfect and still fail at the source-level task.

This distinction suggests that "faithfulness'' should not be treated as a single undifferentiated property. A representation can be faithful for one question and unfaithful for another. Let \[ P,Q:X\rightarrow\{0,1\} \] be two source properties. It is possible for $\tau$ to preserve every distinction required by $P$ while erasing distinctions required by $Q$. In factorization form, there may exist \[ \widehat{P}:Y\rightarrow\{0,1\} \] such that \[ P=\widehat{P}\circ\tau, \] while no \[ \widehat{Q}:Y\rightarrow\{0,1\} \] satisfies \[ Q=\widehat{Q}\circ\tau. \] Calling $\tau$ simply "faithful'' or ``unfaithful'' obscures this dependence on the property under consideration.

A more precise definition is therefore useful.

\begin{definition}[Property-relative faithfulness]
A transformation
[
\tau:X\rightarrow Y
]
is \emph{faithful for a property} $P_X:X\rightarrow Z$ if there exists
[
\widehat{P}:Y\rightarrow Z
]
such that
[
P_X=\widehat{P}\circ\tau.
]
Equivalently,
[
\tau(x_1)=\tau(x_2)
\Longrightarrow
P_X(x_1)=P_X(x_2).
]
\end{definition}

Under this definition, representational fidelity is exactly the sufficient-support condition applied to a transformation. The question is not whether every detail of the source survives. Most useful representations intentionally discard detail. The question is whether the distinctions relevant to the claimed property survive.

This is important because an injective representation is sufficient for every source property in a purely information-theoretic sense, but injectivity is usually unnecessary and often undesirable. Verification representations exist precisely to simplify. They abstract away syntax, implementation detail, irrelevant state, or environmental complexity. The objective is therefore not perfect reproduction of the source. It is controlled forgetting.

Let $\tau$ induce an equivalence relation on $X$:
[
x_1\sim_\tau x_2
\quad\Longleftrightarrow\quad
\tau(x_1)=\tau(x_2).
]
The representation is faithful for $P_X$ exactly when $P_X$ is constant on every equivalence class of $\sim_\tau$. Thus abstraction is safe for a property when everything it collapses is irrelevant to that property.

This provides a useful reinterpretation of abstraction in verification. An abstraction is not safe because it is detailed. It is safe because its identifications respect the property one intends to prove.

Suppose, for example, that source objects contain variables
[
x=(u,v,w),
]
while a representation discards $w$:
[
\tau(u,v,w)=(u,v).
]
For any property of the form
[
P_X(u,v,w)=Q(u,v),
]
the abstraction is sufficient. But if
[
P_X(u,v,w)
]
depends on $w$, then two objects differing only in $w$ can map to the same representation while differing on the target property. No terminal verifier over $(u,v)$ can repair that loss.

This observation scales directly to learned transformations. A language model that rewrites a source program into a verification language is not merely changing syntax. It is choosing which distinctions survive into the formal object. If a distinction affecting source semantics is omitted, weakened, strengthened, or reinterpreted, the terminal proof concerns a different object.

The resulting architecture can be represented as
[
x
\xrightarrow{\tau}
y
\xrightarrow{\mathcal{V}}
c,
]
where $c$ is a certificate or counterexample. It is tempting to concentrate all rigor in the second arrow because $\mathcal{V}$ may be a mature formal tool. Yet the first arrow determines what proposition the second arrow is capable of proving about the source.

The complete proof obligation therefore has at least two layers:
[
\mathcal{V}(y)=\mathrm{accept}
\Longrightarrow
P_Y(y)
]
and
[
P_Y(\tau(x))
\Longrightarrow
P_X(x).
]
The first is verifier soundness. The second is translation reflection. Their composition yields source certification.

If the transformation is itself probabilistic,
[
\tau_\omega:X\rightarrow Y,
]
where $\omega$ denotes sampling randomness or model-dependent variation, the problem becomes more complicated. A source object $x$ can yield different verification objects under repeated transformation:
[
\tau_{\omega_1}(x)\neq\tau_{\omega_2}(x).
]
One must then distinguish the probability that the terminal verifier is sound from the probability that the transformation preserves the target property.

Suppose terminal soundness is exact:
[
\mathcal{V}(y)=\mathrm{accept}
\Longrightarrow
P_Y(y)=1.
]
Suppose, however, that the transformation satisfies the reflection condition only with probability at least $1-\varepsilon$:
[
P_\omega
\left(
P_Y(\tau_\omega(x))=1
\Longrightarrow
P_X(x)=1
\right)
\geq 1-\varepsilon.
]
The resulting source-level guarantee cannot exceed the reliability of the transformation merely because the terminal proof is exact. Formal certainty downstream does not erase probabilistic uncertainty upstream.

Indeed, if the source is incorrect,
[
P_X(x)=0,
]
and the transformation has probability $\varepsilon$ of mapping it into a verification object satisfying $P_Y$, then a perfectly sound terminal verifier can certify the transformed object whenever that representational failure occurs. The proof remains valid. The pipeline conclusion is wrong.

This is an important form of what might be called \emph{certainty laundering}. The term should be used carefully because the formal verifier has done nothing incorrect. The problem is architectural. A probabilistic or semantically uncertain transformation produces an object, a formal stage produces a categorical certificate about that object, and the categorical character of the certificate can obscure the uncertainty introduced before the formal stage.

One can state the issue more neutrally as a conservation principle for uncertainty: a downstream proof cannot retroactively remove uncertainty about whether the proved object preserves the source property. If the uncertainty concerns the relation between $x$ and $\tau(x)$, then proving facts about $\tau(x)$ leaves that relational uncertainty intact.

The distinction is particularly sharp when the transformation has access to the desired verification outcome. Lahiri's staged architecture is relevant here because separating transformation from proof goals can reduce opportunities for a learned system to produce representations optimized to satisfy the verifier rather than faithfully encode the source. The general danger can be expressed without attributing any particular failure mode to the paper.

Let a transformation be
[
\tau:X\times G\rightarrow Y,
]
where $G$ represents information about the desired goal, certificate, or acceptance condition. If $\tau$ can condition on the property it is expected to make provable, then the generated representation can depend upon the target conclusion:
[
\tau(x,g_1)\neq\tau(x,g_2).
]
This does not imply that the transformation is unfaithful, but it introduces a channel through which optimization for provability can diverge from preservation of source semantics.

A goal-blind transformation instead has the form
[
\tau:X\rightarrow Y,
]
or more generally satisfies a conditional independence relation under which the representation is generated without access to the terminal verdict being sought. Such staging does not prove fidelity. It removes one pathway by which the transformation can adapt the representation to the desired conclusion.

This motivates a broader architectural principle. When one component produces the object another component is supposed to verify, independence between representation generation and verdict optimization can be epistemically valuable even when neither component is perfect. Separation does not substitute for proof, but it can make the proof obligation more intelligible.

Representational composition also exposes a distinction between verification and validation that is often blurred in discussions of formal methods. Verification can establish that a formal object satisfies a formal specification. Validation concerns whether that object and specification adequately correspond to the source problem one intended to solve. A perfectly verified model of the wrong system is not a contradiction. It is a successful proof attached to an unsuccessful representation.

For agentic systems, the same pattern appears well beyond source-code translation. Let
[
u
]
be a natural-language instruction, let
[
p=\tau_1(u)
]
be a plan, let
[
c=\tau_2(p)
]
be a sequence of tool calls, and let
[
e=\tau_3(c)
]
be the resulting environmental change. A verifier may check properties at any of these levels:
[
P_U(u),\quad
P_P(p),\quad
P_C(c),\quad
P_E(e).
]
A proof or classification at one level does not automatically establish the corresponding property at another. Each transition
[
u\rightarrow p\rightarrow c\rightarrow e
]
requires a preservation argument for whichever property is intended to survive.

The same is true of summarization. Suppose a long interaction history
[
H
]
is compressed into
[
S=\sigma(H).
]
A verifier operating on $S$ can determine a history property $\Phi(H)$ only if
[
\Phi

\widehat{\Phi}\circ\sigma.
]
A summary can be linguistically excellent while being verification-inadequate because it omits a distinction irrelevant to narrative coherence but essential to the target property.

This is especially consequential in long-running agents where context compression is operationally unavoidable. The system may retain a compact summary rather than the full history. If later decisions or verifications depend upon a property that does not factor through that summary, the information has not merely become harder to retrieve. It has become unrecoverable from the retained representation.

The relationship between abstraction and verification can therefore be represented by a commuting diagram. One desires
[
\begin{array}{ccc}
X & \xrightarrow{\tau} & Y\
P_X\downarrow & & \downarrow P_Y\
Z & = & Z
\end{array}
]
in the sense that
[
P_X

P_Y\circ\tau
]
when exact preservation is appropriate, or at least the one-directional relation required by the intended inference. When the diagram fails to commute, the result established in $Y$ cannot simply be read back into $X$.

For certification, the relevant requirement may be written
[
P_Y\circ\tau
\preceq
P_X,
]
where the order on Boolean predicates is defined pointwise and $0<1$. This expresses
[
P_Y(\tau(x))=1
\Longrightarrow
P_X(x)=1.
]
For bug detection, the desired direction may instead differ. The important point is that the direction of preservation must be chosen according to the inference the system intends to make. ``Semantic equivalence'' is stronger than necessary for some tasks and weaker formulations can suffice, but some explicit relation is unavoidable.

This also clarifies why empirical precision measurements and formal soundness statements answer different questions. If an end-to-end system has measured precision
[
\operatorname{Prec}

P(P_X(x)=1\mid \mathcal{V}(\tau(x))=\mathrm{accept}),
]
that quantity evaluates the complete empirical pipeline on a distribution. It does not by itself identify whether errors arise from $\tau$, $\mathcal{V}$, the specification, the benchmark labels, or distribution shift. Conversely, a proof that $\mathcal{V}$ is sound over $Y$ does not determine end-to-end precision over $X$.

A rigorous evaluation should therefore distinguish at least the terminal-verification question
[
\mathcal{V}(y)=\mathrm{accept}
\Rightarrow
P_Y(y),
]
the representation question
[
P_Y(\tau(x))
\Rightarrow
P_X(x),
]
and the empirical pipeline question
[
P(P_X(x)=1\mid
\mathcal{V}(\tau(x))=\mathrm{accept}).
]
These quantities can differ dramatically without contradiction because they quantify different relations.

The representation problem also interacts with observational refinement from the previous section. Suppose
[
\tau_1:X\rightarrow Y_1
]
and
[
\tau_2:X\rightarrow Y_2
]
are two candidate representations, with
[
\tau_1\succeq\tau_2.
]
If $P_X$ factors through $\tau_2$, then it also factors through $\tau_1$. A finer representation cannot lose a property already recoverable from a coarser one, provided the refinement relation is exact. But the finer representation may be more expensive, more difficult to verify, or more vulnerable to irrelevant complexity. Once again, maximal detail is not the objective. Sufficient support is.

This leads to a useful optimization problem. Let
[
C(\tau)
]
denote the cost of operating on a representation, broadly construed to include computational expense, proof complexity, storage, privacy exposure, and other burdens. One can seek
[
\tau^\star
\in
\arg\min_\tau C(\tau)
]
subject to
[
P_X

\widehat{P}\circ\tau
]
for some $\widehat{P}$. The objective is not to preserve the source completely. It is to find an economical representation sufficient for the property.

For multiple target properties
[
P_1,\ldots,P_m,
]
the constraint becomes
[
P_j

\widehat{P}_j\circ\tau
\qquad
\text{for every }j.
]
As the set of properties expands, the required representation may need to become finer. A representation adequate for functional correctness may be inadequate for timing leakage. One adequate for memory safety may be inadequate for authorization provenance. One adequate for predictive accuracy may be inadequate for intervention, which is precisely the transition considered in the next section.

The representational failure therefore cannot be solved merely by attaching a formal verifier to an artificial-intelligence system. Formal verification changes the epistemic status of claims about the formal object. It does not automatically change the epistemic status of the transformation by which that object was obtained. If the transformation is learned, heuristic, lossy, goal-conditioned, or otherwise fallible, its relation to the source remains part of the verification problem.

Lahiri's work is important in this context because it places the boundary between learned transformation and formal reasoning in view. The lesson is not that neuro-formal systems are inherently unreliable. On the contrary, careful staging can make them more reliable by preventing one form of semantic slippage from being hidden inside the verifier. The deeper lesson is that the rigor of the terminal stage cannot be used as a substitute for rigor about the interface feeding that stage.

The checkpoint has now moved again. In the temporal case, the problem concerned the relation between a current action and the history conditioning it. In the structural case, it concerned the relation between events and the workflow they compose. Here it concerns the relation between an original object and the representation presented for proof. In each case the same mistake is possible: the property established at the checkpoint is allowed to travel across a transformation whose preservation behavior has not itself been established.

Representation is therefore not a neutral prelude to verification. It determines which distinctions reach the verifier, which source objects become indistinguishable, and which propositions a terminal proof can legitimately support. A sound verifier over an insufficient representation is still sound. What fails is the inference that its soundness extends farther than the representation permits.

The appropriate question is consequently not only whether the verifier can prove $P_Y(y)$. It is whether the source-level property one actually cares about factors through the representation $y=\tau(x)$. If it does, formal verification can be transported across the boundary under the established preservation relation. If it does not, no strength of terminal proof can reconstruct what the transformation has already erased.

In this sense, verification begins before formal verification begins. The first proof obligation is often not the theorem checked at the end of the pipeline, but the claim that the object entering the theorem prover is an adequate carrier of the property one intends the theorem to establish.

\section{Epistemic Composition}

The representational problem concerns whether a verified property survives the transformation from a source object into the object actually presented to a verifier. The fourth case introduces a related but distinct boundary. Even when the representation is adequate for prediction, and even when the prediction is accurate, the information sufficient to predict an outcome need not be sufficient to determine whether an intervention should occur. The transition from prediction to action changes not merely the representation but the proposition under consideration.

Ma, Huang, Kim, and Jang~\cite{ma2026} develop this problem in the context of world models for safety-critical embodied systems. Their analysis begins from the observation that world models are commonly evaluated through predictive criteria such as likelihood, reconstruction quality, or visual fidelity. These criteria measure important capabilities, but systems operating in consequential environments require distinctions that predictive performance alone may not preserve. The authors emphasize mismatches between likelihood and risk, prediction and intervention, finite-horizon prediction and accumulated consequences, and predictive representation and decision-relevant representation. Their proposed direction consequently incorporates counterfactual reasoning, safety-critical episodic memory, epistemic uncertainty, recoverability, runtime assurance, and responses such as sensing, revision, deferment, and abstention alongside ordinary environment-changing action.

The distinction can be formalized by separating a predictive model from a decision rule. Let
[
s_t\in\mathcal{S}
]
denote the current state and
[
a_t\in\mathcal{A}
]
a candidate action. A world model may estimate a transition distribution
[
M(s_{t+1}\mid s_t,a_t).
]
If $M$ is perfectly calibrated, one may even suppose
[
M(\cdot\mid s_t,a_t)

P(\cdot\mid s_t,a_t),
]
where $P$ is the true conditional transition law. This is a very strong predictive assumption. It still does not determine whether
[
a_t
]
should be chosen.

A decision rule requires additional structure. In the simplest expected-utility formulation, one introduces a utility or cost function
[
U:\mathcal{S}\times\mathcal{A}\times\mathcal{S}\rightarrow\mathbb{R}
]
and selects an action according to
[
a_t^\star
\in
\arg\max_{a\in\mathcal{A}}
\mathbb{E}_{s'\sim M(\cdot\mid s_t,a)}
\left[
U(s_t,a,s')
\right].
]
The transition model and the utility function perform different roles. The former estimates what may happen. The latter determines how possible consequences enter a choice criterion. Predictive accuracy alone cannot recover the utility function.

This yields an elementary impossibility result.

\begin{proposition}[Prediction does not determine preference]
Let $M$ be a fixed predictive model over the consequences of actions. There exist two decision problems having identical $M$ but different optimal actions.
\end{proposition}

\begin{proof}
Let
[
\mathcal{A}={a,b}
]
and suppose the predictive model deterministically maps
[
a\mapsto s_a,
\qquad
b\mapsto s_b.
]
Thus
[
M(s_a\mid s,a)=1
]
and
[
M(s_b\mid s,b)=1.
]
Define one utility function $U_1$ such that
[
U_1(s,a,s_a)>U_1(s,b,s_b),
]
and another utility function $U_2$ such that
[
U_2(s,a,s_a)<U_2(s,b,s_b).
]
The predictive model is identical in both decision problems, but
[
a^\star_1=a
]
and
[
a^\star_2=b.
]
Therefore the prediction does not determine the preferred intervention.
\end{proof}

The proposition is obvious from decision theory, but it identifies a recurrent scope error in systems where prediction and action are tightly integrated. A model can become arbitrarily accurate about
[
P(s_{t+1}\mid s_t,a_t)
]
without thereby answering which consequences matter, how they should be weighted, what uncertainty should be tolerated, which alternatives should be compared, or whether the system currently possesses sufficient evidence to intervene.

The distinction can be written as a non-implication:
[
P(s_{t+1}\mid s_t,a_t)
\quad\not\Rightarrow\quad
\operatorname{Choose}(a_t).
]
The left-hand side is predictive. The right-hand side is decisional. Treating the former as sufficient support for the latter requires additional structure.

The problem becomes sharper when risk differs from likelihood. Expected value is itself only one decision criterion. In safety-critical systems, low-probability outcomes can dominate decisions because their costs are severe, irreversible, or poorly recoverable. Let a candidate action $a$ induce outcomes
[
s_1,\ldots,s_m
]
with probabilities
[
p_1,\ldots,p_m.
]
A likelihood-oriented model may devote most of its representational capacity to accurately modeling high-probability outcomes. A safety criterion may instead be highly sensitive to a state $s_j$ having small $p_j$ but catastrophic cost.

Let
[
L(a)

\mathbb{E}[\ell(s')\mid a]
]
be an expected loss. Even this quantity can obscure tail-sensitive objectives. A risk functional may instead take the form
[
\rho_a(\ell)
]
where $\rho_a$ could represent a worst-case criterion, a conditional tail expectation, a chance constraint, or another risk-sensitive functional. Two predictive representations sufficient to estimate ordinary likelihood can differ in whether they preserve the information required to estimate $\rho_a$.

The mismatch can again be expressed through factorization. Let
[
\pi_{\mathrm{pred}}(H)
]
be the representation retained by a world model, and let
[
R(H,a)
]
be the decision-relevant risk property. Predictive sufficiency for some forecasting target $F$ means
[
F

\widehat{F}\circ\pi_{\mathrm{pred}}.
]
It does not imply that there exists
[
\widehat{R}
]
such that
[
R

\widehat{R}\circ\pi_{\mathrm{pred}}.
]
A representation can therefore be sufficient for prediction and insufficient for risk.

\begin{proposition}[Predictive sufficiency does not imply decisional sufficiency]
There exist a representation $\pi$, a predictive property $F$, and a decision property $D$ such that
[
F=\widehat{F}\circ\pi
]
for some $\widehat{F}$, while no function $\widehat{D}$ satisfies
[
D=\widehat{D}\circ\pi.
]
\end{proposition}

\begin{proof}
Let the underlying states be
[
x_1=(p,c_1),
\qquad
x_2=(p,c_2),
]
where $p$ is a predictive quantity and $c_i$ is a consequence parameter. Let
[
\pi(x_i)=p
]
for both states. Define
[
F(x_i)=p.
]
Then
[
F=\operatorname{id}\circ\pi,
]
so $\pi$ is sufficient for $F$.

Now define a decision property
[
D(x_i)

\begin{cases}
1,&c_i\leq c^\star,\
0,&c_i>c^\star,
\end{cases}
]
and choose
[
c_1\leq c^\star<c_2.
]
Then
[
\pi(x_1)=\pi(x_2)
]
but
[
D(x_1)\neq D(x_2).
]
By projection insufficiency, $D$ cannot factor through $\pi$.
\end{proof}

The proposition formalizes the difference between improving a world model as a predictor and improving it as a basis for safe action. If the predictive representation collapses states that the decision criterion must distinguish, better prediction within that representation cannot solve the decision problem.

The distinction becomes more consequential when intervention changes the distribution being predicted. Passive prediction concerns a conditional distribution such as
[
P(Y\mid X=x).
]
Intervention concerns a quantity closer to
[
P(Y\mid \operatorname{do}(A=a)),
]
or another counterfactual object describing what would occur under an action. Observational association and interventional consequence coincide only under additional assumptions. A model can therefore be highly successful at predicting what normally follows observed states while remaining inadequate for estimating what would happen if the agent deliberately changed the process.

This gives another scope boundary:
[
\text{observation}
\longrightarrow
\text{intervention}.
]
The transition is not licensed by predictive accuracy alone. It requires assumptions concerning causal structure, confounding, stability, and the consequences of manipulating variables rather than merely observing them.

For an embodied system, the distinction is unavoidable because actions are interventions by definition. A world model used only to anticipate passive sensory input can be evaluated as a forecasting device. A world model used to choose actions must support counterfactual distinctions among candidate interventions.

Let
[
\mathcal{M}{\mathrm{obs}}
]
denote the set of models compatible with the available observational data. It is possible for two models
[
M_1,M_2\in\mathcal{M}{\mathrm{obs}}
]
to agree on all observational predictions while disagreeing on an intervention:
[
P_{M_1}(Y\mid X)

P_{M_2}(Y\mid X)
]
but
[
P_{M_1}(Y\mid\operatorname{do}(A=a))
\neq
P_{M_2}(Y\mid\operatorname{do}(A=a)).
]
If the evidence available to the system cannot distinguish $M_1$ from $M_2$, then observational predictive success does not identify the consequence of intervention.

\begin{proposition}[Observational equivalence can conceal interventional disagreement]
If two environment models are observationally equivalent under the verifier's evidence but assign different consequences to an intervention $a$, then observational evidence alone is insufficient to determine the interventional consequence of $a$.
\end{proposition}

\begin{proof}
Let the observational projection of model $M$ be
[
\pi_{\mathrm{obs}}(M).
]
Suppose
[
\pi_{\mathrm{obs}}(M_1)

\pi_{\mathrm{obs}}(M_2),
]
while
[
I_a(M_1)\neq I_a(M_2),
]
where $I_a(M)$ denotes the predicted consequence of intervention $a$. Then $I_a$ is not constant on the fibers of $\pi_{\mathrm{obs}}$. By the fiber criterion, $I_a$ does not factor through $\pi_{\mathrm{obs}}$. Therefore no procedure receiving only the observational evidence can universally determine the interventional consequence.
\end{proof}

This result connects the epistemic case directly to the paper's general framework. The problem is not that the system has failed to calculate cleverly enough. The evidence supplied to it is compatible with multiple interventionally distinct worlds.

The distinction between prediction and intervention leads naturally to the distinction between uncertainty and warrant. Suppose a system assigns a probability distribution over outcomes for each action. It may still lack sufficient evidence to justify selecting among those actions. The question is then not simply which action maximizes a current objective, but whether the evidential state itself is adequate for committing to an environment-changing action.

Let
[
E_t
]
denote the evidence available at time $t$, and let
[
\Gamma
]
denote a governing criterion specifying conditions under which an intervention is admissible. Define an admissible-action set
[
\mathcal{A}_{\Gamma}(E_t)

\left{
a\in\mathcal{A}:
\Gamma(E_t,a)=1
\right}.
]
Even if a predictive model ranks an action $a^\star$ highest,
[
a^\star
\in
\arg\max_{a\in\mathcal{A}}
Q(a\mid E_t),
]
it does not follow that
[
a^\star\in\mathcal{A}_{\Gamma}(E_t).
]
Optimization and admissibility are separate operations.

Indeed, the admissible set can be empty:
[
\mathcal{A}_{\Gamma}(E_t)=\varnothing.
]
A system whose architecture requires selection of some
[
a\in\mathcal{A}
]
cannot represent this epistemic state faithfully. It must act even when its own governing criterion licenses no available intervention.

This motivates extending the action space. Let
[
\mathcal{A}^{+}

\mathcal{A}
\cup
{\text{sense},\text{revise},\text{defer},\text{abstain}}.
]
The added operations are not ordinary substitutes for environment-changing actions. They alter the epistemic or computational state of the system, preserve option value, or decline commitment.

The distinction can be represented by partitioning actions into
[
\mathcal{A}{\mathrm{world}}
]
and
[
\mathcal{A}{\mathrm{epistemic}},
]
where the former primarily changes the external environment and the latter primarily changes the information available to the decision process. The categories need not be perfectly disjoint in physical systems, since sensing itself can perturb an environment, but the conceptual distinction remains useful.

A sensing operation
[
q\in\mathcal{A}{\mathrm{epistemic}}
]
can transform
[
E_t
\longrightarrow
E{t+1}
]
without committing to the world-changing action under uncertainty. If the additional evidence changes the admissible set,
[
\mathcal{A}{\Gamma}(E_t)
\neq
\mathcal{A}{\Gamma}(E_{t+1}),
]
then sensing has decision value independent of any immediate environmental reward.

This can be formalized through the value of information. Let
[
V(E)

\max_{a\in\mathcal{A}_{\Gamma}(E)}
\mathbb{E}[U(a)\mid E],
]
with an appropriate convention when the admissible set is empty. If a sensing action produces possible future evidence states $E'$ according to
[
P(E'\mid E,q),
]
then the expected informational value of $q$ can be represented as
[
\operatorname{VOI}(q\mid E)

\mathbb{E}_{E'}
[V(E')]

V(E)

C(q),
]
where $C(q)$ is the cost of obtaining the information. When this quantity is positive, immediate intervention can be inferior to acquiring evidence even if a current predictive ranking exists.

This exposes a limitation of architectures that collapse every uncertainty state into a ranking over ordinary actions. A ranking answers
[
\text{which action is best under the current model?}
]
It does not answer
[
\text{is the current model sufficiently discriminating to justify acting?}
]
The latter question concerns the adequacy of the evidence, not merely the relative values produced from it.

One can make the distinction exact through another projection argument. Let
[
r(E)
]
be the ranking over candidate actions induced by evidence $E$. Suppose there exist evidence states $E_1,E_2$ satisfying
[
r(E_1)=r(E_2)
]
but differing in whether intervention is warranted:
[
W(E_1)=1,
\qquad
W(E_2)=0.
]
Then warrant does not factor through the action ranking. No policy receiving only the ranking can distinguish a state in which it should act from a state in which it should defer.

This can occur naturally. Two evidence states may rank
[
a_1>a_2>a_3
]
identically while differing dramatically in confidence. In one state,
[
P(a_1\text{ is preferable}\mid E_1)=0.999,
]
while in another,
[
P(a_1\text{ is preferable}\mid E_2)=0.34,
]
with the remaining probability nearly evenly distributed. The ordinal ranking is the same. The epistemic situation is not.

Consequently,
[
\text{action ordering}
\quad\not\Rightarrow\quad
\text{warrant to act}.
]

The same reasoning applies to scalar confidence if the governing criterion depends upon more than confidence. Two decisions can have identical confidence values while differing in reversibility, consequence severity, availability of further evidence, or distributional uncertainty. A scalar confidence score can therefore be another insufficient projection.

Ma et al.'s emphasis on recoverability is particularly important in this context because intervention decisions should depend not only upon predicted immediate outcomes but also upon the possibility of correcting an error. Let
[
C(a)
]
denote the expected consequence of an action and
[
R(a)
]
a measure of recoverability. Two actions can have identical expected immediate consequences while differing substantially in
[
R(a).
]
If the decision criterion depends upon recoverability but the world-model representation preserves only expected consequence, then the representation is insufficient for the decision property.

A simple model makes this clear. Let
[
a_1,a_2
]
both have expected loss
[
\mathbb{E}[L\mid a_1]

\mathbb{E}[L\mid a_2]

\ell.
]
Suppose, however, that failures following $a_1$ can be reversed at negligible cost while failures following $a_2$ are irreversible. A decision rule sensitive to irreversibility can prefer $a_1$ even though a representation containing only expected loss makes the actions identical. The missing distinction is not predictive accuracy about the expected value. It is the structure of the consequence space.

Finite-horizon prediction creates an analogous insufficiency. Let a model predict states up to horizon $k$:
[
\pi_k(H)

(s_{t+1},\ldots,s_{t+k}).
]
Suppose the target property concerns a cumulative quantity
[
C(H)

\sum_{j=1}^{\infty}
\gamma^{j-1}c(s_{t+j})
]
or a threshold crossing occurring after time $t+k$. If there exist two trajectories identical through the first $k$ future states but differing afterward,
[
\pi_k(H_1)=\pi_k(H_2)
]
while
[
C(H_1)\neq C(H_2),
]
then $C$ does not factor through the finite-horizon prediction.

Extending $k$ can reduce the problem, but no finite increase solves a property with genuinely unbounded dependence unless some sufficient summary exists. As in the temporal section, the correct objective is not horizon maximization but preservation of the distinctions required by the property.

Suppose there exists a terminal statistic
[
z_{t+k}
]
such that all consequences beyond the horizon relevant to $C$ are determined by
[
z_{t+k}.
]
Then the representation
[
(s_{t+1},\ldots,s_{t+k},z_{t+k})
]
may be sufficient even when the raw $k$-step trajectory is not. This again demonstrates that representation design can matter more than sheer horizon length.

The problem of accumulated consequences is particularly acute when individually small effects cross a threshold only through repetition. Let each action contribute
[
c_t\geq 0
]
to a cumulative burden
[
B_T

\sum_{t=1}^{T}c_t.
]
Suppose every local action satisfies
[
c_t<\tau.
]
It does not follow that
[
B_T<\tau.
]
Indeed, for any positive $c_t=c$, one can choose
[
T>\frac{\tau}{c}
]
so that
[
B_T>\tau.
]
Local boundedness and cumulative boundedness are different properties.

If the safety criterion concerns
[
B_T\leq B_{\max},
]
then verifying
[
c_t\leq c_{\max}
]
at every step is sufficient only when one also has a composition bound such as
[
T c_{\max}\leq B_{\max},
]
or a stronger invariant controlling accumulation. Once again, the missing relation is a composition theorem.

This is the epistemic analogue of the temporal and structural cases. A locally acceptable predicted consequence can accumulate into an unacceptable history. The distinction is especially important when effects are delayed, path-dependent, or irreversible.

The four dimensions emphasized by Ma et al. can therefore be understood as different ways in which a predictive representation can fail to support a decision property. Likelihood may omit consequence severity. Passive prediction may omit intervention effects. A finite horizon may omit accumulation. A point prediction may omit epistemic uncertainty. A predicted consequence may omit recoverability. Each omission induces equivalence classes that can identify situations a safe decision rule must distinguish.

The governing question is consequently not whether a world model is ``accurate'' in the abstract. Accuracy is always accuracy with respect to a target. A model can be highly accurate for image prediction and insufficient for collision avoidance, accurate for expected state transition and insufficient for tail risk, accurate for passive observation and insufficient for intervention, or accurate for short-horizon dynamics and insufficient for accumulated harm.

This suggests distinguishing predictive calibration from decisional calibration. Let
[
M
]
produce probabilities over outcomes. Predictive calibration concerns whether events assigned probability $p$ occur with appropriate frequency under the relevant conditions. Decisional calibration concerns whether the information supplied to a decision rule supports interventions at an appropriate evidential threshold given their consequences. Predictive calibration is valuable but does not entail decisional calibration because the latter depends upon utilities, risk, uncertainty, reversibility, and alternative information-gathering actions.

One can make this distinction especially sharp by considering a perfectly calibrated but uninformative predictor. Suppose
[
M(Y=1\mid x)=0.5
]
for every $x$, and suppose the true conditional probability is indeed $0.5$ everywhere. The model is perfectly calibrated. Yet if intervention requires distinguishing cases in which an action is safe from cases in which it is dangerous, the predictor may provide no useful decision discrimination at all. Calibration does not imply resolution.

Conversely, a model may have strong discrimination but poor calibration. The two properties answer different questions. Neither by itself determines whether an intervention is warranted.

The epistemic problem therefore introduces a second-order verification question:
[
\text{Has the system verified enough to act?}
]
This differs from
[
\text{Has the system verified its prediction?}
]
A prediction can be well supported while the action remains underdetermined because the decision requires distinctions absent from the predictive representation.

This motivates an explicit null or deferential branch in consequential systems. Let the available response set be
[
\mathcal{U}

\mathcal{A}{\mathrm{world}}
\cup
\mathcal{A}{\mathrm{epistemic}}
\cup
{\varnothing},
]
where
[
\varnothing
]
denotes no intervention. The null action should not be treated as a failure to choose. It is an ordinary member of the decision architecture whenever the expected value of intervention, after accounting for uncertainty and consequence, does not exceed the value of preserving the current state.

Suppose a decision functional is
[
J(a\mid E)

C_{\mathrm{act}}(a)
+
C_{\mathrm{risk}}(a\mid E)
+
C_{\mathrm{irr}}(a\mid E)

B(a\mid E),
]
where the terms represent action cost, risk cost, irreversibility cost, and expected benefit. Define
[
J(\varnothing\mid E)=0.
]
Then intervention is justified under this model only if
[
\min_{a\in\mathcal{A}_{\mathrm{world}}}
J(a\mid E)<0.
]
If
[
J(a\mid E)\geq 0
]
for every environment-changing action, the null action is optimal.

The mathematical importance of explicitly representing $\varnothing$ is that the system is no longer forced to compare interventions only against one another. It compares intervention against nonintervention. Without the null branch, an optimizer can select the least bad intervention even when every intervention is worse than doing nothing.

Similarly, if an epistemic action $q$ can improve the evidence state at acceptable cost, one may have
[
J(q\mid E)
<
J(a\mid E)
]
for every immediate world-changing action $a$. The correct next operation is then to obtain information rather than to alter the world.

This structure captures Ma et al.'s emphasis on sensing, revision, deferment, and abstention in a general form. These are not peripheral fallback mechanisms appended after the ``real'' decision process. They are part of the action space required to represent epistemic insufficiency honestly.

The result also prevents an important category error. A system that abstains because its evidence is insufficient has not necessarily failed to solve the prediction problem. It may have solved the prediction problem correctly and discovered that the resulting information does not determine an intervention. Refusal can therefore be the output of successful reasoning rather than evidence that reasoning stopped prematurely.

The epistemic failure of composition can now be stated precisely. Predictive success does not automatically compose into decisional sufficiency. A probability distribution over future states is not a permission to produce one of them. A ranking among actions is not evidence that any action is sufficiently warranted. A finite-horizon forecast is not a guarantee about accumulated consequences. A high-confidence prediction is not a measure of recoverability. A well-calibrated model is not necessarily a decision-sufficient model.

Each inference requires an additional relation. If a decision $D$ is to be derived from a predictive representation $\pi_{\mathrm{pred}}$, then one requires a factorization
[
D

\widehat{D}\circ\pi_{\mathrm{pred}}.
]
If no such factorization exists, then the predictive representation is not the correct unit for the decision property, however accurate it may be for its own target.

This is the fourth version of the paper's central argument. The temporal case showed that a present action can omit distinctions encoded in its history. The structural case showed that independent events can omit distinctions encoded in their arrangement. The representational case showed that a transformed object can omit distinctions present in its source. The epistemic case now shows that a prediction can omit distinctions required to determine whether intervention is justified.

The recurring mistake is to treat success on the upstream task as though it established the downstream property. Yet each downstream inference changes the question. From action to history, from events to structure, from source to representation, and from prediction to intervention, the relevant property must survive a transformation. Where preservation has not been established, verification stops at the boundary.

A safety-critical world model should therefore be evaluated not only by asking whether it predicts well, but by asking which decisions its representation can support without identifying decision-distinct situations. The crucial test is not whether the model contains much information. It is whether the property governing intervention factors through the information the model actually provides.

The checkpoint, in this case, is prediction itself. A prediction can be accurate, calibrated, detailed, and still be the wrong boundary at which to declare the decision settled. The fact that the future has been estimated does not mean that the right to alter it has been established.

\section{Four Failures of an Implicit Preservation Law}

The preceding cases concern different systems, different empirical phenomena, and different mathematical objects. Long-horizon interaction is not workflow topology. Workflow topology is not semantic translation. Semantic translation is not intervention under uncertainty. It would therefore be a mistake to explain all four cases through a single causal mechanism. Their commonality lies elsewhere. In each case, a property established on one object is expected to survive a change in representation, scale, composition, or epistemic role. The recurring error is not merely local reasoning. It is the unproved transport of a conclusion across a transformation.

The temporal case begins with actions and asks about histories. The structural case begins with events and asks about their relational composition. The representational case begins with a source object and asks whether a property established after translation reflects back to that source. The epistemic case begins with a predictive representation and asks whether it contains sufficient information to license intervention. These can be written schematically as
[
\begin{aligned}
\text{temporal:}\qquad
&{a_t}_{t=0}^{n-1}
\longrightarrow
H,\
\text{structural:}\qquad
&{v,e}
\longrightarrow
G,\
\text{representational:}\qquad
&x
\xrightarrow{\tau}
y,\
\text{epistemic:}\qquad
&M
\longrightarrow
D.
\end{aligned}
]
The arrows do not denote the same kind of operation. The first is temporal composition, the second relational assembly, the third representational transformation, and the fourth conversion of predictive information into a decision. What matters is that in every case the property on one side of the arrow is not automatically a property on the other.

The common structure can be abstracted. Let
[
X
]
be a source domain and
[
Y
]
a target domain. Let
[
T:X\rightarrow Y
]
be a transformation, broadly understood. Suppose a verifier establishes a predicate
[
V:X\rightarrow{0,1}
]
on the source side, while the desired conclusion concerns a predicate
[
\Phi:Y\rightarrow{0,1}
]
on the target side. The inference
[
V(x)=1
\Longrightarrow
\Phi(T(x))=1
]
is valid only when the triple
[
(V,T,\Phi)
]
satisfies an appropriate preservation condition.

\begin{definition}[Forward preservation]
A transformation
[
T:X\rightarrow Y
]
is \emph{forward-preserving from $V$ to $\Phi$} if
[
V(x)=1
\Longrightarrow
\Phi(T(x))=1
]
for every $x\in X$.
\end{definition}

This definition is deliberately relational. Preservation is not a property of $T$ alone. A transformation can preserve one predicate and fail to preserve another. Likewise, the same predicate can be preserved by one transformation and destroyed by another.

The simplest composition theorem is immediate.

\begin{theorem}[Verified transport]
Let
[
V:X\rightarrow{0,1},
\qquad
T:X\rightarrow Y,
\qquad
\Phi:Y\rightarrow{0,1}.
]
If $T$ is forward-preserving from $V$ to $\Phi$, then
[
V(x)=1
]
licenses
[
\Phi(T(x))=1.
]
If forward preservation has not been established, the validity of $V(x)$ alone supplies no general implication concerning $\Phi(T(x))$.
\end{theorem}

The first statement follows directly from the definition. The second is the substantive methodological point. A verification result has a domain. Transporting that result to another domain requires a bridge. The bridge may be a theorem, a specification, a probabilistic calibration result, a causal assumption, or an empirically validated relation, but it cannot be replaced by the mere fact that the original verification was rigorous.

This distinction also prevents an ambiguity in the phrase ``verification does not compose.'' Verification can compose. Formal methods contain many successful examples in which local invariants, typing judgments, refinement relations, or assume-guarantee contracts yield system-level results. The stronger and more accurate claim is
[
\boxed{
\text{verification does not compose automatically.}
}
]
Composition is an additional property to be established.

For a history
[
H=(s_0,a_0,\ldots,s_n),
]
let
[
V_t(s_t,a_t,s_{t+1})
]
be a family of local predicates. Define the aggregate local-verification predicate
[
L(H)

\bigwedge_{t=0}^{n-1}
V_t(s_t,a_t,s_{t+1}).
]
Then a global property
[
\Phi:\Hist\rightarrow{0,1}
]
follows from local verification exactly when
[
L(H)=1
\Longrightarrow
\Phi(H)=1
]
holds over the relevant class of histories.

The implication itself is the composition theorem.

\begin{definition}[Verification composition law]
A family of local predicates ${V_t}$ is sufficient for a history property $\Phi$ over a history class $\mathcal{C}\subseteq\Hist$ if
[
\forall H\in\mathcal{C},
\qquad
\left[
\bigwedge_t V_t(H)=1
\right]
\Longrightarrow
\Phi(H)=1.
]
\end{definition}

The restriction to a history class $\mathcal{C}$ matters. A composition law can be valid only under environmental, architectural, or adversarial assumptions. For example, a collection of local access-control checks may imply a system-level confidentiality property only when there are no unmodeled communication channels. A local memory-safety property may compose under a language semantics that excludes certain forms of undefined behavior. A distributed protocol proof may rely on bounds concerning faults or message delivery. The theorem is therefore properly understood as conditional:
[
\mathcal{A}
\land
\bigwedge_t V_t
\Longrightarrow
\Phi,
]
where $\mathcal{A}$ collects the assumptions under which composition has been proved.

Omitting $\mathcal{A}$ is another way in which a verification result can silently expand. The proof may be sound within its model while the deployed conclusion is asserted outside the model's admissible domain. Scope concerns assumptions as well as observations.

The four empirical cases can now be restated as challenges to different implicit composition laws. In the temporal case, the tempting inference is
[
\forall t,;
V(a_t)=1
\Longrightarrow
\Phi(H)=1.
]
History-dependent coordination provides counterexamples when $\Phi$ depends upon relations among actions or upon the policy induced by prior interaction. In the structural case, the tempting inference is
[
\forall c\in C(G),;
V(c)=1
\Longrightarrow
\Phi(G)=1,
]
where $C(G)$ denotes locally inspected components. Relational graph properties supply counterexamples whenever incidence matters. In the representational case, the tempting inference is
[
P_Y(\tau(x))=1
\Longrightarrow
P_X(x)=1.
]
The implication requires a reflection theorem for $\tau$. In the epistemic case, the tempting inference is
[
M(\cdot\mid s,a)\text{ is accurate}
\Longrightarrow
\operatorname{Choose}(a),
]
which requires additional decision-relevant information and governing criteria.

The four failures can therefore be summarized as
[
\begin{aligned}
\text{history dependence:}\qquad
&a_t\not!\perp H_{<t},\
\text{structural dependence:}\qquad
&\Phi(G)\neq
f\bigl({\phi(c)}_{c\in C(G)}\bigr)
\quad\text{in general},\
\text{representational dependence:}\qquad
&P_Y(\tau(x))
\not\Rightarrow
P_X(x),\
\text{decisional dependence:}\qquad
&P(s'\mid s,a)
\not\Rightarrow
\operatorname{Choose}(a).
\end{aligned}
]
These expressions should not be read as universal impossibility statements. Each indicates the absence of an implication without further conditions.

The common failure is an implicit preservation law. Something is verified at one scale, and a conclusion is asserted at another. The missing theorem is whatever would make the relevant diagram commute.

Consider the generic diagram
[
\begin{array}{ccc}
X & \xrightarrow{T} & Y\
V\downarrow & & \downarrow\Phi\
{0,1} & \xrightarrow{g} & {0,1}.
\end{array}
]
If there exists a map
[
g:{0,1}\rightarrow{0,1}
]
such that
[
\Phi\circ T

g\circ V,
]
then the target property is determined by the verified source predicate. In the simplest case,
[
g=\operatorname{id},
]
so
[
\Phi(T(x))=V(x).
]
More commonly, exact commutation is unnecessarily strong. One may need only the implication
[
V(x)=1
\Longrightarrow
\Phi(T(x))=1.
]
The relevant preservation requirement depends upon the conclusion being drawn.

This distinction between exact preservation and one-directional preservation is important because verification tasks are asymmetric. Certification requires that accepted transformed objects reflect source correctness. Bug detection may require that a discovered counterexample correspond to a genuine source failure. Safety monitoring may prioritize avoiding false negatives, while anomaly triage may tolerate uncertainty in both directions. The preservation law should therefore be stated in the direction actually required rather than hidden inside a vague appeal to ``equivalence.''

The factorization framework introduced earlier addresses a related but different question. Preservation asks whether a verified property survives a transformation. Factorization asks whether a target property can be determined from a representation at all.

Let
[
\pi:X\rightarrow O
]
be an observational representation and
[
\Phi:X\rightarrow Z
]
a target property. If
[
\Phi

\widehat{\Phi}\circ\pi,
]
then $\pi$ is sufficient support for $\Phi$. If no such $\widehat{\Phi}$ exists, then no verifier restricted to $\pi(x)$ can determine $\Phi(x)$ universally.

Preservation and factorization should not be conflated. Preservation concerns movement of a claim:
[
V(x)
\Longrightarrow
\Phi(T(x)).
]
Factorization concerns recoverability of a property:
[
\Phi(x)

\widehat{\Phi}(\pi(x)).
]
The former is inferential. The latter is informational.

The two interact. If the target property does not factor through the representation supplied to the verifier, then no preservation theorem based solely on that representation can recover the target property universally. Conversely, factorization establishes only that the property is in principle recoverable. It does not prove that the actual verifier computes the correct $\widehat{\Phi}$.

This yields a useful decomposition of verification adequacy.

\begin{definition}[Verification adequacy]
A verifier operating on a representation
[
\pi:X\rightarrow O
]
is adequate for a property
[
\Phi:X\rightarrow Z
]
when both of the following hold. First, $\Phi$ factors through $\pi$:
[
\Phi=\widehat{\Phi}\circ\pi
]
for some $\widehat{\Phi}$. Second, the verifier correctly computes the required $\widehat{\Phi}$ on the relevant domain.
\end{definition}

The first condition is representational sufficiency. The second is computational correctness. These are independent failure modes.

A verifier can fail computationally despite receiving sufficient information. It may be poorly trained, incorrectly implemented, underpowered, or statistically miscalibrated. Conversely, a verifier can be computationally optimal relative to its input and still fail because its input is insufficient.

This distinction can be represented as
[
\text{verification failure}

\text{representation failure}
;\lor;
\text{inference failure},
]
where the disjunction is conceptual rather than probabilistic. The first occurs when the target property does not factor through the verifier's input. The second occurs when factorization exists but the verifier fails to implement the corresponding map correctly.

The distinction has practical consequences for evaluation. If a verifier fails on a benchmark, improving the verifier is appropriate only when the benchmark information available at its interface is sufficient for the target labels. If two differently labeled cases are indistinguishable at the verifier's input, the benchmark is asking the verifier to recover information it has not been given. More model capacity cannot solve the resulting collision.

This suggests a diagnostic procedure for verification systems. Before measuring classifier performance, one can ask whether the representation induces collisions across target-property classes. For a finite dataset
[
{(x_i,\Phi(x_i))}_{i=1}^N,
]
if there exist $i\neq j$ such that
[
\pi(x_i)=\pi(x_j)
]
but
[
\Phi(x_i)\neq\Phi(x_j),
]
then perfect accuracy is impossible for any deterministic verifier operating only on $\pi(x)$. With stochastic verifiers the same collision imposes a nonzero irreducible error under any distribution assigning positive mass to both cases.

The finite result is immediate.

\begin{proposition}[Collision lower bound]
Suppose two source cases $x_1,x_2$ occur with probabilities $p_1,p_2>0$, satisfy
[
\pi(x_1)=\pi(x_2),
]
and have opposite binary labels:
[
\Phi(x_1)\neq\Phi(x_2).
]
Then every verifier whose output depends only upon $\pi(x)$ has error probability at least
[
\min(p_1,p_2)
]
on the distribution restricted to these two cases.
\end{proposition}

\begin{proof}
Because
[
\pi(x_1)=\pi(x_2),
]
the verifier receives the same input for both cases. A deterministic verifier must therefore assign the same label to both and must be wrong on at least one, incurring probability at least
[
\min(p_1,p_2).
]
A randomized verifier assigns the same output distribution to both cases. If it assigns probability $q$ to the label of $x_1$, its expected error on the pair is
[
p_1(1-q)+p_2q.
]
This affine function is minimized at $q=1$ when $p_1>p_2$ and at $q=0$ when $p_2>p_1$, with minimum
[
\min(p_1,p_2).
]
\end{proof}

The proposition gives projection insufficiency a statistical form. Representation-induced collisions create an irreducible error floor. The appropriate remedy is to refine the representation, alter the target property, or accept uncertainty. Training the same verifier harder cannot remove an information-theoretic collision.

The preservation problem becomes more complex when a pipeline contains multiple transformations. Contemporary agentic systems rarely contain a single boundary. An instruction may be summarized, translated into a plan, decomposed into tool calls, routed through multiple agents, transformed into external operations, observed through sensors, compressed back into memory, and then evaluated. Let
[
X_0
\xrightarrow{T_1}
X_1
\xrightarrow{T_2}
X_2
\xrightarrow{T_3}
\cdots
\xrightarrow{T_m}
X_m
]
be such a pipeline.

Suppose each space $X_i$ carries a property
[
P_i:X_i\rightarrow{0,1}.
]
If each transformation preserves the relevant property in the required direction,
[
P_{i-1}(x)=1
\Longrightarrow
P_i(T_i(x))=1,
]
then preservation composes.

\begin{theorem}[Transitivity of preservation]
Suppose for every
[
i\in{1,\ldots,m},
]
the transformation
[
T_i:X_{i-1}\rightarrow X_i
]
satisfies
[
P_{i-1}(x)=1
\Longrightarrow
P_i(T_i(x))=1.
]
Then
[
P_0(x)=1
\Longrightarrow
P_m\bigl(T_m\circ\cdots\circ T_1(x)\bigr)=1.
]
\end{theorem}

\begin{proof}
Apply the preservation implication successively. From
[
P_0(x)=1
]
and preservation of $T_1$,
[
P_1(T_1(x))=1.
]
Preservation of $T_2$ then yields
[
P_2(T_2(T_1(x)))=1.
]
Continuing inductively through $T_m$ gives
[
P_m(T_m\circ\cdots\circ T_1(x))=1.
]
\end{proof}

This theorem states the positive counterpart to the paper's negative examples. Verification can travel arbitrarily far through a pipeline when every boundary has an appropriate preservation law. The problem is not distance. It is the presence of an unverified crossing.

A single missing link breaks the chain of justified inference. Suppose preservation is known for
[
T_1,\ldots,T_{k-1}
]
and
[
T_{k+1},\ldots,T_m
]
but not for $T_k$. Then the complete implication from $P_0$ to $P_m$ is not established. Strong guarantees before and after the gap do not fill it.

This is particularly relevant to hybrid neuro-symbolic architectures. A formal verifier may provide a strong guarantee at $X_m$, while a learned transformation earlier in the pipeline remains unverified. The rigor of later links does not propagate backward through the missing preservation law. Similarly, a carefully verified tool call does not establish a history-level property when the missing link lies between the tool call and the larger trajectory.

The chain model also reveals why adding more checkpoints does not monotonically increase the scope of what has been established. Suppose each transformation is followed by a verifier:
[
X_0
\xrightarrow{T_1}
X_1
\xrightarrow{V_1}
\cdots
\xrightarrow{T_m}
X_m
\xrightarrow{V_m}.
]
Even if
[
V_i(x_i)=1
]
for every checkpoint, the final target property $\Phi(X_0,\ldots,X_m)$ need not follow unless the conjunction of checkpoint predicates is sufficient for $\Phi$.

Formally, define
[
C(x_0,\ldots,x_m)

\bigwedge_{i=1}^{m}V_i(x_i).
]
The desired inference is
[
C=1
\Longrightarrow
\Phi=1.
]
Adding another local verifier replaces $C$ by
[
C'

C\land V_{m+1}.
]
This can reduce the set of trajectories that pass, but it does not establish $\Phi$ unless
[
C'\Rightarrow\Phi
]
has been shown. More checkpoints can strengthen filtering while leaving the fundamental composition gap untouched.

This gives a distinction between \emph{verification density} and \emph{verification scope}. Verification density concerns how many points in a process are checked. Verification scope concerns which properties those checks jointly establish. Increasing density does not necessarily increase scope.

A system can therefore be densely verified and globally under-specified. Every transition can carry a certificate while the history-level property remains unproved. Conversely, a sparse verification architecture can establish a strong global property when its checks are embedded in a compositional proof.

This distinction is useful because engineering discussions often equate continuous monitoring with stronger assurance. Continuous monitoring provides more observations. Whether those observations establish the desired property depends upon their informational and inferential relation to that property.

The same principle applies to redundancy. Suppose $k$ independent verifiers all inspect the same insufficient projection
[
\pi(x).
]
Their agreement can reduce errors arising from imperfect inference, but it cannot recover distinctions absent from $\pi$. If
[
\pi(x_1)=\pi(x_2)
]
while
[
\Phi(x_1)\neq\Phi(x_2),
]
then every verifier receives the same representation for both cases. An ensemble can disagree internally because of stochasticity or differing inductive biases, but no ensemble function of $\pi(x)$ alone can universally determine $\Phi$.

\begin{corollary}[Redundancy does not repair projection insufficiency]
Let
[
f_1,\ldots,f_k
]
be arbitrary verifiers whose outputs depend only upon $\pi(x)$, and let
[
A
]
be any aggregation rule over their outputs. Then
[
A(f_1(\pi(x)),\ldots,f_k(\pi(x)))
]
cannot universally determine $\Phi(x)$ if $\Phi$ does not factor through $\pi$.
\end{corollary}

The proof follows because the entire ensemble is itself a function of $\pi(x)$. Projection insufficiency applies unchanged.

This observation is consequential for architectures that attempt to solve verification failures through debate, voting, critique, or verifier ensembles. Such mechanisms can improve inference over information that is present. They cannot reconstruct distinctions systematically erased before all participants receive their inputs. Diversity of reasoning does not substitute for diversity of evidence.

The same limitation applies to repeated verification over time when each check observes the same insufficient statistic. If the sequence of verifier inputs
[
\pi_1(H),\ldots,\pi_n(H)
]
collectively identifies two histories that differ on $\Phi$, then the entire monitoring process remains insufficient. What matters is the joint projection
[
\Pi(H)

(\pi_1(H),\ldots,\pi_n(H)).
]
The property is recoverable exactly when
[
\Phi

\widehat{\Phi}\circ\Pi
]
for some $\widehat{\Phi}$.

Thus a sequence of individually insufficient observations can become jointly sufficient, but only if their combination separates the property-relevant cases. This gives a positive account of why monitoring over time can help. The benefit is not temporal persistence by itself. It is refinement of the induced partition over possible histories.

Suppose
[
\Pi_n(H)

(\pi_1(H),\ldots,\pi_n(H)).
]
Then
[
\Pi_{n+1}\succeq\Pi_n
]
because $\Pi_n$ can be recovered by discarding the final coordinate. The sequence of joint observations becomes monotonically finer in the informational ordering. Consequently, once a property factors through $\Pi_n$, it also factors through every later $\Pi_m$ with $m>n$, assuming the observations are retained.

This monotonicity of information should not be confused with monotonicity of decisions. Additional evidence can overturn an earlier conclusion. If a verifier acts on partial information, one can have
[
D(\Pi_n(H))=1
]
and
[
D(\Pi_{n+1}(H))=0.
]
The information set has become finer while the warranted decision has reversed. More information is monotonic as refinement, but warrant need not be monotonic.

This distinction becomes important in systems that treat a prior verification decision as permanently settled. A certificate may remain valid if it concerns an immutable mathematical object. A decision about whether to continue an evolving process may cease to be justified when new evidence arrives. Verification results therefore inherit not only a scope but sometimes a temporal validity condition.

The preservation framework can accommodate this by indexing properties to evidence:
[
\Phi(H,E_t).
]
A result established at time $t$ need not imply the corresponding result at $t+1$:
[
\Phi(H,E_t)=1
\quad\not\Rightarrow\quad
\Phi(H,E_{t+1})=1.
]
The target proposition has changed because its evidential argument has changed.

This further clarifies why verification should not be treated as a single terminal event in agentic systems. Some properties are static and certificate-like. Others are conditional, revisable, and dependent upon an evolving history. The unit of verification must therefore include not only the object being judged but, where relevant, the evidence state relative to which the judgment is valid.

The synthesis can now be expressed through three separate questions. The first is informational:
[
\text{Does the target property factor through the representation?}
]
The second is computational:
[
\text{Does the verifier correctly compute the property from that representation?}
]
The third is compositional:
[
\text{Does the verified property survive the transformation or composition across which it is being transported?}
]
Failure at any one of these stages limits the conclusion.

These questions correspond to three distinct proof obligations:
[
\boxed{
\text{support}
\quad+\quad
\text{correct inference}
\quad+\quad
\text{preservation}.
}
]
Support concerns whether the necessary distinctions reach the verifier. Correct inference concerns whether the verifier uses those distinctions properly. Preservation concerns whether the established result remains valid when moved to the object or scale about which the final claim is made.

The four case studies occupy different locations in this decomposition. Long-horizon interaction emphasizes support across time and the causal role of retained history. Workflow anomaly detection emphasizes support for relational structure and the gap between diagnosis and judgment. Neuro-formal verification emphasizes preservation across representation. Safety-critical world models emphasize the transition from predictive support to decisional support. None is reducible to the others, but all become legible within the same account of verification scope.

The central negative result can therefore be sharpened. It is not merely
[
\bigwedge_t V_t=1
\not\Rightarrow
\Phi(H)=1.
]
It is that every inference extending a verification result beyond its immediate domain incurs a preservation obligation. If the extension also passes through an informational projection, it incurs a support obligation. If a verifier must infer the property from that support, it incurs a correctness obligation.

The complete structure can be written schematically as
[
x
\xrightarrow{\pi}
o
\xrightarrow{\widehat{\Phi}}
z
\xrightarrow{T}
z',
]
with separate questions at every arrow. Did $\pi$ preserve the distinctions required by the property? Did the verifier correctly compute $\widehat{\Phi}$? Does the resulting property survive the subsequent transformation $T$? A system can answer two of these questions correctly and still fail on the third.

This is why a checkpoint can be perfectly implemented and still support an invalid conclusion. The defect need not be located at the checkpoint itself. The checkpoint may receive an insufficient representation, or its valid result may later be transported across an unproved transformation. Verification errors can therefore occur before, at, or after the visible act of verification.

The appropriate response is not skepticism toward verification. It is a more demanding account of what verification establishes. A proof should be allowed to say exactly as much as its premises, representation, and preservation theorems support. No less, because local and formal verification remain powerful. No more, because rigor at one boundary cannot substitute for an argument about what survives beyond it.

The general principle can now be stated in its strongest form:
[
\boxed{
\begin{minipage}{0.82\linewidth}
A verification result is portable across a change of scope only to the extent that the property it establishes is supported by the representation at the new scope and preserved by the transformation connecting the two.
\end{minipage}
}
]

The remainder of the argument turns from failure to design. If verification units are property-relative, then the task is not to find one universally correct scale of inspection. It is to characterize the informational support of the property, determine which representations preserve that support, and identify minimal or otherwise efficient verification units relative to explicit architectural constraints.

\section{The Unit Is Determined by the Property}

The preceding analysis rules out a simple replacement for local verification. If output-level verification can be insufficient, it does not follow that trajectory-level verification is universally sufficient. If eventwise inspection can miss workflow structure, it does not follow that the complete workflow graph is the universally correct verification object. If a translated program can omit source-level distinctions, retaining more syntax is useful only when the retained distinctions matter to the property under consideration. If predictive representations can be insufficient for intervention, adding more predictions does not solve the problem unless the additional information separates decision-relevant cases. There is therefore no universally privileged scale at which verification should occur.

The appropriate unit is determined by the property.

This statement can now be made exact. Let
[
\mathcal{X}
]
be a space of possible objects and let
[
\Phi:\mathcal{X}\rightarrow\mathcal{Y}
]
be the property one wishes to determine. A verifier does not necessarily observe $x\in\mathcal{X}$ directly. It receives a representation
[
\pi:\mathcal{X}\rightarrow\mathcal{O}.
]
The representation is sufficient precisely when
[
\Phi

\widehat{\Phi}\circ\pi
]
for some
[
\widehat{\Phi}:\Img(\pi)\rightarrow\mathcal{Y}.
]
Equivalently,
[
\pi(x_1)=\pi(x_2)
\Longrightarrow
\Phi(x_1)=\Phi(x_2).
]

The verification problem is therefore naturally a problem about partitions. Every representation $\pi$ induces an equivalence relation
[
x_1\sim_\pi x_2
\quad\Longleftrightarrow\quad
\pi(x_1)=\pi(x_2).
]
Likewise, the target property induces an equivalence relation
[
x_1\sim_\Phi x_2
\quad\Longleftrightarrow\quad
\Phi(x_1)=\Phi(x_2).
]
The representation is sufficient for $\Phi$ exactly when its equivalence classes do not merge distinct $\Phi$-classes.

\begin{theorem}[Partition criterion for sufficient verification]
A representation
[
\pi:\mathcal{X}\rightarrow\mathcal{O}
]
is sufficient for
[
\Phi:\mathcal{X}\rightarrow\mathcal{Y}
]
if and only if
[
\sim_\pi
;\subseteq;
\sim_\Phi.
]
\end{theorem}

\begin{proof}
Suppose first that $\pi$ is sufficient for $\Phi$. Then there exists $\widehat{\Phi}$ such that
[
\Phi=\widehat{\Phi}\circ\pi.
]
If
[
x_1\sim_\pi x_2,
]
then
[
\pi(x_1)=\pi(x_2),
]
and therefore
[
\Phi(x_1)

\widehat{\Phi}(\pi(x_1))

\widehat{\Phi}(\pi(x_2))

\Phi(x_2).
]
Hence
[
x_1\sim_\Phi x_2.
]
Thus
[
\sim_\pi\subseteq\sim_\Phi.
]

Conversely, suppose
[
\sim_\pi\subseteq\sim_\Phi.
]
Then whenever
[
\pi(x_1)=\pi(x_2),
]
we have
[
\Phi(x_1)=\Phi(x_2).
]
By the fiber criterion, $\Phi$ factors through $\pi$. Therefore $\pi$ is sufficient.
\end{proof}

The theorem gives a direct characterization of the information required for verification. A representation may distinguish more cases than $\Phi$ requires. That is harmless with respect to sufficiency, although it may impose other costs. What it may not do is identify two cases on which $\Phi$ differs.

This immediately yields a canonical representation associated with every property. Define
[
q_\Phi:\mathcal{X}\rightarrow\mathcal{X}/{\sim_\Phi}
]
by mapping each object to its equivalence class under $\sim_\Phi$:
[
q_\Phi(x)=[x]_\Phi.
]
The quotient retains exactly the distinctions made by the property and no others.

\begin{proposition}[Canonical property quotient]
The quotient map
[
q_\Phi:\mathcal{X}\rightarrow\mathcal{X}/{\sim_\Phi}
]
is sufficient for $\Phi$. Moreover, every representation sufficient for $\Phi$ refines $q_\Phi$.
\end{proposition}

\begin{proof}
By construction,
[
q_\Phi(x_1)=q_\Phi(x_2)
]
if and only if
[
\Phi(x_1)=\Phi(x_2).
]
Therefore $\Phi$ is constant on every fiber of $q_\Phi$, and so $\Phi$ factors through $q_\Phi$.

Now let $\pi$ be any representation sufficient for $\Phi$. By the partition criterion,
[
\sim_\pi\subseteq\sim_\Phi.
]
Therefore whenever $\pi(x_1)=\pi(x_2)$, one also has
[
q_\Phi(x_1)=q_\Phi(x_2).
]
Consequently there exists a map
[
r:\Img(\pi)\rightarrow\mathcal{X}/{\sim_\Phi}
]
such that
[
q_\Phi=r\circ\pi.
]
Hence
[
\pi\succeq q_\Phi.
]
\end{proof}

This quotient is the coarsest sufficient representation in the abstract information ordering. It preserves exactly the distinction
[
\Phi(x_1)\neq\Phi(x_2)
]
and forgets everything else.

For a Boolean property
[
\Phi:\mathcal{X}\rightarrow{0,1},
]
the quotient may contain only two classes:
[
\Phi^{-1}(0)
\qquad\text{and}\qquad
\Phi^{-1}(1).
]
In this purely mathematical sense, the target property itself is a minimal sufficient statistic for determining the target property.

That observation is formally correct but operationally incomplete. If the verifier already possessed $\Phi(x)$, there would be nothing left to verify. The canonical quotient therefore describes the information boundary of the problem, not necessarily a realizable verification interface. Engineering verification requires a representation that is both sufficient and obtainable without already solving the target problem.

This motivates a distinction between \emph{informational minimality} and \emph{operational realizability}. The former asks how much information is theoretically required to determine the property. The latter asks which sufficient representation can actually be computed, observed, measured, or maintained under the constraints of the system.

Let
[
\Pi
]
be a class of realizable representations. Then the relevant search is not over all mathematically definable maps
[
\pi:\mathcal{X}\rightarrow\mathcal{O},
]
but over
[
\pi\in\Pi.
]

\begin{definition}[Realizable verification unit]
Let $\Pi$ be an admissible class of representations. A representation
[
\pi\in\Pi
]
is a realizable verification unit for $\Phi$ if $\Phi$ factors through $\pi$.
\end{definition}

A minimal realizable verification unit is then minimal only relative to $\Pi$ and to the refinement ordering. This resolves the ambiguity in saying that verification should operate on ``the smallest object'' containing the necessary information. There may be several incomparable sufficient representations, and practical constraints may exclude the mathematically coarsest one.

For example, a history-level property may be determined either by a particular graph invariant or by a compact sufficient statistic accumulated online. Neither representation need reconstruct the other, even though both determine the same target property. They can therefore be incomparable under refinement while remaining equally sufficient for $\Phi$.

Suppose
[
\pi_1:\mathcal{X}\rightarrow\mathcal{O}_1
]
and
[
\pi_2:\mathcal{X}\rightarrow\mathcal{O}_2
]
both satisfy
[
\Phi

\widehat{\Phi}_1\circ\pi_1

\widehat{\Phi}_2\circ\pi_2.
]
It need not follow that
[
\pi_1\succeq\pi_2
]
or
[
\pi_2\succeq\pi_1.
]
The existence of incomparable sufficient representations is one reason the phrase ``the minimal representation'' should be used only after an ordering and admissible representation class have been specified.

This also separates minimality from optimality. Even if a representation is minimal under informational refinement, it may not be desirable operationally. A slightly finer representation may be dramatically cheaper to compute, easier to audit, more robust to noise, or more compatible with existing system boundaries.

Let
[
C(\pi)
]
be an engineering cost functional. It may incorporate computational expense, memory, latency, privacy exposure, communication burden, proof complexity, fragility, or other system-specific considerations. The design problem becomes
[
\pi^\star
\in
\arg\min_{\pi\in\Pi}
C(\pi)
]
subject to
[
\Phi=\widehat{\Phi}\circ\pi
]
for some $\widehat{\Phi}$.

The optimal verification unit is therefore not necessarily the smallest informational unit. It is a sufficient representation that optimizes an explicitly chosen cost under architectural constraints.

This formulation also makes visible a problem that is otherwise easy to overlook: the cost of verification can depend strongly upon how the same information is represented. Two representations may induce the same partition of $\mathcal{X}$ and therefore contain the same information relative to $\Phi$, while one makes $\Phi$ computationally easy to recover and the other makes it difficult.

Suppose
[
\pi_1
]
and
[
\pi_2
]
are informationally equivalent in the sense that
[
\pi_1\succeq\pi_2
\qquad\text{and}\qquad
\pi_2\succeq\pi_1.
]
Then they preserve the same distinctions, but the corresponding decoding functions
[
\widehat{\Phi}_1
\qquad\text{and}\qquad
\widehat{\Phi}_2
]
can differ greatly in computational complexity.

Thus informational sufficiency is necessary for universal verification, but it is not sufficient for practical verification. The complete design problem concerns at least three dimensions:
[
\text{information},
\qquad
\text{computation},
\qquad
\text{cost}.
]

The first asks whether the representation preserves the target distinction. The second asks whether the property can be computed reliably from the representation. The third asks whether maintaining and evaluating the representation is operationally acceptable.

The concept of minimal support becomes particularly useful when a system must verify several properties simultaneously. Let
[
\Phi_1,\ldots,\Phi_m
]
be target properties on $\mathcal{X}$. A representation $\pi$ is jointly sufficient when each property factors through it:
[
\Phi_j

\widehat{\Phi}j\circ\pi
\qquad
\text{for every }j.
]
Equivalently,
[
\sim\pi
\subseteq
\bigcap_{j=1}^{m}
\sim_{\Phi_j}.
]

\begin{proposition}[Joint verification support]
A representation $\pi$ is sufficient for all properties
[
\Phi_1,\ldots,\Phi_m
]
if and only if
[
\sim_\pi
\subseteq
\bigcap_{j=1}^{m}\sim_{\Phi_j}.
]
\end{proposition}

\begin{proof}
If $\pi$ is sufficient for every $\Phi_j$, then
[
\sim_\pi\subseteq\sim_{\Phi_j}
]
for each $j$, and hence
[
\sim_\pi
\subseteq
\bigcap_j\sim_{\Phi_j}.
]
Conversely, if the inclusion holds, then
[
\sim_\pi\subseteq\sim_{\Phi_j}
]
for every $j$. By the partition criterion, each $\Phi_j$ factors through $\pi$.
\end{proof}

The result formalizes why verification interfaces often become richer as the number of claims they are expected to support increases. A representation sufficient for one property may merge cases that another property requires separating. The joint representation must refine the intersection of the relevant property partitions.

This also reveals a danger in generic ``safety scores.'' Suppose a system is expected to support distinct properties concerning authorization, privacy, physical safety, task correctness, reversibility, and provenance. Compressing all of these into a single scalar
[
S(x)\in\mathbb{R}
]
is sufficient only if every target property factors through $S$. That requires
[
S(x_1)=S(x_2)
\Longrightarrow
\Phi_j(x_1)=\Phi_j(x_2)
]
for every $j$.

There is no reason to assume such a scalar exists in a useful form. A one-dimensional score can contain arbitrarily much information in a set-theoretic sense if unrestricted real numbers are permitted, so dimensionality alone is not the issue. The practical issue is whether the implemented scalar representation preserves the distinctions required by the decision architecture. A scalar score that averages or aggregates independent dimensions can easily identify cases requiring different responses.

For example, let
[
D(x)

(d_{\mathrm{auth}}(x),
d_{\mathrm{risk}}(x))
]
be a two-coordinate diagnostic representation. Suppose a scalar score is
[
S(x)

d_{\mathrm{auth}}(x)
+
d_{\mathrm{risk}}(x).
]
Then cases
[
D(x_1)=(0,1)
]
and
[
D(x_2)=(1,0)
]
have identical scalar score:
[
S(x_1)=S(x_2)=1.
]
If one case requires refusal because authorization is absent while the other requires additional physical-safety inspection, the scalar representation has collapsed two cases requiring different continuations.

The failure is not that $S$ is one-dimensional. It is that the aggregation map is non-injective with respect to distinctions the response policy requires.

This suggests a general relationship between verification and control. Let
[
R:\mathcal{X}\rightarrow\mathcal{U}
]
be the correct response policy, where $\mathcal{U}$ may include permit, refuse, inspect, sense, defer, repair, or other actions. A diagnostic representation
[
D:\mathcal{X}\rightarrow\mathcal{Z}
]
is sufficient for control only if
[
R

\widehat{R}\circ D.
]
If
[
D(x_1)=D(x_2)
]
but
[
R(x_1)\neq R(x_2),
]
then no controller receiving only $D$ can implement the correct policy universally.

The relevant dimensionality of a diagnostic system is therefore determined by the distinctions required for response. A diagnostic coordinate is useful when it separates cases requiring different continuations. Additional coordinates that do not refine the response-relevant partition add no decision information, although they may serve explanatory or auditing purposes.

This point links verification to diagnosis. Verification is often framed as a binary problem:
[
\text{pass}
\quad\text{or}\quad
\text{fail}.
]
But a system that must respond differently to different failure modes requires more than a binary result. If
[
F_1,F_2,\ldots,F_k
]
are distinct fault classes requiring different interventions, then a diagnostic representation must distinguish them at least to the extent that the response function does.

A binary predicate
[
V(x)\in{0,1}
]
can establish that something failed while remaining insufficient to determine what should happen next. The failure state
[
V(x)=0
]
may contain multiple response-distinct conditions.

Suppose
[
x_1,x_2
]
both satisfy
[
V(x_1)=V(x_2)=0,
]
but
[
R(x_1)=\text{repair}
]
and
[
R(x_2)=\text{refuse}.
]
Then $R$ does not factor through $V$. A verifier that produces only pass or fail cannot implement the required response policy without additional evidence.

This distinction is particularly important when repair itself can introduce risk. Detection of failure establishes that the current object violates some criterion. It does not establish that a proposed repair is beneficial, safe, or warranted. Repair therefore introduces another transformation
[
x
\xrightarrow{\rho}
x',
]
and the desired property concerns
[
\Phi(x').
]
A repair operation requires its own preservation or improvement argument.

One may define a repair objective
[
J(\rho\mid x)
]
that incorporates expected improvement, intervention cost, collateral risk, irreversibility, and uncertainty. The existence of a detected fault
[
V(x)=0
]
does not imply
[
J(\rho\mid x)<J(\varnothing\mid x)
]
for any particular repair $\rho$. Detection and repair warrant are therefore distinct verification problems.

The same factorization criterion applies. If the choice among repairs depends upon distinctions absent from the verifier's diagnostic output, then the diagnostic output is insufficient for repair selection.

The unit of verification can thus change across stages of the same system. One representation may suffice to detect that something is wrong. A finer representation may be required to diagnose what is wrong. A still different representation may be required to determine whether intervention is warranted. Yet another may be required to verify the repaired state.

There is no contradiction in these different units because the properties differ:
[
\Phi_{\mathrm{detect}},
\qquad
\Phi_{\mathrm{diagnose}},
\qquad
\Phi_{\mathrm{warrant}},
\qquad
\Phi_{\mathrm{repair}}.
]
Each induces its own partition of the underlying state space.

This property-relative view also explains why the complete history is not automatically the ideal verification object. Let
[
\operatorname{id}{\Hist}:\Hist\rightarrow\Hist
]
be the identity representation. It is trivially sufficient for every history property:
[
\Phi=\Phi\circ\operatorname{id}{\Hist}.
]
But this mathematical sufficiency says nothing about computational tractability, privacy, storage, causal contamination, adversarial exposure, or whether the acting system should have access to the same history as the verifier.

The identity representation is therefore an upper bound on information, not a general engineering prescription.

Indeed, the temporal analysis showed that exposing additional history to an acting system can change its policy. If the same representation is used for generation and verification, increasing verifier support may simultaneously increase actor capability or alter behavior. Verification architecture must therefore distinguish
[
\pi_{\mathrm{act}},
\qquad
\pi_{\mathrm{verify}},
\qquad
\pi_{\mathrm{audit}}.
]
These projections answer different questions and may rationally preserve different distinctions.

For an action-generation property $\Phi_{\mathrm{act}}$, one may require
[
\Phi_{\mathrm{act}}

\widehat{\Phi}{\mathrm{act}}
\circ
\pi{\mathrm{act}}.
]
For an online verification property $\Phi_{\mathrm{verify}}$,
[
\Phi_{\mathrm{verify}}

\widehat{\Phi}{\mathrm{verify}}
\circ
\pi{\mathrm{verify}}.
]
For retrospective reconstruction $\Phi_{\mathrm{audit}}$,
[
\Phi_{\mathrm{audit}}

\widehat{\Phi}{\mathrm{audit}}
\circ
\pi{\mathrm{audit}}.
]
Nothing requires
[
\pi_{\mathrm{act}}

\pi_{\mathrm{verify}}

\pi_{\mathrm{audit}}.
]

This separation has an important security interpretation. Some distinctions may need to be visible to the verifier precisely because they should not be visible to the actor. A verifier may need access to hidden provenance, policy state, authorization metadata, or historical evidence used to evaluate an action. Exposing the same information to the actor could allow strategic adaptation, leakage, or circumvention.

The verifier's informational advantage can therefore be deliberate.

Conversely, the actor may possess local task information that the verifier does not require. There is no general requirement that the verifier reconstruct the actor's entire internal state. It needs sufficient support for the property it is assigned to establish.

The architecture becomes a network of property-relative interfaces rather than a hierarchy in which one omniscient observer sees everything. This is significant because ``global verification'' can otherwise suggest that the solution to local insufficiency is simply universal observability. Universal observability is neither necessary nor always desirable. The objective is sufficient observability.

This can be stated as an information-separation principle:
[
\boxed{
\text{The information required to generate an action, verify an action, and audit an action need not be identical.}
}
]

The principle follows directly from the property-relative account. Generation, verification, and auditing are different functions and can induce different sufficient partitions.

The same reasoning applies across administrative or authority boundaries. Suppose a distributed system contains agents
[
1,\ldots,n,
]
with local observations
[
\pi_i(H).
]
A global property $\Phi(H)$ may fail to factor through any individual $\pi_i$:
[
\forall i,\qquad
\Phi\neq\widehat{\Phi}_i\circ\pi_i.
]
It may nevertheless factor through their joint observation:
[
\Phi

\widehat{\Phi}
\circ
(\pi_1,\ldots,\pi_n).
]
In that case the property is globally decidable from distributed evidence but not locally decidable by any single observer.

This gives a precise notion of distributed verification support.

\begin{definition}[Distributed sufficient support]
A family of observations
[
{\pi_i:\mathcal{X}\rightarrow\mathcal{O}i}{i=1}^{n}
]
is jointly sufficient for $\Phi$ if there exists
[
\widehat{\Phi}:
\Img(\pi_1,\ldots,\pi_n)
\rightarrow\mathcal{Y}
]
such that
[
\Phi(x)

\widehat{\Phi}
\bigl(
\pi_1(x),\ldots,\pi_n(x)
\bigr).
]
\end{definition}

Joint sufficiency does not imply that all observations should be centralized. Cryptographic protocols, secure multiparty computation, distributed proofs, attestations, or other mechanisms can sometimes establish a function of distributed evidence without exposing every underlying observation to one party. The factorization criterion specifies what information must be jointly represented for the property, not how that information must be physically centralized.

This distinction is important because a naive response to distributed insufficiency would be to create an omniscient verifier. Such a verifier can create new concentration risks, privacy problems, attack surfaces, and authority conflicts. A better question is whether the required property can be computed from distributed attestations or selectively disclosed evidence while preserving the relevant distinctions.

The unit of verification is therefore not necessarily a physical object. It can be an informational composition distributed across actors, time, representations, and evidence channels.

The projection-insufficiency result becomes particularly useful here. Suppose every local observer sees
[
\pi_i(H_1)=\pi_i(H_2)
]
for two histories $H_1,H_2$, while
[
\Phi(H_1)\neq\Phi(H_2).
]
Then even the complete tuple of local observations is identical:
[
(\pi_1(H_1),\ldots,\pi_n(H_1))

(\pi_1(H_2),\ldots,\pi_n(H_2)).
]
No aggregation protocol operating only on those observations can determine $\Phi$ universally. Additional evidence or a refined interface is required.

On the other hand, it can happen that every observer is individually insufficient while the tuple is sufficient:
[
\pi_i(H_1)=\pi_i(H_2)
]
for some relevant pairs at each local site, but the combined map
[
\Pi(H)

(\pi_1(H),\ldots,\pi_n(H))
]
separates every pair on which $\Phi$ differs. The property then exists at a distributed informational scale. No single local checkpoint contains it, but their composition does.

This is a more precise alternative to describing such properties as simply ``global.'' Globality can mean several things. A property may require long temporal support, relational structural support, multiple administrative observations, counterfactual evidence, or some combination. The factorization framework identifies the relevant notion by asking which distinctions must survive.

The phrase ``unit of verification'' should therefore not be interpreted as requiring one indivisible object. A unit is whatever representation constitutes sufficient support for the property. It can be a local state, a path, a graph, a history statistic, a tuple of distributed attestations, a formal representation plus a translation certificate, or a predictive model augmented with uncertainty and recoverability information.

This leads to a hierarchy of questions that should precede verifier selection. Given a target property $\Phi$, one should ask which pairs of cases differ on $\Phi$, which observational interfaces distinguish those pairs, which of those interfaces are realizable, which are safe to expose to which components, and which permit tractable inference. Only after these questions have been answered does it make sense to ask which verifier architecture performs best.

In this order, verification begins with the property rather than with the model.

The distinction is consequential for benchmarking. A benchmark normally presents a fixed input representation and a target label. It thereby assumes that the representation is an appropriate unit for the target. Projection insufficiency shows that this assumption is testable. If the benchmark representation collapses target-distinct cases, then part of the measured error is not a failure of reasoning. It is a failure of benchmark support.

Likewise, an apparent ceiling in verifier performance can arise from three different causes:
[
\text{model limitation},
\qquad
\text{data limitation},
\qquad
\text{representation limitation}.
]
Only the first is directly addressed by increasing model capability. More training data can address some forms of the second. Neither solves the third unless the additional data refine the actual representation supplied at inference time.

This gives a practical interpretation of projection insufficiency. When a verifier repeatedly fails, one should not ask only
[
\text{How can the verifier reason better?}
]
One should also ask
[
\text{What distinction would have made these cases separable?}
]
If the missing distinction is absent from the interface, the repair belongs upstream.

The central constructive result of this paper can now be stated compactly.

\begin{theorem}[Property-relative unit of verification]
Let
[
\Phi:\mathcal{X}\rightarrow\mathcal{Y}
]
be a target property and
[
\pi:\mathcal{X}\rightarrow\mathcal{O}
]
an observational representation. A universally correct verifier for $\Phi$ whose complete input is $\pi(x)$ can exist only if
[
\Phi

\widehat{\Phi}\circ\pi
]
for some function $\widehat{\Phi}$. When this condition holds, $\pi$ is informationally sufficient for $\Phi$; when it fails, no increase in downstream verifier capability can make $\pi$ universally sufficient.
\end{theorem}

\begin{proof}
Necessity follows from projection insufficiency. If no such $\widehat{\Phi}$ exists, then by the fiber criterion there exist
[
x_1,x_2
]
such that
[
\pi(x_1)=\pi(x_2)
]
while
[
\Phi(x_1)\neq\Phi(x_2).
]
Any verifier receiving only $\pi(x)$ must receive identical input on $x_1$ and $x_2$ and therefore cannot return both distinct correct values.

If the factorization does hold, then $\widehat{\Phi}$ itself defines a universally correct verifier at the level of mathematical existence:
[
\widehat{\Phi}(\pi(x))

\Phi(x).
]
This establishes informational sufficiency, although it does not imply that $\widehat{\Phi}$ is computationally feasible, learnable, robust, or otherwise operationally realizable.
\end{proof}

The final qualification is essential. The theorem characterizes the information boundary of verification, not the complete engineering problem. It tells us when universal verification is impossible because the representation has erased necessary distinctions. It does not tell us that every sufficient representation admits an efficient verifier.

This separation is useful precisely because it prevents different limitations from being confused. Computational difficulty can motivate better algorithms. Statistical difficulty can motivate better data or calibration. Projection insufficiency requires changing the information supplied to the verifier or weakening the property being claimed.

The four case studies can now be placed within this theorem. An isolated action can be an insufficient projection of a history when coordination depends upon earlier interaction. A bag of events can be an insufficient projection of a workflow when the target depends upon topology. A translated formal object can be an insufficient projection of a source when the translation collapses property-relevant semantics. A predictive representation can be an insufficient projection of a decision situation when intervention depends upon uncertainty, consequence severity, counterfactuals, or recoverability not represented by the predictor.

In every case, the relevant question is the same:
[
\text{Does the claimed property factor through what the verifier can see?}
]

If the answer is no, the verifier is being asked to distinguish what its interface has made indistinguishable. If the answer is yes, the remaining questions concern inference, computation, calibration, and preservation. This provides a clean division of labor among problems that are otherwise often grouped together under the broad heading of verification.

The appropriate unit of verification is therefore neither inherently local nor inherently global. It is neither necessarily an output nor necessarily a trajectory. It is not defined by the amount of information retained, the physical size of the object, or the sophistication of the verifier. It is defined relative to the property.

A verification unit is adequate when it preserves the distinctions upon which that property depends. A minimal unit is one that preserves no unnecessary distinctions relative to an explicit information ordering and admissible representation class. An optimal unit is a sufficient one that also satisfies the practical constraints of the system.

This reframes the engineering problem. Instead of asking how far the checkpoint should be expanded, one asks what the checkpoint must be able to distinguish. Instead of assuming that more context is always safer, one asks which information must persist, which information must be exposed, and to whom. Instead of treating verifier intelligence as the primary bottleneck, one asks whether the verification interface makes the target property identifiable in the first place.

The resulting principle is simple:
[
\boxed{
\text{A verifier cannot recover a distinction that its representation has already erased.}
}
]
Everything else follows from deciding which distinctions the claimed property actually requires.

\section{Beyond the Checkpoint}

The checkpoint is an attractive architectural primitive because it converts an open-ended process into a bounded question. An object arrives at a boundary, a predicate is evaluated, and the object either proceeds or does not. The appeal is both computational and institutional. A checkpoint localizes responsibility, limits the amount of information that must be inspected at once, creates a place at which records can be produced, and transforms a continuous process into a sequence of auditable decisions.

Nothing in the preceding argument makes checkpoints obsolete. The conclusion is instead that a checkpoint must be interpreted according to the property it actually establishes.

Adelman's \emph{Beyond the Checkpoint}~\cite{adelman2014} provides a useful methodological counterpoint because the checkpoint in her analysis is not merely a physical location. It participates in a wider organization of mobility, surveillance, categorization, anticipation, and authority. The present argument concerns computational verification rather than the political and cultural analysis developed in that work, and the analogy should not be extended beyond what it supports. Its methodological usefulness lies in a narrower observation: concentrating attention on the visible site of inspection can obscure the larger system through which the significance of inspection is produced.

The same risk appears in computational verification. An output filter is visible. A formal proof is visible. A workflow anomaly score is visible. A safety classifier is visible. A pass or refusal event is visible. These become natural places to locate the concept of verification because they are discrete moments at which a judgment is emitted.

Yet the preceding sections show that the epistemic basis of that judgment can extend in both directions away from the checkpoint. Upstream transformations determine what information reaches it. Downstream composition determines what conclusions can legitimately be transported away from it.

The checkpoint can therefore be represented as an interface:
[
X
\xrightarrow{\pi}
O
\xrightarrow{V}
C
\xrightarrow{\kappa}
U,
]
where $X$ is the underlying object or process, $\pi$ is the representation supplied to the verifier, $V$ is the verification procedure, $C$ is the resulting certificate or diagnostic state, and $\kappa$ maps that result into a continuation or intervention $U$.

Each arrow carries a distinct obligation.

The first asks whether
[
\pi(X)
]
preserves the distinctions required by the property. The second asks whether
[
V
]
correctly infers the intended property from those observations. The third asks whether the certificate supports the continuation selected from it.

A visible verification event therefore sits inside a larger verification path:
[
\text{world}
\longrightarrow
\text{representation}
\longrightarrow
\text{judgment}
\longrightarrow
\text{continuation}.
]
Errors can arise at every transition.

Suppose the intended property is
[
\Phi:X\rightarrow{0,1}.
]
If
[
\Phi=\widehat{\Phi}\circ\pi,
]
then the representation is sufficient. If the verifier correctly computes $\widehat{\Phi}$, then
[
V(\pi(x))=\Phi(x).
]
But suppose the continuation policy is
[
\kappa:{0,1}\rightarrow\mathcal{U}.
]
Even a perfectly correct verification result does not establish that $\kappa$ is the correct response policy.

For example,
[
\Phi(x)=0
]
may establish that the current object fails a specification. It does not by itself determine whether the appropriate response is
[
\text{refuse},
\qquad
\text{repair},
\qquad
\text{retry},
\qquad
\text{request evidence},
\qquad
\text{defer},
]
or some other continuation.

Thus
[
\text{correct detection}
\quad\not\Rightarrow\quad
\text{correct intervention}.
]

This gives the checkpoint three logically distinct questions:
[
\begin{aligned}
&\text{What can be observed?}\
&\text{What can be concluded?}\
&\text{What does that conclusion warrant?}
\end{aligned}
]
Collapsing these questions into a single pass/fail boundary conceals the composition laws connecting them.

A more explicit architecture separates observation, diagnosis, and disposition. Let
[
\pi:X\rightarrow O
]
produce evidence,
[
D:O\rightarrow Z
]
produce a diagnostic state, and
[
R:Z\times O\rightarrow\mathcal{U}
]
produce a response. The architecture becomes
[
x
\xrightarrow{\pi}
o
\xrightarrow{D}
z
\xrightarrow{R(\cdot,o)}
u.
]
The response function can retain access to the underlying evidence rather than receiving only a compressed pass/fail result.

This is useful whenever the diagnostic state is sufficient for detection but insufficient for disposition. If two objects satisfy
[
D(\pi(x_1))

D(\pi(x_2))
]
while the warranted responses differ,
[
R^\star(x_1)\neq R^\star(x_2),
]
then $D$ is not a sufficient interface for response selection. Either the diagnostic representation must be refined or the response stage must retain access to additional evidence.

The familiar binary checkpoint
[
V(x)\in{\mathrm{pass},\mathrm{fail}}
]
is therefore a special case rather than the canonical form of verification. It is appropriate when the target property and the continuation policy genuinely factor through that binary distinction. Where they do not, binary verification compresses too early.

This can be understood as a problem of premature quotienting. A binary verifier partitions the object space into
[
V^{-1}(0)
\qquad\text{and}\qquad
V^{-1}(1).
]
Every distinction within each class is erased from the verifier output. If downstream behavior must distinguish two members of the same class, then the binary certificate is not a sufficient continuation interface.

The checkpoint may still be useful for one purpose while being insufficient for another. A binary result may suffice to halt obviously invalid execution while a richer diagnostic record is retained for repair, audit, or appeal. The mistake lies not in compression itself but in assuming that one compressed output should serve every downstream property.

This suggests an architectural separation among at least four functions:
[
\text{detection},
\qquad
\text{diagnosis},
\qquad
\text{disposition},
\qquad
\text{audit}.
]
Detection asks whether a specified condition holds. Diagnosis asks which distinctions explain the condition. Disposition asks what should happen next. Audit asks what evidence should remain available for later reconstruction or challenge.

The representations sufficient for these functions can differ:
[
\pi_{\mathrm{det}},
\qquad
\pi_{\mathrm{diag}},
\qquad
\pi_{\mathrm{disp}},
\qquad
\pi_{\mathrm{audit}}.
]
A well-designed system need not force them to coincide.

This separation is particularly important for persistent histories. Let the complete event history at time $t$ be
[
H_{\leq t}.
]
A system can retain a record
[
R_t(H_{\leq t}),
]
expose a restricted projection to the actor
[
M_t(H_{\leq t}),
]
supply another projection to the online verifier
[
E_t(H_{\leq t}),
]
and preserve a richer projection for audit
[
A_t(H_{\leq t}).
]
These functions may satisfy
[
M_t\neq E_t\neq A_t\neq R_t.
]

Retention therefore should not be equated with availability.

An event may be physically or logically retained without being exposed to an acting component. It may be exposed to a verifier without being admissible evidence for a particular decision. It may be admissible evidence without being sufficient to warrant an intervention. These are distinct relations.

For an event $e$, one can distinguish predicates
[
\operatorname{Retained}(e),
]
[
\operatorname{Visible}i(e),
]
[
\operatorname{Admissible}{\Phi}(e),
]
and
[
\operatorname{Warranting}{u}(e).
]
No general equivalence connects them:
[
\operatorname{Retained}(e)
\not\Rightarrow
\operatorname{Visible}i(e),
]
[
\operatorname{Visible}i(e)
\not\Rightarrow
\operatorname{Admissible}{\Phi}(e),
]
and
[
\operatorname{Admissible}{\Phi}(e)
\not\Rightarrow
\operatorname{Warranting}{u}(e).
]

The distinctions matter because ``the system knows'' is otherwise dangerously ambiguous. Information may exist somewhere in the architecture while being unavailable to the component whose behavior depends upon it. Conversely, a component may have access to information that should not be used for a particular decision.

Verification requires not merely information possession but a declared evidential relation.

Let
[
\mathcal{E}{\Phi}(x)
]
denote the evidence admissible for determining property $\Phi$ about object $x$. A verification procedure should ideally be understood as
[
V\Phi:
\mathcal{E}{\Phi}(x)
\rightarrow
\mathcal{C}{\Phi},
]
where $\mathcal{C}_{\Phi}$ is a space of property-relative certificates or diagnostic results.

This makes explicit that evidence can be property-relative. The same observation may be relevant to one claim and irrelevant, inadmissible, or misleading for another.

A provenance record provides a simple example. Provenance can be decisive when the property concerns authorization:
[
\Phi_{\mathrm{auth}}(x).
]
The same provenance may be irrelevant to a purely mathematical property:
[
\Phi_{\mathrm{math}}(x).
]
Conversely, a numerical proof may establish $\Phi_{\mathrm{math}}$ while saying nothing about whether the computation was authorized to occur.

Combining all available evidence indiscriminately is therefore not equivalent to better verification. Evidence must be interpreted relative to the proposition being established.

The distinction also clarifies the role of audit logs. A log can preserve evidence required to reconstruct an execution without granting that evidence causal influence over the execution itself. This is an important architectural possibility:
[
\text{retain more than the actor can condition upon}.
]
Such asymmetry allows retrospective verification to operate over a richer representation than action generation.

Suppose
[
M(H_1)=M(H_2)
]
for the actor-facing projection, while
[
A(H_1)\neq A(H_2)
]
for the audit projection. The actor cannot condition its current action on the distinction between $H_1$ and $H_2$ through $M$, while the auditor can later recover that distinction through $A$.

This permits historical accountability without necessarily creating the same historical dependence in the acting policy.

The distinction is especially relevant to the long-horizon case. Simply giving an agent its entire retained history can increase the set of strategies it can condition upon. Retaining the history externally for verification does not have the same consequence. Memory architecture is therefore part of verification architecture.

The checkpoint should also be separated from the authority to modify the object it evaluates. A verifier can possess epistemic authority concerning a predicate without possessing operational authority to intervene.

Let
[
V(x)=z
]
be a diagnostic result and let
[
\mathcal{U}(x,z)
]
be the set of interventions technically reachable from the checkpoint. Let
[
\mathcal{W}(x,z)
\subseteq
\mathcal{U}(x,z)
]
be the subset for which intervention is warranted.

Then
[
u\in\mathcal{U}(x,z)
]
does not imply
[
u\in\mathcal{W}(x,z).
]
Reachability is not warrant.

This distinction matters in automated repair. A verifier may detect a violation and have the technical capability to rewrite the object immediately. Yet the fact that repair is reachable does not establish that repair is the correct continuation. The system may lack sufficient diagnostic resolution to know which repair preserves other properties.

Suppose a current object satisfies a vector of properties
[
\mathbf{\Phi}(x)

(\Phi_1(x),\ldots,\Phi_m(x)).
]
A repair transformation
[
\rho:x\mapsto x'
]
may improve one coordinate:
[
\Phi_j(x)=0,
\qquad
\Phi_j(x')=1,
]
while degrading another:
[
\Phi_k(x)=1,
\qquad
\Phi_k(x')=0.
]
A verifier that observes only $\Phi_j$ cannot determine this tradeoff.

Thus successful repair of the detected property does not imply successful repair of the object:
[
\Phi_j(\rho(x))=1
\quad\not\Rightarrow\quad
\mathbf{\Phi}(\rho(x))
\succeq
\mathbf{\Phi}(x).
]

A repair architecture consequently needs a declared repair objective. Let
[
J(x,\rho(x))
]
measure the relevant change across properties, costs, and risks. The intervention problem becomes
[
\rho^\star
\in
\arg\min_{\rho\in\mathcal{R}}
J(x,\rho(x)),
]
possibly including the identity transformation
[
\rho_0(x)=x.
]
Including the identity is the repair analogue of including the null action. It permits the system to represent the conclusion that a detected problem does not yet warrant modification.

This is a useful safeguard against what might otherwise be called compulsory repair. If the response space contains only modifications, the system must alter the object once a fault is detected, even when every available repair carries greater expected cost than leaving the object unchanged pending further evidence.

The same principle applies to refusal. Refusal should be represented as one possible disposition rather than as the semantic meaning of every failed verification. A failed check can establish
[
\neg\Phi(x)
]
without establishing
[
\operatorname{Refuse}(x).
]
To obtain the latter, one requires a rule
[
\neg\Phi(x)
\Longrightarrow
\operatorname{Refuse}(x),
]
or a more elaborate decision procedure using additional evidence.

In some systems that implication is intentionally part of the specification. For example, failure of an authentication predicate may be defined to require refusal. In such a case the composition law is explicit. The problem arises when the same inference is imported into domains where the relation has not been specified.

The distinction can be summarized by separating a predicate from its policy:
[
\Phi(x)
\qquad\text{versus}\qquad
\kappa(\Phi(x),E).
]
The predicate describes the object. The policy determines a continuation given that description and perhaps additional evidence $E$.

This separation also supports appeal and revision. If a verifier stores only the disposition
[
u=\text{refuse},
]
the evidential basis for the decision may be lost. If it instead stores a certificate
[
c=(\Phi,z,E,\mathcal{A}),
]
where $z$ is the diagnostic result and $\mathcal{A}$ records relevant assumptions, a later process can reconsider the disposition without pretending that the original observation never occurred.

A verification record should therefore distinguish
[
\text{observation},
\quad
\text{inference},
\quad
\text{assumption},
\quad
\text{disposition}.
]
This separation makes revision possible when one element changes.

Suppose a verifier at time $t$ produces
[
c_t=(o_t,\phi_t,\mathcal{A}_t,u_t).
]
New evidence at time $t+1$ may leave
[
o_t
]
unchanged while invalidating an assumption in
[
\mathcal{A}_t.
]
The earlier observation remains historically true, but the disposition it supported may no longer be warranted.

This is another reason append-only histories can be preferable to destructive state replacement for audit. Revision need not mean erasure. A later record can state that a previous conclusion has been superseded while preserving the evidence and reasoning that produced it.

Let
[
L=(e_1,e_2,\ldots,e_n)
]
be an append-only verification ledger. A revision can be represented by a new event
[
e_{n+1}

\operatorname{Revise}(e_k,c'),
]
rather than deleting $e_k$. The current state can be computed from the ledger while the historical sequence remains reconstructible.

This separates
[
\text{current validity}
]
from
[
\text{historical occurrence}.
]
A certificate can cease to govern future continuation without ceasing to have existed.

Persistent histories therefore support a richer concept of verification than a mutable flag. A flag answers
[
\text{What is the current verdict?}
]
A ledger can also answer
[
\text{What was observed?}
]
[
\text{What was concluded?}
]
[
\text{Under which assumptions?}
]
[
\text{What changed?}
]
and
[
\text{Why is the current verdict different?}
]

This matters whenever verification itself is subject to verification.

A system that can reconstruct its own decision path permits second-order questions such as whether the correct evidence was available, whether the correct verifier was invoked, whether the preservation assumptions held, and whether the disposition remained within its authorized scope.

One can represent such a certificate as
[
C=
\langle
\Phi,
\pi,
o,
V,
z,
\mathcal{A},
\kappa,
u
\rangle,
]
where $\Phi$ is the property claimed, $\pi$ the observation interface, $o=\pi(x)$ the observation, $V$ the verifier, $z$ its result, $\mathcal{A}$ the assumptions under which the result is interpreted, $\kappa$ the continuation rule, and $u$ the resulting disposition.

This certificate does not prove itself correct. It makes the structure of the claim inspectable.

The benefit is that scope becomes explicit. A bare statement
[
\text{verified}
]
is underspecified. A structured certificate states, in effect:
[
\text{property }\Phi
\text{ was evaluated from representation }\pi
\text{ under assumptions }\mathcal{A}.
]
The verifier's conclusion can then be compared with the claim someone later attempts to derive from it.

This suggests a general discipline of \emph{scope-carrying verification}. A verification result should carry enough metadata to identify the proposition established and the conditions under which it was established.

In its simplest mathematical form, instead of returning
[
V(o)\in{0,1},
]
the verifier returns
[
V(o)

(\Phi,z,\mathcal{A}),
]
where $z$ is the result. Downstream systems are then expected to use the certificate only through rules whose premises match
[
\Phi
]
and
[
\mathcal{A}.
]

This resembles type information. A value cannot be safely substituted merely because some computation produced it. Its type constrains the contexts in which it can be used. Similarly, a verification result should not be freely substituted for arbitrary stronger claims. Its scope constrains the inferences in which it can participate.

One can make the analogy formal. Let
[
C_\Phi
]
denote a certificate establishing property $\Phi$. A downstream rule has the form
[
C_{\Phi_1},\ldots,C_{\Phi_k},\mathcal{A}
\vdash
C_\Psi.
]
The rule itself is a composition theorem. It states that certificates for the premises, under assumptions $\mathcal{A}$, are sufficient to establish $\Psi$.

Under such an architecture, global assurance is not produced by declaring a single omniscient checkpoint. It is constructed through explicit certificate transformations.

This provides a positive interpretation of compositional verification:
[
C_{\Phi_1}
+
\cdots
+
C_{\Phi_k}
+
\text{preservation rule}
\Longrightarrow
C_\Psi.
]
The preservation rule is not administrative glue. It is part of the proof.

A checkpoint architecture becomes robust when the boundaries between certificates are treated as formally as the certificates themselves.

This is the central sense in which verification must move beyond the checkpoint. The phrase does not mean eliminating checkpoints or replacing them with continuous total surveillance. It means refusing to treat the visible moment of inspection as the complete epistemic object.

The relevant system includes what happened before the checkpoint, what representation crossed it, what distinctions that representation preserved, what proposition the verifier actually established, under what assumptions the proposition holds, and what downstream transformation converts that proposition into action.

The four case studies now appear as four ways this extended system can matter. Long-horizon interaction shows that the current checkpoint may omit causal history. Workflow analysis shows that independent checkpoints may omit relational structure. Neuro-formal verification shows that a checkpoint may operate on a transformed object whose relation to the source remains part of the proof obligation. Safety-critical world models show that even an accurate predictive checkpoint may not contain the distinctions required for intervention.

The design response is therefore not a universal demand for larger checkpoints. It is a demand for explicit interfaces among verification units.

For each claimed property $\Phi$, the system should identify a representation
[
\pi_\Phi
]
through which $\Phi$ factors. For each transformation
[
T
]
across which a certificate is transported, the system should identify the preservation relation required. For each disposition
[
u,
]
the system should identify the evidential and normative rule connecting the verified property to that disposition.

The resulting architecture can be summarized as
[
\boxed{
\text{evidence}
\rightarrow
\text{property}
\rightarrow
\text{certificate}
\rightarrow
\text{warrant}
\rightarrow
\text{continuation}.
}
]
No arrow should be treated as free.

This perspective changes what it means to ask whether a system is verified. The question is incomplete unless one asks:
[
\text{Verified with respect to what property?}
]
[
\text{From which representation?}
]
[
\text{Under which assumptions?}
]
[
\text{Across which transformations?}
]
[
\text{For which subsequent inference or intervention?}
]

These are not qualifications added after verification. They specify the verification claim itself.

The checkpoint remains useful precisely when its limits are explicit. A local test can provide strong assurance about a local property. A graph verifier can provide strong assurance about a structural property. A theorem prover can provide strong assurance about a formal representation. A world model can provide strong evidence about future states. None of these results becomes weaker when its scope is stated accurately.

What becomes weaker is only the unsupported inference beyond that scope.

The objective is therefore not to make verification less decisive. It is to make the boundary of decision coincide with the boundary of evidence. Where the property extends farther, the verification unit must extend or compose accordingly. Where the evidence does not support intervention, the architecture must be capable of retaining, inspecting, deferring, or refusing to conclude.

Beyond the checkpoint lies not an absence of verification but its full structure.

\section{Conclusion}

Verification is often imagined as the place where uncertainty ends. A system produces an output, a verifier inspects it, and a verdict converts an uncertain process into an accepted or rejected result. This architecture is useful precisely because it creates boundaries. It permits complex processes to be decomposed, intermediate objects to be certified, responsibilities to be separated, and errors to be intercepted before they propagate.

The difficulty begins when the scope of the conclusion silently exceeds the scope of the verification.

The four technical cases examined here expose different forms of this expansion. Long-horizon agent interaction shows that properties of a trajectory need not be properties of any isolated action. Workflow-level anomaly detection shows that properties of relational organization need not be determined by properties of individual events. Neuro-formal verification shows that a theorem about a transformed formal object need not automatically be a theorem about its source. Safety-critical world models show that an accurate prediction need not contain sufficient information to determine whether intervention is warranted.

These are not instances of one empirical phenomenon. They share a logical structure:
[
\text{property established at }X
\quad\Longrightarrow?\quad
\text{property asserted at }Y.
]
The arrow is where the proof obligation resides.

The central mathematical result is correspondingly simple. Let
[
\pi:\mathcal{X}\rightarrow\mathcal{O}
]
be the representation available to a verifier and
[
\Phi:\mathcal{X}\rightarrow\mathcal{Y}
]
the property it is expected to determine. Universal verification from $\pi$ is possible only if
[
\Phi=\widehat{\Phi}\circ\pi
]
for some
[
\widehat{\Phi}:\Img(\pi)\rightarrow\mathcal{Y}.
]
Equivalently,
[
\pi(x_1)=\pi(x_2)
\Longrightarrow
\Phi(x_1)=\Phi(x_2).
]
If two property-distinct cases become identical at the verification interface, no downstream increase in model intelligence, computation, debate, redundancy, or formal sophistication can reconstruct the erased distinction from that interface alone.

This establishes an informational boundary on verification.

It also establishes why the appropriate verification unit cannot be specified independently of the property. An action can be sufficient for one property and insufficient for another. A workflow graph can be necessary for a relational property and excessive for an event-local property. A formal translation can be faithful for memory safety while omitting distinctions relevant to timing or provenance. A predictive representation can be sufficient for forecasting while insufficient for intervention.

The question is therefore not
[
\text{Should verification be local or global?}
]
but
[
\text{Which distinctions must survive for this property to be decidable?}
]

This formulation avoids the opposite mistake of treating maximal observation as the solution to local insufficiency. The complete history is trivially richer than an isolated event, but retaining, exposing, and processing the complete history can introduce computational, privacy, strategic, and causal costs. More information can alter an actor's behavior even when it improves an auditor's reconstruction. The actor, verifier, and auditor therefore need not receive identical projections of the same history.

The design objective is sufficient support, not maximal support.

The argument also separates three questions that are frequently compressed into one:
[
\boxed{
\text{support}
\quad+\quad
\text{correct inference}
\quad+\quad
\text{preservation}.
}
]
Support asks whether the verifier receives the distinctions required by the property. Correct inference asks whether it computes the property correctly from those distinctions. Preservation asks whether the established result survives the transformation across which a broader conclusion is drawn.

A failure of support cannot be repaired merely by improving inference. A failure of inference does not show that the representation was insufficient. A failure of preservation does not show that the original verifier was unsound. Keeping these failure modes separate is necessary for locating what must actually change.

This distinction becomes especially important when verification results enter institutional decision systems. The logical structure developed here is not limited to software. A quantitative indicator can be correctly calculated while the policy conclusion attached to it requires premises that the indicator itself does not establish.

Clara E. Mattei's \emph{Escape from Capitalism} provides a useful political-economic example of this distinction. Mattei challenges presentations of austerity, unemployment, inflation management, balanced budgets, and related economic categories as if the associated policy conclusions followed neutrally from economic measurement itself. Her broader argument is political and historical: economic categories and expert institutions, in her account, operate within a social order and can help present particular distributions of constraint as economic necessity rather than political choice. The present paper need not adopt that larger thesis in order to extract a narrower verification problem.

Suppose an institution observes a fiscal quantity
[
d_t

\frac{\text{deficit}_t}{\text{output}_t}
]
and verifies correctly that
[
d_t>\tau.
]
That establishes a proposition about the measured fiscal state. It does not, without an additional rule, establish
[
\operatorname{CutSpending}.
]
The complete inference has the form
[
d_t>\tau
]
together with
[
d_t>\tau
\Longrightarrow
\operatorname{CutSpending}.
]
The first proposition can be empirical. The second is a policy rule.

No improvement in the accuracy with which $d_t$ is measured proves the second proposition. A more precise checkpoint does not convert a decision rule into an empirical fact.

The same distinction applies to unemployment. Let
[
u_t
]
denote a measured unemployment rate. Suppose a model establishes a relation between unemployment, wage pressure, and some target inflation measure. Even if the empirical relation were estimated perfectly within a specified population and period, the inference
[
u_t<u^\star
\Longrightarrow
\operatorname{IncreaseUnemployment}
]
would contain additional premises concerning which target should govern, which costs matter, how those costs are distributed, which alternative interventions are admissible, and what authority the institution possesses to impose them.

The distinction can be represented by separating a descriptive state
[
z\in\mathcal{Z}
]
from an institutional response
[
a\in\mathcal{A}.
]
An empirical procedure estimates
[
z=\pi(x).
]
A policy rule then maps
[
\kappa:\mathcal{Z}\times\Gamma\rightarrow\mathcal{A},
]
where
[
\Gamma
]
contains governing objectives, constraints, institutional authority, distributional commitments, and other premises relevant to the decision.

The actual decision is therefore
[
a=\kappa(z,\Gamma),
]
not simply
[
a=z.
]

Two institutions can agree completely about the measured state
[
z
]
while disagreeing about
[
\Gamma
]
and consequently select different interventions. Such disagreement does not necessarily imply that one side has mismeasured the economy. It may concern the rule connecting measurement to action.

This gives a useful institutional analogue of the epistemic composition problem:
[
\text{economic indicator}
\quad\not\Rightarrow\quad
\text{policy necessity}.
]

Mattei's discussion of austerity makes this distinction especially visible. In her account, fiscal retrenchment is not simply equivalent to ``spending less.'' It involves choices about which expenditures are reduced, which interests remain protected, and who bears the resulting adjustment. A scalar fiscal aggregate can therefore be insufficient to characterize the distributional structure of the intervention.

Let total expenditure be
[
E

\sum_{i=1}^{n}E_i,
]
where
[
E_i
]
is expenditure directed toward category $i$. A target requiring
[
E'<E
]
does not determine the vector
[
(E_1',E_2',\ldots,E_n').
]
There may be many allocations satisfying
[
\sum_i E_i'=E'
]
with radically different effects.

Thus
[
\Delta E<0
]
does not uniquely determine
[
\Delta E_i.
]

A verification procedure concerned only with aggregate expenditure can establish that a fiscal target has been met while remaining informationally incapable of determining how the adjustment was distributed.

Suppose two budgets $B_1$ and $B_2$ satisfy
[
E(B_1)=E(B_2)
]
and
[
d(B_1)=d(B_2),
]
but differ in their distributional consequences:
[
D(B_1)\neq D(B_2).
]
Then a representation
[
\pi(B)=(E(B),d(B))
]
is insufficient for $D$.

By the fiber criterion,
[
\pi(B_1)=\pi(B_2)
]
together with
[
D(B_1)\neq D(B_2)
]
proves that no function
[
\widehat{D}
]
can satisfy
[
D=\widehat{D}\circ\pi
]
universally.

The result does not determine which distributional policy should be chosen. It establishes something prior to that political choice: aggregate fiscal indicators alone cannot determine a property that varies while those indicators remain fixed.

This distinction provides a precise way to discuss a recurrent issue in political economy without pretending that mathematics resolves the political dispute. If two policies occupy the same equivalence class under the reported indicator but produce different outcomes under a property citizens care about, then the indicator is not sufficient support for choosing between them.

The relevant remedy is not to assign the indicator a political meaning in the opposite direction. It is to expose the missing dimensions.

For a policy $p$, one might therefore distinguish
[
F(p)
]
for its fiscal effect,
[
I(p)
]
for its inflation effect,
[
U(p)
]
for its employment effect,
[
W(p)
]
for its distributional effect,
[
R(p)
]
for its reversibility, and
[
A(p)
]
for whatever authorization or institutional criterion governs its adoption.

A policy decision
[
K(p)
]
that depends upon all six cannot factor through $F(p)$ alone:
[
K
\neq
\widehat{K}\circ F
]
whenever there exist policies with identical fiscal effects but different values on decision-relevant coordinates.

Nor does replacing $F$ with a weighted score automatically solve the problem. If
[
S(p)

\alpha F(p)
+\beta I(p)
+\gamma U(p)
+\delta W(p),
]
then $S$ is sufficient for the final decision only if policies assigned the same score never require different decisions under the governing criterion. The aggregation weights themselves encode a valuation rule. They are not supplied by the measured quantities.

The mathematical form is identical to the anomaly-score problem encountered earlier. A scalar can be diagnostically useful without being a complete warrant for intervention.

Mattei's discussion of economic expertise adds another layer. Her argument contests the presentation of institutions and economic doctrines as politically neutral when their recommendations embody substantive commitments. In the language developed here, this can be reframed without deciding the larger ideological question: the transition
[
\text{measurement}
\rightarrow
\text{model}
\rightarrow
\text{objective}
\rightarrow
\text{policy}
]
contains multiple transformations, and the empirical validity of an upstream stage does not make the downstream transformations assumption-free.

Consider
[
x
\xrightarrow{\pi}
z
\xrightarrow{M}
\widehat{y}
\xrightarrow{J}
r
\xrightarrow{\kappa}
a,
]
where $x$ is the economic state, $z$ a measured representation, $M$ a predictive model, $\widehat{y}$ a predicted consequence, $J$ a valuation or objective functional, $r$ a ranking, and $\kappa$ a policy rule.

A dispute can occur at any arrow.

One party may dispute the measurement:
[
\pi.
]
Another may accept the measurement but dispute the model:
[
M.
]
Another may accept both but dispute the objective:
[
J.
]
Another may accept the objective while disputing whether the institution has authority to implement
[
\kappa.
]

Calling the complete sequence ``what the data require'' collapses these distinct propositions.

The distinction is not peculiar to economics. Medicine separates test results from treatment decisions. Engineering separates sensor measurements from control policies. Courts distinguish evidence from judgment and judgment from remedy. Machine learning distinguishes prediction from intervention. In each case, the quality of the upstream evidence matters greatly, but evidence acquires operational force through additional rules.

Mattei's examples are useful precisely because economic discourse can make those rules difficult to see. Terms such as "necessity,'' "discipline,'' "stability,'' and "soundness'' can appear to describe measured states while also functioning as conclusions about what must be done. The verification framework asks that these roles be separated.

If
[
\Phi_{\mathrm{desc}}(x)
]
is descriptive and
[
\Phi_{\mathrm{act}}(x)
]
is action-guiding, then the implication
[
\Phi_{\mathrm{desc}}(x)
\Longrightarrow
\Phi_{\mathrm{act}}(x)
]
is a substantive bridge. It should be inspectable as such.

This does not imply that the bridge is illegitimate. Some decision rules are well justified. Some institutional constraints are explicit and democratically established. Some interventions may be strongly supported by empirical evidence. The formal point is narrower: the justification resides partly in the bridge, not solely in the measurement preceding it.

This distinction also prevents the framework from being used as a disguised argument for policy paralysis. The conclusion is not
[
\text{measurement never warrants action}.
]
It is
[
\text{measurement warrants action through an explicit decision relation}.
]
When the relation is well specified and its premises are satisfied, action can be strongly justified. When the relation is disputed, increasing the precision of the measurement does not by itself resolve the dispute.

The same principle applies to nonintervention. Choosing
[
\varnothing
]
is itself a policy when the existing state has continuing effects. The null action should therefore be evaluated rather than treated as normatively empty. An architecture that makes intervention prove itself while allowing the status quo to proceed without evaluation embeds an asymmetry into the decision rule.

Let
[
J(a\mid E)
]
be the criterion applied to interventions. Consistency may require evaluating
[
J(\varnothing\mid E)
]
under the same relevant dimensions. Otherwise the checkpoint compares proposed change against an unscored baseline.

This point connects back to Mattei without requiring agreement with her political conclusions. One of the implications of treating an economic order as natural or inevitable is that continuation can disappear from the space of explicit decisions. The existing arrangement becomes the baseline against which alternatives must justify themselves.

In verification terms, however, continuation is still a transition:
[
x_t
\xrightarrow{\varnothing}
x_{t+1}.
]
If the environment evolves, nonintervention can have consequences:
[
x_{t+1}\neq x_t.
]
The absence of an active control operation is not necessarily the absence of a trajectory.

This yields a more general symmetry principle:
[
\boxed{
\text{When intervention and continuation both have consequences, both belong to the decision problem.}
}
]

The principle is relevant far beyond political economy. A software system can continue running rather than being repaired. An autonomous vehicle can continue rather than maneuver. A medical system can withhold treatment. A moderator can leave content unchanged. A financial authority can leave policy settings unchanged. In each case the null action has a trajectory and can therefore be compared with alternatives.

The checkpoint should not conceal this comparison.

The broader lesson is that verification systems are embedded in continuation systems. A verifier does not merely label the present. Its output changes which futures remain reachable.

Let
[
\mathcal{R}(x)
]
denote the set of continuations reachable from state $x$. A checkpoint result $c$ may induce a restricted set
[
\mathcal{R}_c(x)
\subseteq
\mathcal{R}(x).
]
A refusal removes some paths. A repair creates others. A request for evidence delays commitment. An acceptance permits a continuation that might otherwise have been blocked.

Verification therefore participates in the topology of future possibility.

This makes the distinction between diagnosis and warrant especially important. A false inference at the checkpoint does not merely produce an incorrect label. It can remove admissible futures or admit inadmissible ones.

One can represent the consequence of a verification architecture by
[
\Gamma_V(x)

\mathcal{R}_{V(\pi(x))}(x).
]
Two verifiers with identical classification accuracy can nevertheless have different consequences if their errors remove different future options. Verification quality cannot therefore always be reduced to a scalar accuracy measure.

A false refusal and a false acceptance may have different costs. A reversible false intervention and an irreversible false intervention may have different costs. An error that preserves the possibility of later correction differs from one that destroys the evidence required to discover the error.

The relevant objective may consequently depend upon the continuation structure:
[
J(V)

\mathbb{E}
\left[
L\bigl(
x,
V(\pi(x)),
\Gamma_V(x)
\bigr)
\right].
]
This evaluates verification partly by what its judgments do to reachable futures.

The framework developed here therefore ends at a somewhat different place from where it began. The initial problem appeared to concern the scale of verification. Should a system inspect outputs, trajectories, workflows, source programs, or world models?

The answer is that scale is secondary.

The primary object is the property.

Once the property is specified, one can ask which distinctions determine it. Those distinctions determine the informational support required. The architecture must then preserve that support through whatever transformations intervene between the underlying process and the verifier. The verifier must correctly infer the property from that representation. Any subsequent intervention requires an additional relation connecting the verified property to the action.

Thus:
[
\boxed{
\text{property}
\rightarrow
\text{sufficient support}
\rightarrow
\text{verification}
\rightarrow
\text{preservation}
\rightarrow
\text{warranted continuation}.
}
]

The four technical cases show what happens when one of these relations is assumed rather than established. Mattei's political-economic examples show why the same discipline matters when measurements enter institutions: the fact that a quantity is real, measurable, or correctly estimated does not make the policy mapping attached to it part of the measurement itself. That mapping is another object of analysis.

The point is not that every verification system must become indefinitely expansive. Quite the opposite. Once the target property is specified precisely, irrelevant information can be discarded deliberately. A good verification interface is not one that sees everything. It is one that preserves what the claimed conclusion requires and does not pretend that omitted distinctions remain available.

This is the difference between a checkpoint and a verification boundary.

A checkpoint is a place where inspection occurs.

A verification boundary specifies how far the resulting claim travels.

The distinction matters because modern computational systems increasingly combine learned representations, formal tools, long-lived memory, multi-agent interaction, autonomous intervention, and institutional decision rules. In such systems there may be no single location at which the relevant property exists in complete form. It may depend upon a history, a relation among agents, a source-to-model correspondence, distributed evidence, or a counterfactual comparison among futures.

The verification architecture must follow that dependence.

A locally sound test remains valuable. A formally proved transformed program remains valuable. A highly accurate world model remains valuable. A carefully measured economic indicator remains valuable. Their value does not depend upon allowing them to imply more than they establish.

Rigor is not increased by enlarging the rhetoric surrounding a result. It is increased by preserving the boundary between what has been shown and what still requires a bridge.

The central conclusion is therefore not a rejection of checkpoints but a rule for using them:
[
\boxed{
\text{No conclusion should cross a boundary that its supporting property has not been shown to cross.}
}
]

Where composition is proved, local verification can support global assurance. Where representation is faithful, formal proof can support source-level claims. Where predictive information is decision-sufficient, prediction can support intervention. Where an indicator is connected to policy through explicit governing premises, measurement can support action.

Where those bridges are absent, the checkpoint marks the end of the verified claim, not the beginning of an automatic necessity.

The unit of verification is the smallest realizable informational structure through which the property of interest can be recovered. The unit of decision may be larger, because decision can require properties beyond those established by the verifier. The unit of audit may be larger still, because later reconstruction can require evidence deliberately withheld from both actor and online decision process.

Recognizing these differences replaces a single universal checkpoint with a compositional architecture of scoped claims.

That architecture does not eliminate uncertainty. It makes uncertainty locatable. It does not eliminate judgment. It makes the premises of judgment visible. It does not prevent action. It distinguishes action supported by the available evidence from action smuggled into the meaning of the evidence itself.

Verification, in this account, is not the moment at which a system declares that a question has been settled. It is the disciplined determination of which question has actually been settled, from which information, under which assumptions, and how far the answer is entitled to travel.


\newpage 
\section*{Appendices}

\appendix

\section{Factorization, Quotients, and Minimal Verification Support}
\label{app:factorization}

This appendix develops the mathematical structure underlying the notion of a verification unit. The central object is not a verifier but a triple
[
(\mathcal X,\pi,\Phi),
]
where $\mathcal X$ is a state or history space,
[
\pi:\mathcal X\to\mathcal O
]
is an observational representation, and
[
\Phi:\mathcal X\to\mathcal Y
]
is the target property. The principal question is whether $\Phi$ is determined by $\pi$.

The deterministic theory reduces to the factorization problem
[
\Phi=\widehat{\Phi}\circ\pi.
]
This elementary equation generates a useful structure of equivalence relations, quotients, refinement orderings, minimal sufficient representations, joint verification units, and composition laws.

\subsection{Observation-Induced Equivalence Relations}

Every representation
[
\pi:\mathcal X\to\mathcal O
]
induces an equivalence relation
[
x\sim_\pi x'
\quad\Longleftrightarrow\quad
\pi(x)=\pi(x').
]
Likewise, every property
[
\Phi:\mathcal X\to\mathcal Y
]
induces
[
x\sim_\Phi x'
\quad\Longleftrightarrow\quad
\Phi(x)=\Phi(x').
]

The equivalence classes
[
[x]_\pi

{x'\in\mathcal X:x'\sim_\pi x}
]
are the fibers of $\pi$. They contain precisely the source objects that are observationally indistinguishable to any verifier restricted to $\pi$.

The equivalence classes
[
[x]_\Phi

{x'\in\mathcal X:\Phi(x')=\Phi(x)}
]
contain the objects that need not be distinguished for the sole purpose of evaluating $\Phi$.

The fundamental relation is therefore
[
\sim_\pi\subseteq\sim_\Phi.
]

\begin{theorem}[Factorization theorem]
Let
[
\pi:\mathcal X\to\mathcal O
]
and
[
\Phi:\mathcal X\to\mathcal Y.
]
The following are equivalent:

\begin{enumerate}
\item There exists
[
\widehat{\Phi}:\Img(\pi)\to\mathcal Y
]
such that
[
\Phi=\widehat{\Phi}\circ\pi.
]

\item For every $x,x'\in\mathcal X$,

[
\pi(x)=\pi(x')
\Longrightarrow
\Phi(x)=\Phi(x').
]

\item

[
\sim_\pi\subseteq\sim_\Phi.
]

\item $\Phi$ is constant on every fiber

[
\pi^{-1}(o),
\qquad
o\in\Img(\pi).
]
\end{enumerate}
\end{theorem}

\begin{proof}
The equivalence of (2), (3), and (4) follows immediately from the definitions.

Suppose (1) holds. If
[
\pi(x)=\pi(x'),
]
then
[
\Phi(x)

\widehat{\Phi}(\pi(x))

\widehat{\Phi}(\pi(x'))

\Phi(x'),
]
so (2) holds.

Conversely, suppose (2) holds. Define
[
\widehat{\Phi}(o)=\Phi(x)
]
for any $x$ satisfying
[
\pi(x)=o.
]
Such an $x$ exists because
[
o\in\Img(\pi).
]
Condition (2) guarantees that the value is independent of the representative chosen. Thus $\widehat{\Phi}$ is well defined and
[
\widehat{\Phi}(\pi(x))

\Phi(x)
]
for every $x\in\mathcal X$. Hence
[
\Phi=\widehat{\Phi}\circ\pi.
]
\end{proof}

This theorem gives a complete deterministic criterion for informational sufficiency.

\subsection{Indistinguishability and Impossibility}

The negative result follows immediately.

\begin{theorem}[Deterministic projection impossibility]
Suppose there exist
[
x_0,x_1\in\mathcal X
]
such that
[
\pi(x_0)=\pi(x_1)
]
but
[
\Phi(x_0)\neq\Phi(x_1).
]
Then there exists no function
[
f:\Img(\pi)\to\mathcal Y
]
satisfying
[
f(\pi(x))=\Phi(x)
]
for every $x\in\mathcal X$.
\end{theorem}

\begin{proof}
Suppose such an $f$ existed. Then
[
f(\pi(x_0))

\Phi(x_0)
]
and
[
f(\pi(x_1))

\Phi(x_1).
]
But
[
\pi(x_0)=\pi(x_1),
]
so
[
f(\pi(x_0))

f(\pi(x_1)).
]
Hence
[
\Phi(x_0)=\Phi(x_1),
]
contradicting the hypothesis.
\end{proof}

The theorem is independent of computational complexity. It remains true if $f$ is allowed unbounded time, unbounded memory, arbitrary search, arbitrary model capacity, or an oracle for every proposition other than the missing distinction itself.

The obstruction is informational:
[
\pi(x_0)=\pi(x_1).
]

\subsection{Canonical Quotient Representation}

Define
[
q_\Phi:\mathcal X\to\mathcal X/{\sim_\Phi}
]
by
[
q_\Phi(x)=[x]_\Phi.
]

Then $\Phi$ factors canonically through the quotient:
[
\mathcal X
\xrightarrow{q_\Phi}
\mathcal X/{\sim_\Phi}
\xrightarrow{\widetilde{\Phi}}
\mathcal Y,
]
where
[
\widetilde{\Phi}([x]_\Phi)=\Phi(x).
]

The map is well defined because all representatives of $[x]_\Phi$ have the same $\Phi$-value.

\begin{theorem}[Universal property of the property quotient]
Let $\pi:\mathcal X\to\mathcal O$ be sufficient for $\Phi$. Then there exists a unique map
[
r:\Img(\pi)\to\mathcal X/{\sim_\Phi}
]
such that
[
q_\Phi=r\circ\pi.
]
\end{theorem}

\begin{proof}
Define
[
r(\pi(x))=[x]\Phi.
]
To show that $r$ is well defined, suppose
[
\pi(x)=\pi(x').
]
Since $\pi$ is sufficient for $\Phi$,
[
\Phi(x)=\Phi(x'),
]
and therefore
[
[x]\Phi=[x']_\Phi.
]
Hence the definition does not depend upon the representative.

For every $x$,
[
(r\circ\pi)(x)

r(\pi(x))

[x]_\Phi

q_\Phi(x).
]
Thus
[
q_\Phi=r\circ\pi.
]

Uniqueness follows because every element of $\Img(\pi)$ has the form $\pi(x)$, so any map satisfying the factorization must obey
[
r(\pi(x))=q_\Phi(x).
]
\end{proof}

The quotient
[
\mathcal X/{\sim_\Phi}
]
is therefore canonical up to isomorphism as the coarsest abstract representation retaining exactly the distinctions required by $\Phi$.

\subsection{The Refinement Preorder}

Let $\Pi(\mathcal X)$ denote a collection of representations of $\mathcal X$. Define
[
\pi_1\succeq\pi_2
]
when there exists a map $r$ satisfying
[
\pi_2=r\circ\pi_1.
]

Thus $\pi_1$ is at least as informative as $\pi_2$.

\begin{proposition}
The relation $\succeq$ is a preorder.
\end{proposition}

\begin{proof}
Reflexivity follows from
[
\pi=\operatorname{id}\circ\pi.
]

For transitivity, suppose
[
\pi_1\succeq\pi_2
]
and
[
\pi_2\succeq\pi_3.
]
Then there exist $r_{12}$ and $r_{23}$ such that
[
\pi_2=r_{12}\circ\pi_1
]
and
[
\pi_3=r_{23}\circ\pi_2.
]
Therefore
[
\pi_3

r_{23}\circ r_{12}\circ\pi_1,
]
so
[
\pi_1\succeq\pi_3.
]
\end{proof}

Antisymmetry need not hold as literal equality. Two representations may encode exactly the same information in different codomains. Define
[
\pi_1\equiv\pi_2
]
when
[
\pi_1\succeq\pi_2
]
and
[
\pi_2\succeq\pi_1.
]
The quotient of representations by $\equiv$ is partially ordered by informational refinement.

\begin{theorem}[Upward closure of sufficiency]
If
[
\pi_1\succeq\pi_2
]
and $\pi_2$ is sufficient for $\Phi$, then $\pi_1$ is sufficient for $\Phi$.
\end{theorem}

\begin{proof}
Since $\pi_2$ is sufficient,
[
\Phi=\widehat{\Phi}_2\circ\pi_2.
]
Since
[
\pi_2=r\circ\pi_1,
]
we obtain
[
\Phi

\widehat{\Phi}_2\circ r\circ\pi_1.
]
Defining
[
\widehat{\Phi}_1=\widehat{\Phi}_2\circ r
]
gives
[
\Phi=\widehat{\Phi}_1\circ\pi_1.
]
\end{proof}

Hence the sufficient representations for a fixed property form an upward-closed subset of the information preorder.

Let
[
\operatorname{Suff}(\Phi)

{\pi:\Phi\text{ factors through }\pi}.
]
Then
[
\pi\in\operatorname{Suff}(\Phi),
\qquad
\pi'\succeq\pi
]
imply
[
\pi'\in\operatorname{Suff}(\Phi).
]

\subsection{Minimal Sufficient Representations}

Relative to a representation class $\Pi$, define
[
\operatorname{Suff}_\Pi(\Phi)

\operatorname{Suff}(\Phi)\cap\Pi.
]

\begin{definition}[Minimal sufficient representation]
A representation
[
\pi^\star\in\operatorname{Suff}\Pi(\Phi)
]
is minimal if there is no
[
\pi\in\operatorname{Suff}\Pi(\Phi)
]
such that
[
\pi^\star\succeq\pi
]
and
[
\pi\not\equiv\pi^\star.
]
\end{definition}

Minimal elements need not be unique.

Indeed, $\Pi$ can contain incomparable representations
[
\pi_1,\pi_2
]
such that both are sufficient but neither satisfies
[
\pi_1\succeq\pi_2
]
nor
[
\pi_2\succeq\pi_1.
]

This is one reason a verification unit cannot generally be described as \emph{the} minimal representation without specifying the admissible representation family and its ordering.

If arbitrary representations are allowed, however, the property quotient is minimal.

\begin{theorem}[Abstract minimality of the property quotient]
For every sufficient representation $\pi$,
[
\pi\succeq q_\Phi.
]
Hence $q_\Phi$ is a least element of $\operatorname{Suff}(\Phi)$ up to informational equivalence.
\end{theorem}

This follows immediately from the universal property proved above.

The theorem is informational rather than computational. Constructing
[
q_\Phi(x)
]
may require knowing $\Phi(x)$ already. Thus the quotient identifies the minimal information that must be retained, not necessarily a feasible method for obtaining that information.

\subsection{Joint Properties}

Let
[
\Phi_i:\mathcal X\to\mathcal Y_i,
\qquad
i=1,\ldots,m.
]
Define the joint property
[
\mathbf{\Phi}(x)

\bigl(
\Phi_1(x),
\ldots,
\Phi_m(x)
\bigr).
]

Then
[
x\sim_{\mathbf{\Phi}}x'
]
if and only if
[
\Phi_i(x)=\Phi_i(x')
]
for every $i$. Consequently,
[
\sim_{\mathbf{\Phi}}

\bigcap_{i=1}^{m}\sim_{\Phi_i}.
]

\begin{theorem}[Joint sufficiency]
A representation $\pi$ is sufficient for every $\Phi_i$ if and only if it is sufficient for $\mathbf{\Phi}$.
\end{theorem}

\begin{proof}
If $\pi$ is sufficient for every $\Phi_i$, there exist maps $\widehat{\Phi}_i$ such that
[
\Phi_i=\widehat{\Phi}_i\circ\pi.
]
Define
[
\widehat{\mathbf{\Phi}}(o)

\bigl(
\widehat{\Phi}_1(o),
\ldots,
\widehat{\Phi}_m(o)
\bigr).
]
Then
[
\mathbf{\Phi}

\widehat{\mathbf{\Phi}}\circ\pi.
]

Conversely, if
[
\mathbf{\Phi}

\widehat{\mathbf{\Phi}}\circ\pi,
]
then composing with the $i$th coordinate projection
[
p_i:
\prod_j\mathcal Y_j
\to
\mathcal Y_i
]
gives
[
\Phi_i

p_i\circ\widehat{\mathbf{\Phi}}\circ\pi.
]
Thus each $\Phi_i$ factors through $\pi$.
\end{proof}

This establishes formally that adding verification objectives can only refine, never coarsen, the canonical sufficient partition:
[
\sim_{\mathbf{\Phi}}

\bigcap_i\sim_{\Phi_i}
\subseteq
\sim_{\Phi_j}
]
for every $j$.

\subsection{Composition of Representations}

Consider
[
\mathcal X
\xrightarrow{\pi}
\mathcal O
\xrightarrow{\rho}
\mathcal Z.
]
The composite representation is
[
\sigma=\rho\circ\pi.
]

Since
[
\pi(x)=\pi(x')
\Longrightarrow
\sigma(x)=\sigma(x'),
]
we have
[
\pi\succeq\sigma.
]

Post-processing cannot increase information in this deterministic ordering.

\begin{corollary}[No deterministic recovery after destructive compression]
If $\Phi$ does not factor through $\sigma=\rho\circ\pi$, then no function applied only after $\sigma$ can universally recover $\Phi$.
\end{corollary}

More generally, let
[
f:\mathcal Z\to\mathcal W.
]
Then
[
f\circ\sigma
]
is still a function of $\sigma$. If
[
\sigma(x_0)=\sigma(x_1)
]
for a property-distinct pair, every subsequent deterministic computation receives identical input on the pair.

Thus
[
\text{computation after information loss}
]
cannot reverse
[
\text{information loss}.
]

This is the deterministic analogue of a data-processing principle.

\subsection{Local Predicates and Global Properties}

Let
[
H=(x_1,\ldots,x_n)\in\mathcal X^n
]
be a finite history. Suppose each component has a local predicate
[
V_i:\mathcal X\to{0,1}.
]
Define
[
L(H)

\bigwedge_{i=1}^{n}V_i(x_i).
]

The local-verification projection is
[
\pi_L(H)

\bigl(
V_1(x_1),
\ldots,
V_n(x_n)
\bigr).
]

A global property
[
\Phi:\mathcal X^n\to{0,1}
]
is decidable from the local verification results if and only if
[
\Phi

\widehat{\Phi}\circ\pi_L.
]

In particular, the implication
[
L(H)=1
\Longrightarrow
\Phi(H)=1
]
is weaker than full factorization. It constrains only the fiber
[
\pi_L^{-1}(1,\ldots,1).
]

\begin{definition}[Positive-fiber sufficiency]
The local predicates are positively sufficient for $\Phi$ if
[
\pi_L(H)=(1,\ldots,1)
\Longrightarrow
\Phi(H)=1.
]
\end{definition}

Thus a conventional compositional verification theorem need not make $\Phi$ fully recoverable from all local verdict patterns. It need only establish that the all-pass fiber lies inside the globally valid set:
[
\pi_L^{-1}(1,\ldots,1)
\subseteq
\Phi^{-1}(1).
]

This yields a geometric interpretation of compositional verification.

Let
[
A_i

{H:V_i(x_i)=1}.
]
Then
[
L^{-1}(1)

\bigcap_{i=1}^{n}A_i.
]
The desired composition theorem is
[
\bigcap_{i=1}^{n}A_i
\subseteq
\Phi^{-1}(1).
]

Failure of composition means
[
\left(
\bigcap_{i=1}^{n}A_i
\right)
\cap
\Phi^{-1}(0)
\neq
\varnothing.
]

Any element of this intersection is a constructive counterexample:
[
H^\star
\in
\bigcap_i A_i
\cap
\Phi^{-1}(0).
]

It satisfies every local verifier while violating the global property.

\subsection{Certificates as Logical Objects}

Suppose verifier $i$ produces a certificate
[
C_i
]
corresponding to proposition
[
P_i.
]
A global certificate $C_\Phi$ is derivable only if the proof system contains a rule
[
P_1,\ldots,P_n,\mathcal A
\vdash
\Phi,
]
where $\mathcal A$ denotes auxiliary assumptions.

The existence of certificates
[
C_1,\ldots,C_n
]
does not itself establish the sequent.

Formally, let
[
\operatorname{Th}(S)
]
denote the deductive closure of a proposition set $S$. Then
[
\Phi
]
is licensed exactly when
[
\Phi
\in
\operatorname{Th}
\left(
{P_1,\ldots,P_n}\cup\mathcal A
\right).
]

Hence
[
\forall i,;P_i
]
and
[
\Phi
]
remain logically distinct until a derivation connects them.

This provides a proof-theoretic counterpart to the factorization framework. Factorization asks whether the available information distinguishes the relevant cases. Deduction asks whether the available propositions entail the desired proposition.

Both are necessary in a complete verification architecture.

\subsection{Transformation Chains}

Consider spaces
[
\mathcal X_0,\ldots,\mathcal X_n
]
and transformations
[
T_i:\mathcal X_{i-1}\to\mathcal X_i.
]
Let
[
\Phi_i:\mathcal X_i\to{0,1}.
]

Suppose each transformation satisfies the reflection condition
[
\Phi_i(T_i(x))=1
\Longrightarrow
\Phi_{i-1}(x)=1.
]

\begin{theorem}[Backward transport through a transformation chain]
If
[
\Phi_n
\left(
T_n\circ\cdots\circ T_1(x)
\right)
=1,
]
then
[
\Phi_0(x)=1.
]
\end{theorem}

\begin{proof}
Let
[
x_0=x
]
and recursively define
[
x_i=T_i(x_{i-1}).
]
By hypothesis,
[
\Phi_n(x_n)=1.
]
Reflection for $T_n$ yields
[
\Phi_{n-1}(x_{n-1})=1.
]
Repeated application gives
[
\Phi_{n-2}(x_{n-2})=1,
\ldots,
\Phi_0(x_0)=1.
]
\end{proof}

The theorem formalizes source-level certification through a sequence of transformations. Every link requires the correct implication direction.

If reflection fails at one index $k$, the chain theorem cannot be invoked, even if every other transformation is perfectly verified.

\subsection{Preservation Defect}

The binary distinction between preservation and non-preservation can be generalized by defining a defect set.

For
[
T:\mathcal X\to\mathcal Y,
]
source property $P_X$, and target property $P_Y$, define the reflection-defect set
[
\Delta_T

\left{
x\in\mathcal X:
P_Y(T(x))=1
\land
P_X(x)=0
\right}.
]

Exact reflection is equivalent to
[
\Delta_T=\varnothing.
]

Given a probability measure $\mu$ on $\mathcal X$, define the defect mass
[
\delta_\mu(T)

\mu(\Delta_T).
]

Then
[
\delta_\mu(T)=0
]
is distribution-relative almost-sure reflection, whereas
[
\Delta_T=\varnothing
]
is universal reflection.

These should not be conflated.

A transformation can satisfy
[
\delta_\mu(T)=0
]
under one deployment distribution while possessing source objects outside the support of $\mu$ on which reflection fails.

This distinction will become important in the probabilistic appendix.

\subsection{Verification Interfaces as Partitions}

The partition perspective allows an entire verification architecture to be described without reference to the internal implementation of its verifier.

Let
[
\mathcal P_\pi

{
\pi^{-1}(o):o\in\Img(\pi)
}
]
be the partition induced by $\pi$.

Let
[
\mathcal P_\Phi

{
\Phi^{-1}(y):y\in\Img(\Phi)
}.
]

Then
[
\pi\text{ is sufficient for }\Phi
]
if and only if
[
\mathcal P_\pi
]
refines
[
\mathcal P_\Phi.
]

Using the conventional partition-refinement notation
[
\mathcal P_\pi\preceq\mathcal P_\Phi,
]
where the left partition is finer, sufficiency becomes
[
\boxed{
\mathcal P_\pi\preceq\mathcal P_\Phi.
}
]

The verification-unit problem can therefore be posed as a search over partitions:
[
\text{find }\mathcal P
]
such that
[
\mathcal P\preceq\mathcal P_\Phi
]
while minimizing some cost
[
C(\mathcal P).
]

If $\mathfrak P$ is the admissible family of realizable partitions, the optimization problem is
[
\mathcal P^\star
\in
\arg\min_{\mathcal P\in\mathfrak P}
C(\mathcal P)
]
subject to
[
\mathcal P\preceq\mathcal P_\Phi.
]

This formulation separates the semantic constraint from the engineering objective.

The semantic constraint says:
[
\text{do not merge property-distinct cases}.
]

The engineering objective says:
[
\text{among representations satisfying that constraint, choose an acceptable one}.
]

\subsection{Monotonicity Under Property Refinement}

Suppose two properties
[
\Phi:\mathcal X\to\mathcal Y
]
and
[
\Psi:\mathcal X\to\mathcal Z
]
satisfy
[
\Phi=f\circ\Psi
]
for some $f$.

Then $\Psi$ is at least as discriminating as $\Phi$:
[
\sim_\Psi\subseteq\sim_\Phi.
]

\begin{proposition}
If $\pi$ is sufficient for $\Psi$, then $\pi$ is sufficient for $\Phi$.
\end{proposition}

\begin{proof}
If
[
\Psi=\widehat{\Psi}\circ\pi
]
and
[
\Phi=f\circ\Psi,
]
then
[
\Phi

f\circ\widehat{\Psi}\circ\pi.
]
Thus $\Phi$ factors through $\pi$.
\end{proof}

The converse does not generally hold. A representation sufficient for a coarse property can be insufficient for a refinement of that property.

For example, a binary property
[
\Phi(x)=
\mathbf{1}{\text{failure}}
]
may factor through a pass/fail diagnostic while a finer property
[
\Psi(x)
\in
{
\text{none},
\text{authorization},
\text{semantic},
\text{structural},
\text{epistemic}
}
]
does not.

Thus diagnostic refinement can require representational refinement.

\subsection{Decision Sufficiency as a Special Case}

Let
[
R:\mathcal X\to\mathcal U
]
be a response policy. A representation $\pi$ is decision-sufficient exactly when
[
R=\widehat{R}\circ\pi.
]

All preceding results therefore apply with
[
\Phi=R.
]

In particular,
[
\pi(x_1)=\pi(x_2)
]
and
[
R(x_1)\neq R(x_2)
]
prove that $\pi$ is insufficient for decision.

Suppose a diagnostic property
[
D:\mathcal X\to\mathcal Z
]
is sufficient for detecting a condition but the response depends upon additional variables:
[
R(x)=r(D(x),W(x)).
]
If $W$ does not factor through $D$, then
[
R
]
need not factor through $D$.

Hence
[
\text{diagnostic sufficiency}
\not\Rightarrow
\text{decision sufficiency}.
]

This is the formal structure underlying the distinction among detection, diagnosis, warrant, and disposition developed in the main text.

\subsection{A Lattice-Theoretic View}

When $\mathcal X$ is finite, the set of all partitions of $\mathcal X$ forms a lattice under refinement. Denote it
[
\operatorname{Part}(\mathcal X).
]

For partitions $\mathcal P$ and $\mathcal Q$, their meet
[
\mathcal P\wedge\mathcal Q
]
is the common refinement generated by intersections of blocks:
[
\mathcal P\wedge\mathcal Q

{
P\cap Q:
P\in\mathcal P,;
Q\in\mathcal Q,;
P\cap Q\neq\varnothing
}.
]

If $\Phi_1$ and $\Phi_2$ are properties, then
[
\mathcal P_{(\Phi_1,\Phi_2)}

\mathcal P_{\Phi_1}
\wedge
\mathcal P_{\Phi_2}.
]

Thus adding verification objectives corresponds to moving downward toward finer partitions in the refinement lattice.

If a representation $\pi$ must support properties
[
\Phi_1,\ldots,\Phi_m,
]
its induced partition must satisfy
[
\mathcal P_\pi
\preceq
\bigwedge_{i=1}^{m}
\mathcal P_{\Phi_i}.
]

The meet therefore characterizes the canonical joint verification requirement.

Conversely, removing a property can permit coarsening. If an expensive distinction matters only to $\Phi_k$, dropping $\Phi_k$ from the verification contract may permit the interface to identify states previously required to remain distinct.

This gives verification contracts an algebraic interpretation: changing the set of promised properties changes the minimum refinement required of the observation interface.

\subsection{Verification Debt from Unproved Edges}

Consider a directed architecture graph
[
\mathcal G=(V,E).
]
Associate each node
[
v\in V
]
with a property
[
P_v
]
and each directed edge
[
e=(u,v)
]
with a transformation
[
T_e.
]

Define
[
\chi(e)

\begin{cases}
1,
&
P_u(x)\Rightarrow P_v(T_e(x))
\text{ has been established},\
0,
&
\text{otherwise}.
\end{cases}
]

For a path
[
p=(e_1,\ldots,e_k),
]
define its preservation certificate
[
\chi(p)

\prod_{j=1}^{k}\chi(e_j).
]

Then
[
\chi(p)=1
]
if and only if every edge on the path possesses the required preservation result.

A verified node at the beginning of the path licenses the target property at the end only when
[
\chi(p)=1.
]

The number
[
D(p)

\sum_{j=1}^{k}(1-\chi(e_j))
]
counts unproved preservation edges along the path. It can be interpreted as a simple structural measure of verification debt.

This quantity is not a probability of failure. It measures missing proof obligations.

One unproved edge is sufficient to prevent the path-level implication from being established:
[
D(p)>0
\Longrightarrow
\text{path-level preservation not yet proved}.
]

Thus a pipeline with many locally verified nodes can remain globally uncertified because its verification graph contains uncertified edges.

This formalizes the distinction between verification density and verification scope:
[
\text{node coverage}
\neq
\text{edge preservation}.
]

\subsection{Summary}

The deterministic theory can be condensed into five statements.

First, a representation $\pi$ supports a property $\Phi$ exactly when
[
\Phi=\widehat{\Phi}\circ\pi.
]

Second, this is equivalent to
[
\sim_\pi\subseteq\sim_\Phi.
]

Third, the quotient
[
q_\Phi:\mathcal X\to\mathcal X/{\sim_\Phi}
]
is the canonical least informative representation sufficient for $\Phi$, considered abstractly.

Fourth, increasing verifier capability cannot overcome a violation of
[
\sim_\pi\subseteq\sim_\Phi.
]

Fifth, transporting a verified property across a transformation requires a separate preservation relation for that transformation.

The resulting architecture has two mathematically distinct kinds of edges:
[
\boxed{
\text{observation edges}
\qquad\text{and}\qquad
\text{preservation edges}.
}
]
Observation edges determine which distinctions reach a verifier. Preservation edges determine which conclusions survive movement between objects, representations, scales, or stages.

A complete verification argument requires both.

\section{Minimal Countermodels to Naive Composition}
\label{app:countermodels}

The claims in the main text are primarily impossibility claims of the form
[
A\not\Rightarrow B.
]
Such claims are established by counterexample. This appendix constructs minimal or near-minimal countermodels for the four principal failures of composition and then derives a general counterexample-construction theorem.

The purpose is not to model the empirical systems examined in the main text in full detail. It is to isolate the logical structure responsible for each failure. If a proposed implication fails in a two-state or three-node system, increasing the complexity of the surrounding architecture cannot make that implication valid without introducing additional assumptions.

\subsection{Temporal Countermodel}

Let the action alphabet be
[
\mathcal A={0,1}.
]
Consider histories of length two:
[
\Hist_2=\mathcal A^2.
]
Define the local verifier
[
V:\mathcal A\to{0,1}
]
by
[
V(0)=V(1)=1.
]
Thus every individual action is locally admissible.

Define a history property
[
\Phi_{\mathrm{eq}}(a_1,a_2)

\mathbf 1{a_1=a_2}.
]
Then
[
\Phi_{\mathrm{eq}}(0,0)=1,
\qquad
\Phi_{\mathrm{eq}}(0,1)=0.
]

Nevertheless,
[
V(0)=V(0)=1
]
for the first history and
[
V(0)=V(1)=1
]
for the second.

Hence
[
V(a_1)\land V(a_2)=1
]
does not determine
[
\Phi_{\mathrm{eq}}(a_1,a_2).
]

More explicitly,
[
(0,0)
\in
V^{-1}(1)\times V^{-1}(1)
]
and
[
(0,1)
\in
V^{-1}(1)\times V^{-1}(1),
]
but
[
\Phi_{\mathrm{eq}}(0,0)
\neq
\Phi_{\mathrm{eq}}(0,1).
]

Thus
[
\bigwedge_{t=1}^{2}V(a_t)=1
\not\Rightarrow
\Phi_{\mathrm{eq}}(H)=1.
]

The example requires only two time steps. A one-step history cannot exhibit a genuinely relational temporal property between distinct actions. In this sense, history length two is minimal for this form of temporal non-composition.

A second example shows that the problem remains even when the local verifier is nontrivial.

Let
[
\mathcal A={0,1,2}
]
and define
[
V(a)=
\mathbf 1{a\neq2}.
]
The locally admissible set is therefore
[
{0,1}.
]
Define
[
\Phi_{\mathrm{alt}}(a_1,a_2)

\mathbf 1{a_1\neq a_2}.
]
Then
[
H_1=(0,1)
]
and
[
H_2=(0,0)
]
both satisfy
[
V(a_1)=V(a_2)=1,
]
while
[
\Phi_{\mathrm{alt}}(H_1)=1
]
and
[
\Phi_{\mathrm{alt}}(H_2)=0.
]

The failure therefore does not depend upon a vacuous verifier. It follows whenever the locally accepted region contains multiple actions whose relations matter to the history-level property.

\begin{proposition}[Criterion for temporal non-composition]
Let
[
A_{\mathrm{acc}}

{a\in\mathcal A:V(a)=1}.
]
If there exist histories
[
H,H'\in A_{\mathrm{acc}}^n
]
such that
[
\Phi(H)\neq\Phi(H'),
]
then local acceptance of every action is insufficient for $\Phi$.
\end{proposition}

\begin{proof}
For every coordinate of both histories,
[
V(a_t)=1.
]
Thus the complete vector of local pass/fail results is
[
(1,\ldots,1)
]
for both $H$ and $H'$. Since $\Phi$ differs between them, $\Phi$ does not factor through the vector of local verification results.
\end{proof}

This proposition identifies the exact requirement for all-pass composition:
[
\Phi
]
must be constant on
[
A_{\mathrm{acc}}^n
]
if the all-pass verdict alone is supposed to establish $\Phi$.

That is a strong condition. It says that every history composed exclusively of locally accepted actions must belong to the same relevant global property class.

\subsection{A Finite-Memory Countermodel}

Local actions can also be identical while their interpretation differs because of prior state.

Let
[
\mathcal S={q_0,q_1}
]
and
[
\mathcal A={a}.
]
There is only one visible action. Define the local verifier by
[
V(a)=1.
]

Let the global property be
[
\Phi(q,a)=
\begin{cases}
1,&q=q_0,\
0,&q=q_1.
\end{cases}
]
If the verifier observes only $a$, then
[
\pi(q_0,a)=a

\pi(q_1,a),
]
while
[
\Phi(q_0,a)\neq\Phi(q_1,a).
]

Thus even perfect verification of the action token cannot establish the property because the property depends upon hidden historical state.

This is the smallest possible state-based counterexample: two latent states and one observable action.

The example demonstrates that output diversity is unnecessary for temporal insufficiency. The same action can be admissible in one history and inadmissible in another.

\subsection{Structural Countermodel}

Let the vertex set be
[
V={1,2,3}.
]
Consider directed graphs
[
G_1=(V,E_1)
]
and
[
G_2=(V,E_2)
]
with
[
E_1={(1,2),(2,3)}
]
and
[
E_2={(1,2),(1,3)}.
]

Suppose all vertices share the same local label
[
\lambda_V(v)=\alpha
]
and all edges share the same local label
[
\lambda_E(e)=\beta.
]

Then the eventwise multisets are identical:
[
M_V(G_1)=M_V(G_2)

{!{\alpha,\alpha,\alpha}!}
]
and
[
M_E(G_1)=M_E(G_2)

{!{\beta,\beta}!}.
]

Define
[
\Phi_{\mathrm{path}}(G)

\mathbf 1
{
\text{$G$ contains a directed path of length two}
}.
]
Then
[
\Phi_{\mathrm{path}}(G_1)=1
]
because
[
1\to2\to3,
]
while
[
\Phi_{\mathrm{path}}(G_2)=0.
]

Hence
[
(M_V(G_1),M_E(G_1))

(M_V(G_2),M_E(G_2))
]
but
[
\Phi_{\mathrm{path}}(G_1)
\neq
\Phi_{\mathrm{path}}(G_2).
]

Therefore no verifier operating solely on the multisets of local vertex and edge attributes can universally determine $\Phi_{\mathrm{path}}$.

Three vertices are minimal for distinguishing the existence of a path of length two.

The same construction generalizes.

\begin{theorem}[Incidence-erasure theorem]
Let $\pi$ be a graph representation invariant under reassignment of edges among fixed vertices while preserving the multiset of vertex attributes and the multiset of edge attributes. Then any graph property $\Phi$ that changes under such a reassignment does not factor through $\pi$.
\end{theorem}

\begin{proof}
By hypothesis there exist graphs $G_1,G_2$ related by an incidence reassignment such that
[
\pi(G_1)=\pi(G_2)
]
while
[
\Phi(G_1)\neq\Phi(G_2).
]
The result follows immediately from the fiber criterion.
\end{proof}

Properties potentially affected include reachability, cycle structure, connectivity, path multiplicity, bottlenecks, centrality, and many forms of workflow dependency.

The theorem does not say that every such property is always lost. It says that a representation invariant under a transformation cannot determine a property that is not invariant under that transformation.

This suggests a general symmetry principle.

\begin{theorem}[Invariance mismatch]
Let a group $\mathcal G$ act on $\mathcal X$. Suppose
[
\pi(gx)=\pi(x)
]
for every
[
g\in\mathcal G
]
and
[
x\in\mathcal X.
]
If $\Phi$ is not invariant under the action, so that for some $g,x$,
[
\Phi(gx)\neq\Phi(x),
]
then $\Phi$ does not factor through $\pi$.
\end{theorem}

\begin{proof}
Choose $g,x$ satisfying
[
\Phi(gx)\neq\Phi(x).
]
By invariance of $\pi$,
[
\pi(gx)=\pi(x).
]
Thus the same representation corresponds to two different $\Phi$-values, violating the fiber criterion.
\end{proof}

This theorem subsumes many structural insufficiency results. Whenever a representation deliberately quotients out a symmetry that the target property does not respect, verification becomes impossible from that representation alone.

\subsection{Representational Countermodel}

Let the source space be
[
X={x_0,x_1}
]
and the verification space
[
Y={y}.
]
Define the transformation
[
\tau(x_0)=\tau(x_1)=y.
]

Let the source property be
[
P_X(x_0)=0,
\qquad
P_X(x_1)=1.
]
Let
[
P_Y(y)=1.
]

A sound verifier on $Y$ accepts $y$:
[
\mathcal V(y)=1.
]
Since
[
P_Y(y)=1,
]
the verifier is perfectly sound with respect to the target representation.

But
[
\mathcal V(\tau(x_0))

\mathcal V(y)

1
]
while
[
P_X(x_0)=0.
]

Thus
[
\mathcal V(\tau(x))=1
\not\Rightarrow
P_X(x)=1.
]

This is a minimal representational countermodel. It contains two source states and one transformed state. The transformation collapses exactly one source distinction, and that distinction is precisely the target property.

The construction can be generalized.

\begin{theorem}[Collapse counterexample]
Let
[
\tau:X\to Y.
]
If $\tau$ is non-injective, choose
[
x_0\neq x_1
]
such that
[
\tau(x_0)=\tau(x_1).
]
Then there exists a Boolean source property
[
P:X\to{0,1}
]
that cannot be universally verified from $\tau(x)$.
\end{theorem}

\begin{proof}
Define
[
P(x_0)=0,
\qquad
P(x_1)=1,
]
and assign arbitrary values elsewhere. Since
[
\tau(x_0)=\tau(x_1)
]
but
[
P(x_0)\neq P(x_1),
]
$P$ does not factor through $\tau$.
\end{proof}

Hence every non-injective representation loses some possible property.

This does not make non-injective representations undesirable. Nearly every useful abstraction is non-injective. The theorem shows instead why sufficiency must always be property-relative.

A representation can be lossy and still be perfectly sufficient for the particular property being verified.

\subsection{Specification Countermodel}

Representational failure can occur even when the implementation is represented exactly if the specification itself collapses distinctions in the intended requirement.

Let the space of informal requirements be
[
I={i_0,i_1}
]
and let the formal-specification space contain one formula
[
F={\varphi}.
]
Suppose
[
\sigma(i_0)=\sigma(i_1)=\varphi.
]

Let
[
W:I\to{0,1}
]
be a property of the intended requirement with
[
W(i_0)=0,
\qquad
W(i_1)=1.
]

A theorem prover can prove
[
\varphi
]
perfectly. It cannot determine which intended requirement generated $\varphi$, because
[
\sigma(i_0)=\sigma(i_1).
]

Thus formal proof of $\varphi$ cannot distinguish
[
W(i_0)
]
from
[
W(i_1).
]

This is the specification analogue of source-translation insufficiency:
[
\text{proof correctness}
\not\Rightarrow
\text{specification adequacy}.
]

The formal theorem can be completely correct while the informal claim attached to it exceeds the distinctions encoded in the specification.

\subsection{Bounded-Verification Countermodel}

Let a transition system contain states
[
s_0,s_1,\ldots,s_{k+1}
]
with deterministic transitions
[
s_i\to s_{i+1}.
]
Suppose
[
P(s_i)=1
]
for
[
0\leq i\leq k
]
but
[
P(s_{k+1})=0.
]

A verifier restricted to executions of length at most $k$ establishes
[
\forall i\leq k,\quad P(s_i)=1.
]
Nevertheless,
[
\forall i\geq0,\quad P(s_i)=1
]
is false.

Thus
[
\text{no counterexample of depth }\leq k
]
does not imply
[
\text{no counterexample}.
]

The inference becomes valid only under an additional completeness result, such as a bound $k^\star$ satisfying
[
\text{if a counterexample exists, one exists with depth }\leq k^\star.
]

Then for
[
k\geq k^\star,
]
bounded verification can establish the unbounded property.

This is another example in which the local verifier need not be weakened. Its result can be exact within the declared horizon. The invalid step is the extension beyond that horizon without a completeness theorem.

\subsection{Epistemic Countermodel I: Prediction and Preference}

Let
[
\mathcal A={a,b}
]
and
[
\mathcal S'={s_a,s_b}.
]
Define a perfect deterministic predictor
[
M(s_a\mid a)=1,
\qquad
M(s_b\mid b)=1.
]

Now define two utility functions:
[
U_1(s_a)=1,
\qquad
U_1(s_b)=0,
]
and
[
U_2(s_a)=0,
\qquad
U_2(s_b)=1.
]

The predictive model $M$ is identical in both decision problems. Yet
[
\arg\max_{c\in{a,b}}
\mathbb E_M[U_1\mid c]

{a},
]
whereas
[
\arg\max_{c\in{a,b}}
\mathbb E_M[U_2\mid c]

{b}.
]

Therefore no function of $M$ alone can determine the preferred action across both decision problems.

The counterexample is minimal: two actions and two outcomes suffice.

\subsection{Epistemic Countermodel II: Identical Ranking, Different Warrant}

Let two evidence states be
[
E_1,E_2.
]
Suppose both produce the same ranking
[
a>b.
]

Under $E_1$, let
[
P(a\text{ is superior}\mid E_1)=0.99.
]
Under $E_2$, let
[
P(a\text{ is superior}\mid E_2)=0.51.
]

Define an intervention-warrant property
[
W(E)

\mathbf 1
{
P(a\text{ is superior}\mid E)\geq0.95
}.
]
Then
[
W(E_1)=1
]
and
[
W(E_2)=0.
]

But the ranking projection satisfies
[
r(E_1)=r(E_2).
]

Therefore warrant does not factor through ranking:
[
W\neq\widehat W\circ r.
]

A ranking can be sufficient for choosing the maximizer conditional on acting while being insufficient for determining whether action is warranted.

\subsection{Epistemic Countermodel III: Expected Value and Tail Risk}

Let actions
[
a,b
]
have losses distributed as follows:
[
L_a=
\begin{cases}
0,&\text{with probability }0.99,\
100,&\text{with probability }0.01,
\end{cases}
]
and
[
L_b=1
]
with probability one.

Then
[
\mathbb E[L_a]=1
]
and
[
\mathbb E[L_b]=1.
]

Thus an expected-loss representation
[
\pi(a)=\mathbb E[L_a]
]
gives
[
\pi(a)=\pi(b).
]

Define instead the worst-case property
[
R(c)=\sup\operatorname{supp}(L_c).
]
Then
[
R(a)=100,
\qquad
R(b)=1.
]

Hence expected loss is insufficient for worst-case risk.

Likewise, for an admissibility threshold
[
\tau=10,
]
define
[
A(c)

\mathbf 1{R(c)\leq10}.
]
Then
[
A(a)=0,
\qquad
A(b)=1,
]
despite
[
\pi(a)=\pi(b).
]

The example demonstrates that equality of expectations does not imply equality under nonlinear risk functionals.

\subsection{Epistemic Countermodel IV: Finite Horizon}

Let two future trajectories be
[
H_1=(0,0,\ldots,0,0,\ldots)
]
and
[
H_2=(0,0,\ldots,0,1,1,\ldots),
]
where the first $k$ coordinates are identical and the trajectories diverge at coordinate $k+1$.

Define the $k$-horizon projection
[
\pi_k(H)

(H_1,\ldots,H_k).
]
Then
[
\pi_k(H_1)=\pi_k(H_2).
]

Define the eventual-event property
[
\Phi(H)

\mathbf 1
{
\exists t>0:H_t=1
}.
]
Then
[
\Phi(H_1)=0,
\qquad
\Phi(H_2)=1.
]

Therefore
[
\Phi
]
does not factor through any fixed finite-horizon projection $\pi_k$ on a trajectory class permitting arbitrary delayed divergence.

This gives a general result.

\begin{theorem}[No fixed finite horizon for eventual properties]
Let $\Hist$ contain, for every $k$, two histories agreeing through time $k$ but differing on a property $\Phi$. Then $\Phi$ does not factor through any finite-horizon projection
[
\pi_k(H)=H_{\leq k}.
]
\end{theorem}

The proof is immediate from the fiber criterion.

The theorem does not imply that the entire infinite history must be retained. A finite sufficient statistic may exist even when no fixed raw horizon is sufficient.

\subsection{Accumulation Countermodel}

Let local contributions satisfy
[
0\leq c_t\leq\varepsilon
]
for every $t$, where
[
\varepsilon>0.
]
Suppose the local verifier accepts exactly when
[
c_t\leq\varepsilon.
]
Thus every contribution passes.

Define the global property
[
\Phi_T

\mathbf 1
\left{
\sum_{t=1}^{T}c_t\leq B
\right}.
]

Choose
[
c_t=\varepsilon
]
for every $t$. Then
[
\sum_{t=1}^{T}c_t

T\varepsilon.
]
For
[
T>\frac{B}{\varepsilon},
]
we obtain
[
T\varepsilon>B.
]
Thus
[
\forall t,\quad V(c_t)=1
]
while
[
\Phi_T=0.
]

The missing composition law is a budget invariant. For example, if one additionally verifies
[
\sum_{j=1}^{t}c_j\leq B
]
at every $t$, then the final cumulative property follows immediately.

Alternatively, if the horizon is known and
[
T\varepsilon\leq B,
]
local bounds imply the global bound.

Thus the same local predicate can be sufficient or insufficient depending upon the composition assumptions.

\subsection{Null-Action Countermodel}

Let the available interventions be
[
\mathcal A={a,b}.
]
Suppose their costs are
[
J(a)=5,
\qquad
J(b)=10.
]
An optimizer restricted to $\mathcal A$ chooses
[
a.
]

Now extend the action set by the null action:
[
\mathcal A^+

{a,b,\varnothing}
]
with
[
J(\varnothing)=0.
]
Then
[
\arg\min_{c\in\mathcal A^+}J(c)

{\varnothing}.
]

Thus
[
a=\arg\min_{c\in\mathcal A}J(c)
]
does not imply that $a$ is preferable to nonintervention.

The example shows that an optimization problem can produce a mathematically correct solution while the feasible set itself omits a decision-relevant alternative.

Verification of optimality is always relative to the declared feasible set.

\subsection{Repair Countermodel}

Let the system possess two binary properties
[
P,Q.
]
Consider an initial state $x$ with
[
(P(x),Q(x))=(0,1).
]
A repair $\rho$ yields
[
(P(\rho(x)),Q(\rho(x)))=(1,0).
]

A verifier observing only $P$ concludes
[
P(x)=0
]
and
[
P(\rho(x))=1.
]
Thus the repair is a strict improvement with respect to $P$.

But if the desired global property is
[
\Phi(x)=P(x)\land Q(x),
]
then
[
\Phi(x)=0
]
and
[
\Phi(\rho(x))=0.
]
The repair has not improved the global property at all.

If instead the objective is
[
J(x)=w_PP(x)+w_QQ(x),
]
then
[
J(\rho(x))-J(x)

w_P-w_Q.
]
The repair improves $J$ only if
[
w_P>w_Q.
]

Thus successful correction of the detected coordinate does not determine whether the intervention improves the full objective.

\subsection{Economic-Aggregation Countermodel}

The institutional example in the conclusion can also be stated as a simple insufficiency theorem.

Let a budget be
[
B=(b_1,b_2)
]
with total expenditure
[
E(B)=b_1+b_2.
]
Consider
[
B_1=(9,1)
]
and
[
B_2=(1,9).
]
Then
[
E(B_1)=E(B_2)=10.
]

Define a distribution-sensitive property
[
D(B)

\mathbf 1{b_1\geq b_2}.
]
Then
[
D(B_1)=1,
\qquad
D(B_2)=0.
]

Therefore
[
D
]
does not factor through aggregate expenditure $E$.

No greater precision in measuring
[
E
]
can recover the distributional distinction between $B_1$ and $B_2$.

More generally, let
[
S:\mathbb R^n\to\mathbb R
]
be an aggregate statistic. If $S$ is non-injective, then there exist distribution-sensitive properties not recoverable from $S$.

This is again a property-relative statement. A scalar aggregate may be completely sufficient for one question and insufficient for another.

\subsection{General Counterexample-Construction Theorem}

The examples above are instances of one construction.

\begin{theorem}[Counterexample construction by fiber splitting]
Let
[
\pi:\mathcal X\to\mathcal O
]
be a representation. Let
[
F_o=\pi^{-1}(o)
]
be a fiber containing at least two distinct elements:
[
|F_o|\geq2.
]
Then there exists a Boolean property
[
\Phi:\mathcal X\to{0,1}
]
for which $\pi$ is insufficient.
\end{theorem}

\begin{proof}
Choose distinct
[
x_0,x_1\in F_o.
]
Define
[
\Phi(x_0)=0,
\qquad
\Phi(x_1)=1,
]
and define $\Phi$ arbitrarily elsewhere. Since
[
\pi(x_0)=\pi(x_1)=o
]
while
[
\Phi(x_0)\neq\Phi(x_1),
]
$\Phi$ does not factor through $\pi$.
\end{proof}

\begin{corollary}
A representation is sufficient for every possible property on $\mathcal X$ if and only if it is injective.
\end{corollary}

\begin{proof}
If $\pi$ is injective, then every $x$ can be recovered from $\pi(x)$ on $\Img(\pi)$, so every property factors through $\pi$.

If $\pi$ is non-injective, the fiber-splitting theorem constructs a property for which it is insufficient.
\end{proof}

This establishes a sharp tradeoff. Universal property preservation requires complete discrimination:
[
\pi(x)=\pi(x')
\Longrightarrow
x=x'.
]
Any genuine abstraction sacrifices universal sufficiency in exchange for compression.

The appropriate question is therefore never whether an abstraction loses information. Unless it is injective, it does. The relevant question is whether it loses information required by the declared property.

\subsection{Counterexample Construction for Local Verification}

Let
[
\mathcal X
]
be a component space and
[
A\subseteq\mathcal X
]
the locally accepted set:
[
A={x:V(x)=1}.
]
For $n$ components, every locally accepted system lies in
[
A^n.
]

\begin{theorem}[Universal local-to-global composition criterion]
For a Boolean global property
[
\Phi:\mathcal X^n\to{0,1},
]
the implication
[
\bigwedge_{i=1}^{n}V(x_i)=1
\Longrightarrow
\Phi(x_1,\ldots,x_n)=1
]
holds universally if and only if
[
A^n\subseteq\Phi^{-1}(1).
]
\end{theorem}

\begin{proof}
The antecedent
[
\bigwedge_iV(x_i)=1
]
holds exactly when
[
(x_1,\ldots,x_n)\in A^n.
]
Thus the implication holds for every system exactly when every element of $A^n$ belongs to $\Phi^{-1}(1)$.
\end{proof}

Consequently, a counterexample can be found by solving the feasibility problem
[
H^\star
\in
A^n\cap\Phi^{-1}(0).
]

If the intersection is nonempty, any member is a witness to non-composition.

This formulation can be useful computationally. Verification of a claimed local-to-global theorem can itself be reduced to checking whether
[
A^n\cap\Phi^{-1}(0)=\varnothing.
]

Thus a composition theorem is equivalent to an emptiness result.

\subsection{Constraint Form}

Suppose local acceptance is expressed by constraints
[
g_i(x_i)\leq0
]
and global failure by
[
h(x_1,\ldots,x_n)>0.
]
A counterexample exists exactly when the constraint system
[
g_i(x_i)\leq0
\qquad
\forall i
]
and
[
h(x_1,\ldots,x_n)>0
]
is feasible.

The verification problem can therefore be written as
[
\text{find }x_1,\ldots,x_n
]
subject to
[
g_i(x_i)\leq0
]
and
[
h(x_1,\ldots,x_n)>0.
]

If no feasible solution exists, then the local constraints imply the global property.

If a feasible solution exists, it is a constructive witness against composition.

This provides a bridge between the abstract theory and SMT, SAT, mixed-integer optimization, model checking, and counterexample-guided verification.

\subsection{Approximate Local Constraints}

Let local quantities satisfy
[
g_i(x_i)\leq\varepsilon_i
]
rather than exact constraints. Suppose a global defect obeys a bound
[
h(x_1,\ldots,x_n)
\leq
F(\varepsilon_1,\ldots,\varepsilon_n).
]

Then a global tolerance
[
h\leq\delta
]
follows whenever
[
F(\varepsilon_1,\ldots,\varepsilon_n)
\leq\delta.
]

The entire composition problem is therefore transferred to the function $F$.

For additive error,
[
F(\varepsilon_1,\ldots,\varepsilon_n)

\sum_{i=1}^{n}\varepsilon_i.
]
For multiplicative amplification,
[
F(\varepsilon_1,\ldots,\varepsilon_n)

\prod_{i=1}^{n}(1+\varepsilon_i)-1.
]
For a contractive process with factor
[
0\leq\lambda<1,
]
one may obtain
[
F
\leq
\sum_{i=1}^{n}\lambda^{n-i}\varepsilon_i.
]

Thus approximate local verification composes only through an explicit error-propagation law.

Without such a law,
[
\forall i,\quad
g_i(x_i)\leq\varepsilon
]
does not specify the size of the global defect.

\subsection{A General Non-Composition Schema}

The countermodels in this appendix can be expressed through a single schema.

Let
[
\pi:\mathcal X\to\mathcal O
]
be the information retained at a checkpoint. Let
[
\Phi:\mathcal X\to\mathcal Y
]
be the property asserted beyond the checkpoint.

To disprove universal validity, it suffices to construct
[
x_0,x_1\in\mathcal X
]
such that
[
\boxed{
\pi(x_0)=\pi(x_1)
\qquad\text{and}\qquad
\Phi(x_0)\neq\Phi(x_1).
}
]

For temporal verification, $x_0$ and $x_1$ are histories.

For structural verification, they are graphs.

For representational verification, they are source objects.

For epistemic verification, they are decision situations or evidence states.

For institutional verification, they can be policy states sharing the same reported indicator.

The mathematical operation is always the same: locate a fiber of the verification representation that intersects more than one target-property class.

The counterexample is not merely evidence that a particular verifier failed. It proves that every verifier restricted to that representation must fail on at least one member of the fiber.

This distinction is central. Empirical verifier failure says
[
\text{this procedure produced the wrong answer}.
]
A fiber-splitting counterexample says
[
\text{no procedure using only this information can answer correctly on all cases}.
]

The second is a stronger diagnosis because it identifies the interface, rather than the intelligence of the verifier, as the limiting factor.

The next appendix extends this deterministic result to probability distributions. There the exact condition
[
\Phi=\widehat{\Phi}\circ\pi
]
will be replaced by conditional uncertainty
[
H(\Phi\mid\pi),
]
Bayes error, mutual information, and approximate sufficiency. This will distinguish universal impossibility from distribution-relative predictability and show formally why excellent empirical accuracy can coexist with structural projection insufficiency.

\section{Probabilistic Sufficiency, Bayes Error, and Information Loss}
\label{app:probabilistic}

The preceding appendices considered universal deterministic verification. A representation
[
\pi:\mathcal X\to\mathcal O
]
was sufficient for a property
[
\Phi:\mathcal X\to\mathcal Y
]
exactly when
[
\Phi=\widehat{\Phi}\circ\pi.
]
That criterion is distribution-free. It requires correctness for every
[
x\in\mathcal X.
]

Most empirical verification systems are evaluated differently. Objects are sampled from a probability distribution, properties may themselves be uncertain, and performance is summarized by expected loss, accuracy, calibration, or other statistical quantities. A representation that is not universally sufficient can nevertheless predict a property with arbitrarily high probability under a particular distribution.

This appendix develops the distinction formally.

Let
[
X:\Omega\to\mathcal X
]
be a random variable on probability space
[
(\Omega,\mathcal F,\mathbb P).
]
Define the observation
[
O=\pi(X)
]
and target
[
Y=\Phi(X).
]

For simplicity, assume initially that $\mathcal X$, $\mathcal O$, and $\mathcal Y$ are finite.

The verification problem is now to infer $Y$ from $O$.

\subsection{Almost-Sure Sufficiency}

The deterministic factorization condition can be weakened from universal equality to almost-sure equality.

\begin{definition}[Almost-sure sufficiency]
A representation $\pi$ is almost surely sufficient for $\Phi$ under distribution $\mathbb P$ if there exists
[
f:\mathcal O\to\mathcal Y
]
such that
[
\mathbb P[f(O)=Y]=1.
]
\end{definition}

This permits failure on source states having probability zero.

\begin{theorem}[Almost-sure fiber criterion]
For finite $\mathcal X$, $\pi$ is almost surely sufficient for $\Phi$ if and only if every positive-probability observational fiber contains only one positive-probability target value.

Equivalently, for every
[
o\in\mathcal O
]
with
[
\mathbb P(O=o)>0,
]
there exists
[
y_o\in\mathcal Y
]
such that
[
\mathbb P(Y=y_o\mid O=o)=1.
]
\end{theorem}

\begin{proof}
Suppose $f(O)=Y$ almost surely. Conditional on any $o$ with positive probability,
[
Y=f(o)
]
almost surely. Thus the conditional distribution of $Y$ is concentrated on one value.

Conversely, suppose every positive-probability observation $o$ determines a unique target value $y_o$ almost surely. Define
[
f(o)=y_o
]
on the support of $O$ and arbitrarily elsewhere. Then
[
f(O)=Y
]
almost surely.
\end{proof}

This is weaker than universal factorization. It is possible that
[
\pi(x_0)=\pi(x_1),
\qquad
\Phi(x_0)\neq\Phi(x_1),
]
while
[
\mathbb P(X=x_1)=0.
]
The deterministic representation is insufficient, but the conflict is invisible under the deployment distribution.

This distinction is fundamental:
[
\boxed{
\text{distribution-relative success}
\neq
\text{universal sufficiency}.
}
]

\subsection{Conditional Entropy}

For discrete variables, exact almost-sure sufficiency has an information-theoretic characterization.

Define the conditional entropy
[
H(Y\mid O)

-\sum_{o,y}
P(o,y)
\log P(y\mid o).
]

\begin{theorem}[Entropy characterization]
For finite $\mathcal Y$,
[
H(Y\mid O)=0
]
if and only if $Y$ is almost surely a deterministic function of $O$.
\end{theorem}

\begin{proof}
Conditional entropy decomposes as
[
H(Y\mid O)

\sum_o
P(o)H(Y\mid O=o).
]
Every term is nonnegative. Therefore
[
H(Y\mid O)=0
]
if and only if
[
H(Y\mid O=o)=0
]
for every $o$ with
[
P(o)>0.
]
A finite discrete random variable has zero entropy exactly when its distribution is concentrated at one value. Hence every positive-probability $o$ determines $Y$ almost surely.
\end{proof}

Thus the probabilistic analogue of exact factorization is
[
\boxed{
H(\Phi(X)\mid\pi(X))=0.
}
]

If
[
H(Y\mid O)>0,
]
then irreducible target uncertainty remains after the observation is known.

No deterministic verifier restricted to $O$ can eliminate that uncertainty.

\subsection{Mutual Information}

Recall
[
I(Y;O)

H(Y)-H(Y\mid O).
]
Therefore exact almost-sure sufficiency implies
[
I(Y;O)=H(Y).
]

All target information is then present in the observation.

Conversely, for finite $Y$,
[
I(Y;O)=H(Y)
]
implies
[
H(Y\mid O)=0,
]
so $O$ is almost surely sufficient.

Hence
[
\boxed{
O\text{ is almost surely sufficient for }Y
\iff
I(Y;O)=H(Y).
}
]

The ratio
[
\eta(O;Y)

\frac{I(Y;O)}{H(Y)}
]
when
[
H(Y)>0
]
can be interpreted as the fraction of target entropy removed by observing $O$.

It satisfies
[
0\leq\eta(O;Y)\leq1.
]

However, $\eta$ should not be interpreted as verification accuracy. Two observation systems can have identical mutual information while producing different decision losses under different loss functions.

Information and decision quality remain distinct.

\subsection{Bayes Error}

For classification under zero-one loss, define a verifier
[
f:\mathcal O\to\mathcal Y.
]
Its error probability is
[
R(f)

\mathbb P[f(O)\neq Y].
]

The optimal achievable error from $O$ is the Bayes error
[
R^\star(O)

\inf_f
\mathbb P[f(O)\neq Y].
]

For finite target spaces,
[
R^\star(O)

1-
\sum_o
P(o)
\max_y P(y\mid o).
]

\begin{theorem}[Bayes classifier]
An optimal verifier is
[
f^\star(o)
\in
\arg\max_y P(Y=y\mid O=o),
]
and its risk is
[
R^\star(O)

\sum_o
P(o)
\left[
1-\max_yP(y\mid o)
\right].
]
\end{theorem}

\begin{proof}
Conditional on $O=o$, predicting $y$ produces conditional error
[
1-P(Y=y\mid O=o).
]
This is minimized by maximizing
[
P(Y=y\mid O=o).
]
Averaging the resulting minimum conditional error over $o$ gives the formula.
\end{proof}

This establishes a representation-dependent lower bound:
[
\boxed{
R(f)\geq R^\star(O)
}
]
for every verifier receiving only $O$.

Increasing model capability can reduce
[
R(f)-R^\star(O),
]
the gap between the implemented verifier and the optimal verifier on that representation.

It cannot reduce
[
R^\star(O)
]
without changing the information available in $O$, changing the target distribution, or changing the task.

This yields a useful decomposition:
[
R(f)

R^\star(O)
+
\left[
R(f)-R^\star(O)
\right].
]

The first term is representation-relative irreducible classification error. The second is verifier suboptimality.

A system can therefore fail for two mathematically different reasons:
[
\boxed{
\text{representation error floor}
+
\text{inference error above the floor}.
}
]

\subsection{Binary Bayes Error Within a Fiber}

For a Boolean target
[
Y\in{0,1},
]
let
[
p_o=P(Y=1\mid O=o).
]
Then the conditional Bayes error is
[
r^\star(o)

\min{p_o,1-p_o}.
]

Therefore
[
R^\star(O)

\sum_o
P(o)\min{p_o,1-p_o}.
]

This directly generalizes the collision lower bound of Appendix~\ref{app:factorization}.

Suppose a single observational fiber contains two source states
[
x_0,x_1
]
with opposite labels and probabilities
[
p_0,p_1.
]
Any classifier receiving only that fiber identity must assign the same prediction distribution to both. Its error contribution is bounded below by
[
\min{p_0,p_1}.
]

If
[
p_1\ll p_0,
]
the Bayes error can be extremely small even though deterministic sufficiency fails.

For example, suppose
[
\pi(x_0)=\pi(x_1)=o,
]
\Phi(x_0)=0,
\qquad
\Phi(x_1)=1,

and
[
P(X=x_0)=0.999999,
\qquad
P(X=x_1)=10^{-6}.
]
The optimal verifier predicts
[
0
]
and achieves
[
R^\star=10^{-6}.
]

Thus
[
99.9999%
]
accuracy is compatible with a mathematically decisive failure of universal factorization.

This is why empirical accuracy alone cannot establish universal support.

\subsection{Approximate Sufficiency}

Exact sufficiency requires
[
R^\star(O)=0.
]
For statistical systems one may instead define an error tolerance.

\begin{definition}[$\varepsilon$-sufficiency]
An observation $O$ is $\varepsilon$-sufficient for target $Y$ under zero-one loss if
[
R^\star(O)\leq\varepsilon.
]
\end{definition}

Equivalently,
[
\inf_fP[f(O)\neq Y]\leq\varepsilon.
]

This definition is explicitly distribution-relative and loss-relative.

A representation can be
[
10^{-6}\text{-sufficient}
]
under one distribution and
[
0.4\text{-sufficient}
]
under another.

Therefore any approximate sufficiency claim should specify at least
[
(\pi,\Phi,P,L,\varepsilon),
]
where $P$ is the relevant distribution and $L$ the relevant loss function.

Without these qualifiers, ``approximately sufficient'' is incomplete.

\subsection{Loss-Relative Sufficiency}

Zero-one classification is only one decision problem.

Let
[
\mathcal D
]
be a decision space and
[
L:\mathcal Y\times\mathcal D\to\mathbb R
]
a loss function.

A decision rule is
[
\delta:\mathcal O\to\mathcal D.
]
Its risk is
[
R_L(\delta)

\mathbb E[L(Y,\delta(O))].
]

Define
[
R_L^\star(O)

\inf_\delta
\mathbb E[L(Y,\delta(O))].
]

Two representations can have similar classification accuracy and radically different risk under $L$.

For asymmetric binary losses, let
[
L(1,0)=c_{\mathrm{FN}},
\qquad
L(0,1)=c_{\mathrm{FP}},
]
with correct classifications having zero loss.

Conditional on observation $o$, predicting $1$ has expected loss
[
c_{\mathrm{FP}}P(Y=0\mid o),
]
while predicting $0$ has expected loss
[
c_{\mathrm{FN}}P(Y=1\mid o).
]

Prediction $1$ is optimal exactly when
[
c_{\mathrm{FP}}P(Y=0\mid o)
\leq
c_{\mathrm{FN}}P(Y=1\mid o).
]

Writing
[
p_o=P(Y=1\mid o),
]
this becomes
[
c_{\mathrm{FP}}(1-p_o)
\leq
c_{\mathrm{FN}}p_o,
]
or
[
p_o
\geq
\frac{c_{\mathrm{FP}}}
{c_{\mathrm{FP}}+c_{\mathrm{FN}}}.
]

Thus the same posterior distribution can imply different optimal actions under different loss structures.

Probability does not determine policy without a loss or decision criterion.

\subsection{Data Processing}

Suppose the verification pipeline is
[
X
\xrightarrow{\pi}
O
\xrightarrow{\rho}
Z.
]
Then
[
Z=\rho(O)
]
is a deterministic post-processing of $O$.

The random variables form a Markov chain
[
Y-X-O-Z.
]

The data-processing inequality gives
[
I(Y;Z)\leq I(Y;O).
]

Thus deterministic post-processing cannot increase the mutual information about $Y$.

More generally, for a stochastic channel
[
K(z\mid o),
]
if
[
Y\to O\to Z
]
forms a Markov chain, then again
[
I(Y;Z)\leq I(Y;O).
]

This gives an information-theoretic form of the earlier statement that downstream reasoning cannot recreate distinctions destroyed upstream.

A downstream verifier may reorganize information, denoise it under distributional assumptions, or make target-relevant structure easier to compute. But if the only information entering the stage is $O$, it cannot create new information about $Y$ beyond what $O$ already contains.

\subsection{Blackwell-Style Decision Ordering}

Mutual information is not the only way to compare representations. A stronger comparison asks whether one observation can simulate another.

Suppose
[
O_2
]
can be generated from
[
O_1
]
through a stochastic channel
[
K(o_2\mid o_1).
]
Then
[
Y\to O_1\to O_2
]
forms a Markov chain.

For every decision problem with loss $L$,
[
R_L^\star(O_1)
\leq
R_L^\star(O_2).
]

\begin{proposition}[Decision monotonicity under garbling]
If $O_2$ is obtained from $O_1$ by a stochastic transformation independent of $Y$ conditional on $O_1$, then no decision problem based on $Y$ has lower Bayes risk using $O_2$ than using $O_1$.
\end{proposition}

\begin{proof}
Any decision rule
[
\delta_2(O_2)
]
can be simulated using $O_1$ by first sampling
[
O_2\sim K(\cdot\mid O_1)
]
and then applying $\delta_2$. Therefore every risk achievable from $O_2$ is also achievable from $O_1$. Taking infima gives
[
R_L^\star(O_1)\leq R_L^\star(O_2).
]
\end{proof}

This is stronger than a statement about one chosen loss function. It says that a garbled representation cannot dominate its source across decision problems.

\subsection{Sufficient Statistics and Verification}

Suppose the underlying source variable is $X$, but the verifier need not reconstruct $X$. It needs only infer
[
Y=\Phi(X).
]

A statistic
[
T=T(X)
]
is target-sufficient when
[
H(Y\mid T)=0
]
in the exact deterministic-target setting.

More generally, if $Y$ is itself stochastic conditional on $X$, one may require
[
Y\perp X\mid T.
]

This means
[
P(Y\mid X,T)

P(Y\mid T).
]

Since $T$ is a function of $X$, this reduces to
[
P(Y\mid X)

P(Y\mid T)
]
almost surely.

Thus $T$ retains all information in $X$ relevant to predicting $Y$.

This probabilistic definition generalizes deterministic factorization.

If
[
Y=\Phi(X)
]
deterministically, then
[
Y\perp X\mid T
]
implies
[
H(Y\mid T)=0.
]

\subsection{Conditional Mutual Information}

The amount of target information present in $X$ but absent from representation $O$ can be quantified by
[
I(Y;X\mid O).
]

If
[
Y=\Phi(X)
]
deterministically, then
[
H(Y\mid X,O)=0,
]
and therefore
[
I(Y;X\mid O)

H(Y\mid O).
]

Hence, in the deterministic-target case,
[
\boxed{
I(Y;X\mid O)=H(Y\mid O).
}
]

This quantity measures the target-relevant uncertainty remaining after observing $O$ that would disappear if the full source $X$ were known.

Exact sufficiency is therefore equivalent to
[
I(Y;X\mid O)=0.
]

This provides a natural measure of representational deficit:
[
D_{\mathrm{info}}(\pi;\Phi)

I(\Phi(X);X\mid\pi(X)).
]

For deterministic $\Phi$,
[
D_{\mathrm{info}}

H(\Phi(X)\mid\pi(X)).
]

The deficit depends upon the deployment distribution. A representation can have zero deficit under one distribution and positive deficit under another.

\subsection{Distribution Shift}

Let
[
P
]
be a training or benchmark distribution and
[
Q
]
a deployment distribution.

A representation may satisfy
[
H_P(Y\mid O)=0
]
while
[
H_Q(Y\mid O)>0.
]

This can occur when a conflicting state exists in the same observational fiber but receives zero probability under $P$ and positive probability under $Q$.

Consider
[
x_0,x_1
]
with
[
\pi(x_0)=\pi(x_1)=o,
]
\Phi(x_0)=0,
\qquad
\Phi(x_1)=1.


Let
[
P(X=x_0)=1,
\qquad
P(X=x_1)=0.
]
Then
[
H_P(Y\mid O)=0.
]

Now let
[
Q(X=x_0)=Q(X=x_1)=\frac12.
]
Under $Q$,
[
P_Q(Y=0\mid O=o)

P_Q(Y=1\mid O=o)

\frac12,
]
and therefore
[
H_Q(Y\mid O=o)=1
]
bit.

The representation has not changed. The support of the deployment distribution has.

This illustrates why universal factorization is stronger than benchmark sufficiency. Universal factorization protects against arbitrary redistribution of probability mass over the same source space:
[
\pi(x)=\pi(x')
\Longrightarrow
\Phi(x)=\Phi(x').
]

If this condition holds, then
[
H_Q(Y\mid O)=0
]
for every distribution $Q$ on $\mathcal X$.

\begin{theorem}[Distribution-robust exact sufficiency]
For finite $\mathcal X$, the following are equivalent:

\begin{enumerate}
\item $\Phi$ factors through $\pi$ universally.
\item For every probability distribution $Q$ on $\mathcal X$,
[
H_Q(\Phi(X)\mid\pi(X))=0.
]
\end{enumerate}
\end{theorem}

\begin{proof}
If $\Phi=\widehat{\Phi}\circ\pi$, then
[
\Phi(X)

\widehat{\Phi}(\pi(X))
]
for every $X$, hence the conditional entropy is zero under every distribution.

Conversely, suppose universal factorization fails. Then there exist
[
x_0,x_1
]
with the same observation and different target values. Define $Q$ by
[
Q(x_0)=Q(x_1)=\frac12.
]
Conditional on their common observation, the target is non-degenerate, so
[
H_Q(Y\mid O)>0.
]
Thus condition 2 fails.
\end{proof}

This theorem gives universal deterministic sufficiency a statistical interpretation: it is exactly sufficiency under every possible redistribution over the source space.

\subsection{Support-Restricted Sufficiency}

Universal robustness can be stronger than necessary when the source space contains physically impossible states.

Let
[
\mathcal R\subseteq\mathcal X
]
be the reachable set.

A representation need only separate property-distinct cases within $\mathcal R$ if states outside $\mathcal R$ are impossible under the system dynamics.

Define the restricted property
[
\Phi|{\mathcal R}
]
and restricted representation
[
\pi|{\mathcal R}.
]

Then the appropriate condition is
[
\pi(x)=\pi(x')
\Longrightarrow
\Phi(x)=\Phi(x')
\qquad
\forall x,x'\in\mathcal R.
]

Equivalently,
[
\Phi|_{\mathcal R}

\widehat{\Phi}
\circ
\pi|_{\mathcal R}.
]

This is stronger than sufficiency under one empirical distribution but weaker than sufficiency over the entire syntactic state space.

It therefore occupies an important intermediate position:
[
\text{benchmark support}
\subseteq
\text{reachable support}
\subseteq
\text{syntactic state space}.
]

Claims should specify which domain is intended.

\subsection{Rare Collisions}

Suppose a representation is almost sufficient except for a set
[
C

\left{
x:
\exists x'
\text{ with }
\pi(x')=\pi(x),
;
\Phi(x')\neq\Phi(x)
\right}.
]

Call $C$ the collision set.

If
[
P(X\in C)=\delta,
]
then all irreducible verification error is concentrated within $C$.

However,
[
R^\star(O)
\leq\delta
]
rather than necessarily equaling $\delta$, because some collision fibers may have highly imbalanced target probabilities.

A simple upper bound follows by choosing any representative label for each fiber. A sharper expression is the Bayes-risk formula.

For binary $Y$,
[
R^\star(O)

\sum_{o\in\mathcal O_C}
P(o)
\min{p_o,1-p_o},
]
where $\mathcal O_C$ denotes ambiguous observation values.

This makes clear why rare ambiguous cases can disappear in aggregate evaluation while remaining structurally important.

If the rare cases correspond to high-consequence events, zero-one error can be a poor metric.

\subsection{Risk-Weighted Projection Defect}

Let
[
c:\mathcal X\to[0,\infty)
]
assign consequence weight to source states.

Define a weighted verification loss
[
L_c(x,\hat y)

c(x)
\mathbf 1{\hat y\neq\Phi(x)}.
]

Then the representation-relative Bayes risk is
[
R_c^\star(O)

\inf_f
\mathbb E
\left[
c(X)
\mathbf 1
{
f(O)\neq\Phi(X)
}
\right].
]

Within a fiber $O=o$, a binary classifier choosing label $1$ incurs conditional weighted loss
[
\mathbb E[
c(X)\mathbf 1{Y=0}
\mid O=o
],
]
while choosing label $0$ incurs
[
\mathbb E[
c(X)\mathbf 1{Y=1}
\mid O=o
].
]

The optimal prediction minimizes these weighted masses rather than ordinary class probability.

Thus a representation with extremely low unweighted Bayes error can have large consequence-weighted Bayes risk.

This formalizes a basic safety concern:
[
\text{rare}
\not\Rightarrow
\text{negligible}.
]

Probability mass and consequence magnitude are separate quantities.

\subsection{Calibration Is Not Sufficiency}

Suppose a verifier outputs
[
q(O)\in[0,1]
]
interpreted as
[
P(Y=1\mid O).
]

Perfect calibration means, informally,
[
P(Y=1\mid q(O)=p)=p
]
for relevant values of $p$.

Calibration does not imply that $O$ is sufficient.

Consider an observation carrying no information:
[
O=o
]
always, with
[
P(Y=1)=0.3.
]
The constant predictor
[
q(o)=0.3
]
is perfectly calibrated.

Yet
[
H(Y\mid O)=H(Y)>0.
]

Thus the verifier can accurately represent its uncertainty while remaining unable to resolve the target.

This is not a defect in calibration. It is exactly what a calibrated system should do when its representation is insufficient.

Consequently,
[
\boxed{
\text{calibration}
\not\Rightarrow
\text{sufficiency}.
}
]

Likewise,
[
\text{sufficiency}
\not\Rightarrow
\text{calibration}
]
for an implemented probabilistic predictor, because a sufficient representation can still be processed by a badly calibrated model.

Representation quality and predictor calibration are distinct.

\subsection{Accuracy Is Not Calibration}

A classifier can also have high accuracy and poor calibration.

Suppose
[
P(Y=1)=0.99
]
and a model always predicts class $1$ with confidence
[
0.51.
]
Its classification accuracy is
[
0.99,
]
but its confidence is not calibrated to the empirical frequency.

Conversely, a calibrated model can have modest classification accuracy when the underlying problem is intrinsically uncertain.

Thus at least three quantities must be distinguished:
[
\text{accuracy},
\qquad
\text{calibration},
\qquad
\text{representational sufficiency}.
]

None is equivalent to either of the others.

\subsection{Selective Verification and Abstention}

Suppose the verifier may abstain:
[
f:\mathcal O
\to
\mathcal Y\cup{\bot}.
]

Define coverage
[
\operatorname{Cov}(f)

P(f(O)\neq\bot)
]
and selective risk
[
R_{\mathrm{sel}}(f)

P(f(O)\neq Y\mid f(O)\neq\bot).
]

A verifier can reduce selective risk by abstaining on ambiguous fibers.

For binary targets, define
[
p_o=P(Y=1\mid O=o).
]
For confidence threshold
[
\alpha\in\left[\frac12,1\right],
]
consider
[
f_\alpha(o)

\begin{cases}
1,&p_o\geq\alpha,\
0,&p_o\leq1-\alpha,\
\bot,&\text{otherwise}.
\end{cases}
]

Then every non-abstained prediction has conditional error at most
[
1-\alpha.
]
Therefore
[
R_{\mathrm{sel}}(f_\alpha)
\leq
1-\alpha.
]

The cost is reduced coverage.

This gives a formal reason to include deferment or abstention in verification architectures. When the representation does not distinguish a fiber sufficiently for a required confidence level, the system need not force a binary decision.

The tradeoff becomes
[
\text{coverage}
\quad\text{versus}\quad
\text{conditional error}.
]

A system lacking abstention implicitly fixes coverage at
[
1,
]
whether or not its representation supports reliable decisions everywhere.

\subsection{Value of Additional Observation}

Let the current observation be
[
O
]
and suppose an additional measurement
[
Z
]
can be acquired.

The original Bayes risk is
[
R_L^\star(O).
]
After observing $Z$, it becomes
[
R_L^\star(O,Z).
]

Since the verifier can always ignore $Z$,
[
R_L^\star(O,Z)
\leq
R_L^\star(O).
]

Define the decision value of the additional observation as
[
\operatorname{VOI}_L(Z\mid O)

R_L^\star(O)

R_L^\star(O,Z).
]

Then
[
\operatorname{VOI}_L(Z\mid O)\geq0.
]

If obtaining $Z$ has cost
[
C(Z),
]
a simple acquisition criterion is
[
\operatorname{VOI}_L(Z\mid O)>C(Z).
]

This formalizes the distinction between acting and sensing. When current support is insufficient, acquiring additional information can dominate committing immediately.

Importantly,
[
I(Y;Z\mid O)>0
]
does not by itself imply positive net decision value after measurement cost. Nor does zero mutual information imply usefulness for every other purpose. Information gain and decision value remain property- and objective-relative.

\subsection{Information Gain from Multiple Checkpoints}

Let
[
O_1,\ldots,O_n
]
be observations obtained at multiple checkpoints. Define
[
O_{\leq n}

(O_1,\ldots,O_n).
]

Conditional entropy is monotone:
[
H(Y\mid O_{\leq n+1})
\leq
H(Y\mid O_{\leq n}).
]

Indeed,
[
H(Y\mid O_{\leq n})

H(Y\mid O_{\leq n+1})

I(Y;O_{n+1}\mid O_{\leq n})
\geq0.
]

Thus retained observations cannot increase information-theoretic uncertainty about $Y$.

But they may add no information:
[
I(Y;O_{n+1}\mid O_{\leq n})=0.
]

Repeated measurements of the same insufficient projection can therefore saturate without reaching sufficiency.

If
[
O_i=\pi(X)
]
deterministically for every $i$, then
[
(O_1,\ldots,O_n)
]
contains no more information than one copy of $\pi(X)$:
[
H(Y\mid O_1,\ldots,O_n)

H(Y\mid O_1).
]

This gives an information-theoretic proof of the earlier redundancy result.

\subsection{Independent Noise and Repeated Observation}

Redundancy can help when repeated observations contain independent noise rather than identical information.

Suppose
[
Y\in{0,1}
]
and each checkpoint observes
[
O_i=Y\oplus N_i,
]
where
[
N_i\sim\operatorname{Bernoulli}(\varepsilon)
]
independently and
[
0<\varepsilon<\frac12.
]

Each observation is imperfect:
[
P(O_i\neq Y)=\varepsilon.
]

For odd $n$, majority vote fails exactly when more than half of the $N_i$ equal $1$:
[
P(\widehat Y_n\neq Y)

\sum_{j=(n+1)/2}^{n}
\binom{n}{j}
\varepsilon^j(1-\varepsilon)^{n-j}.
]

This probability tends to zero as
[
n\to\infty.
]

Thus redundancy can repair observation noise.

Contrast this with structural projection loss. If
[
O_i=\pi(X)
]
for all $i$ and $\pi$ merges property-distinct states, repeated copies do not separate them.

The distinction is:
[
\boxed{
\text{noise can be averaged;}
\qquad
\text{erasure cannot be averaged away}.
}
]

This is a central architectural distinction between unreliable observation and insufficient observation.

\subsection{Correlated Verification Errors}

Redundant verifiers also fail to provide the expected benefit when their errors are strongly correlated.

Let
[
E_i

\mathbf 1{
\widehat Y_i\neq Y
}.
]
If the $E_i$ are independent with
[
P(E_i=1)=\varepsilon<\frac12,
]
majority voting improves reliability rapidly.

But suppose instead
[
E_1=\cdots=E_n
]
almost surely. Then
[
P(\text{majority error})

\varepsilon
]
for every $n$.

Adding verifiers produces no improvement.

Shared representations can create exactly this dependence. If all verifiers operate on the same lossy projection, every member of an ambiguous fiber can induce a common-mode error.

Therefore ensemble size is not by itself a measure of verification independence.

The relevant quantity is the structure of residual error after conditioning on the shared representation.

\subsection{A Decomposition of Ensemble Error}

Suppose verifier $i$ receives
[
O
]
and has internal randomness
[
R_i.
]
Its prediction is
[
\widehat Y_i=f_i(O,R_i).
]

Conditioned on $O=o$, independent verifier randomness can reduce algorithmic variance. But if
[
P(Y\mid O=o)
]
is non-degenerate, uncertainty about $Y$ remains common to all verifiers.

By the law of total variance, for binary $Y$,
[
\operatorname{Var}(Y)

\mathbb E[
\operatorname{Var}(Y\mid O)
]
+
\operatorname{Var}(
\mathbb E[Y\mid O]
).
]

The first term
[
\mathbb E[
\operatorname{Var}(Y\mid O)
]
]
is unresolved target variability within observational fibers.

An ensemble operating only on $O$ cannot eliminate it.

This provides another statistical form of projection insufficiency.

\subsection{Sequential Evidence}

Let
[
E_t=(O_1,\ldots,O_t)
]
be accumulated evidence.

For a fixed target $Y$,
[
H(Y\mid E_{t+1})
\leq
H(Y\mid E_t).
]

Define information gain
[
G_{t+1}

I(Y;O_{t+1}\mid E_t).
]
Then
[
H(Y\mid E_t)

H(Y\mid E_{t+1})

G_{t+1}.
]

Summing,
[
H(Y)

H(Y\mid E_T)

\sum_{t=1}^{T}
I(Y;O_t\mid E_{t-1}).
]

Thus total information acquired about the target decomposes into incremental conditional information gains.

This provides a mathematical model for verification as evidence accumulation rather than a single terminal checkpoint.

A stopping rule can be defined by a required uncertainty threshold:
[
\tau

\inf
\left{
t:
H(Y\mid E_t)\leq\eta
\right}.
]

Alternatively, under decision loss $L$, one can stop when
[
R_L^\star(E_t)
]
falls below an admissible threshold.

The verification boundary then becomes endogenous: it occurs when sufficient evidence has accumulated, not necessarily at a predetermined physical stage.

\subsection{Sequential Probability Ratios}

For binary hypotheses
[
H_0,H_1,
]
let observations
[
O_1,O_2,\ldots
]
have likelihoods
[
p_0(o)

P(O=o\mid H_0)
]
and
[
p_1(o)

P(O=o\mid H_1).
]

Define the cumulative log-likelihood ratio
[
\Lambda_t

\sum_{i=1}^{t}
\log
\frac{p_1(O_i)}{p_0(O_i)}.
]

Evidence accumulates additively in log space.

A sequential verifier can use thresholds
[
A<0<B
]
and choose
[
\begin{cases}
H_0,&\Lambda_t\leq A,\
H_1,&\Lambda_t\geq B,\
\text{continue observing},&A<\Lambda_t<B.
\end{cases}
]

The important structural point is the existence of the third outcome:
[
\text{continue observing}.
]

A checkpoint architecture that requires a terminal binary judgment at every observation discards this possibility.

Sequential verification instead treats uncertainty as a state from which further evidence acquisition is an admissible continuation.

\subsection{Approximate Composition Across Stages}

Consider a chain of properties
[
Y_0,Y_1,\ldots,Y_n
]
where stage $i$ attempts to preserve the truth of the preceding stage.

Let
[
E_i
]
denote the event that preservation fails at stage $i$:
[
E_i

{
Y_i\text{ gives an accepted result while the required predecessor property fails}
}.
]

Suppose
[
P(E_i)\leq\varepsilon_i.
]

Without assumptions about dependence, the union bound gives
[
P\left(
\bigcup_{i=1}^{n}E_i
\right)
\leq
\sum_{i=1}^{n}\varepsilon_i.
]

Therefore the end-to-end failure probability satisfies
[
P(\text{some preservation failure})
\leq
\min
\left{
1,
\sum_{i=1}^{n}\varepsilon_i
\right}.
]

If preservation failures are independent,
[
P(\text{no preservation failure})

\prod_{i=1}^{n}(1-\varepsilon_i),
]
and hence
[
P(\text{at least one failure})

1-
\prod_{i=1}^{n}(1-\varepsilon_i).
]

For small $\varepsilon_i$,
[
1-\prod_i(1-\varepsilon_i)
\approx
\sum_i\varepsilon_i.
]

This illustrates why individually strong stages need not yield equally strong end-to-end guarantees.

For equal stage error
[
\varepsilon_i=\varepsilon,
]
the independent-chain reliability is
[
(1-\varepsilon)^n.
]

As
[
n\to\infty,
]
this tends to zero for every fixed
[
\varepsilon>0.
]

Thus a long pipeline of high-reliability but imperfect transformations can accumulate substantial end-to-end failure probability.

\subsection{Conditional Stage Guarantees}

The preceding calculation can be too crude because stage errors may be conditional.

Suppose stage $i$ satisfies
[
P(E_i\mid E_1^c,\ldots,E_{i-1}^c)
\leq
\varepsilon_i.
]
Then
[
P\left(
\bigcap_{i=1}^{n}E_i^c
\right)

\prod_{i=1}^{n}
P(
E_i^c
\mid
E_1^c,\ldots,E_{i-1}^c
)
]
and therefore
[
P(\text{all stages preserve})
\geq
\prod_{i=1}^{n}(1-\varepsilon_i).
]

Hence
[
P(\text{some failure})
\leq
1-
\prod_{i=1}^{n}(1-\varepsilon_i).
]

This bound does not require unconditional independence. It requires conditional guarantees along the successful path.

This form is often closer to what a compositional verification argument actually needs.

\subsection{Uncertainty Laundering}

Consider a pipeline
[
X
\xrightarrow{T}
Z
\xrightarrow{V}
C,
]
where $V$ is an exact formal verifier.

Suppose
[
C=1
]
if and only if
[
P_Z(Z)=1.
]
Thus
[
P(C=1\mid P_Z(Z)=1)=1.
]

Suppose, however, that the source-reflection property is imperfect:
[
P(
P_X(X)=1
\mid
P_Z(T(X))=1
)

1-\varepsilon.
]

Then observing the exact certificate
[
C=1
]
does not generally establish source correctness with probability one. Under the stated relation,
[
P(P_X(X)=1\mid C=1)

1-\varepsilon
]
when the certificate is equivalent to target-property satisfaction and no additional selection effect intervenes.

The final certificate is exact about $Z$ while the source-level conclusion remains probabilistic.

Representing the entire result merely as
[
\text{VERIFIED}
]
can therefore obscure where uncertainty resides.

The correct decomposition is
[
\underbrace{
P_Z(T(X))=1
}{\text{formally established}}
]
together with
[
\underbrace{
P_X(X)\mid P_Z(T(X))
}{\text{translation/reflection uncertainty}}.
]

Formal certainty at the terminal stage does not retroactively convert probabilistic upstream preservation into logical certainty.

\subsection{Distributional Composition and Covariate Shift}

Suppose a verifier is evaluated under distribution
[
P
]
and deployed under
[
Q.
]
Let its loss be
[
\ell(x).
]
Then
[
R_P

\mathbb E_P[\ell(X)]
]
and
[
R_Q

\mathbb E_Q[\ell(X)].
]

If
[
Q\ll P,
]
then
[
R_Q

\mathbb E_P
\left[
\frac{dQ}{dP}(X)\ell(X)
\right].
]

A small value of
[
R_P
]
does not imply a small value of
[
R_Q
]
unless the density ratio
[
\frac{dQ}{dP}
]
is controlled on high-loss regions.

For bounded loss
[
0\leq\ell\leq1
]
and
[
\frac{dQ}{dP}\leq M,
]
one obtains
[
R_Q
\leq
M R_P.
]

Without such a bound, deployment can concentrate probability mass on rare benchmark failures.

This is especially relevant when representation collisions are rare under $P$ but common under $Q$.

\subsection{Total Variation Bound}

For bounded loss
[
0\leq\ell\leq1,
]
the difference in risk under two distributions satisfies
[
|R_P-R_Q|
\leq
d_{\mathrm{TV}}(P,Q),
]
where
[
d_{\mathrm{TV}}(P,Q)

\sup_A|P(A)-Q(A)|.
]

Thus a known distributional distance can bound degradation.

But this result also clarifies the limitation of benchmark-only claims. If no useful bound on
[
d_{\mathrm{TV}}(P,Q)
]
or another appropriate divergence is available, benchmark risk alone does not determine deployment risk.

\subsection{KL Divergence and Event Transfer}

Let
[
D_{\mathrm{KL}}(Q\Vert P)
]
be the Kullback-Leibler divergence. Pinsker's inequality gives
[
d_{\mathrm{TV}}(P,Q)
\leq
\sqrt{
\frac12
D_{\mathrm{KL}}(Q\Vert P)
}.
]

Therefore, for bounded loss,
[
|R_P-R_Q|
\leq
\sqrt{
\frac12
D_{\mathrm{KL}}(Q\Vert P)
}.
]

This provides one possible route from a distribution-shift assumption to a verification-risk bound.

The important methodological point is that the bridge must be stated. A benchmark result
[
R_P\leq\varepsilon
]
becomes a deployment statement only after a relation between $P$ and $Q$ has been supplied.

Once again, the missing object is a composition law.

\subsection{Approximate Factorization}

Exact factorization requires
[
Y=f(O)
]
almost surely.

An approximate analogue seeks
[
f^\star
\in
\arg\min_f
\mathbb E[
d(Y,f(O))
]
]
for some distortion function
[
d:\mathcal Y\times\mathcal Y\to[0,\infty).
]

Define the representation distortion
[
D^\star(\pi)

\inf_f
\mathbb E[
d(
\Phi(X),
f(\pi(X))
)
].
]

Then exact sufficiency under $d$ occurs when
[
D^\star(\pi)=0.
]

For zero-one distortion,
[
D^\star(\pi)=R^\star(O).
]

This gives a general optimization formulation:
[
\min_{\pi\in\Pi}
C(\pi)
]
subject to
[
D^\star(\pi)\leq\varepsilon.
]

The exact theory in Appendix~\ref{app:factorization} is recovered at
[
\varepsilon=0.
]

Approximate verification can therefore be viewed as a constrained representation-design problem:
[
\boxed{
\text{minimize representation cost subject to bounded target distortion}.
}
]

\subsection{Information-Constrained Verification}

A complementary formulation constrains information rather than representation cost.

Suppose one seeks an encoded representation
[
Z
]
of source $X$ minimizing
[
I(X;Z)
]
subject to a target-performance constraint
[
R_L^\star(Z)\leq\varepsilon.
]

The optimization is
[
\inf_{P(Z\mid X)}
I(X;Z)
]
subject to
[
R_L^\star(Z)\leq\varepsilon.
]

This asks for the least source information that must be retained to support a specified decision quality.

For exact deterministic verification,
[
\varepsilon=0,
]
the representation must at least separate target-distinct source cases on the relevant support.

For approximate verification,
[
\varepsilon>0,
]
some target distinctions can be deliberately merged if their resulting expected loss remains below tolerance.

This gives mathematical form to a tradeoff that otherwise remains qualitative:
[
\text{compression}
\quad\leftrightarrow\quad
\text{verification loss}.
]

\subsection{Rate-Like Lower Bound}

Suppose
[
Y=\Phi(X)
]
and exact recovery is required:
[
H(Y\mid Z)=0.
]
Then
[
I(Y;Z)=H(Y).
]

Since
[
Y\to X\to Z
]
forms a Markov chain,
[
I(Y;Z)\leq I(X;Z).
]
Therefore
[
\boxed{
I(X;Z)\geq H(Y).
}
]

Any exact representation must carry at least the entropy of the target under the deployment distribution.

This lower bound can be weak because the architecture may require more information than the final target label in order to update the verification state prospectively. That issue is addressed by continuation equivalence in the next appendix.

For retrospective one-shot determination, however, the inequality provides a basic information lower bound.

\subsection{Fano-Type Relation}

For a target taking
[
M=|\mathcal Y|
]
possible values, let
[
P_e

P(\widehat Y\neq Y).
]
Fano's inequality gives
[
H(Y\mid O)
\leq
h_2(P_e)
+
P_e\log(M-1),
]
where
[
h_2(p)

-p\log p-(1-p)\log(1-p)
]
is binary entropy.

Equivalently, substantial residual conditional entropy imposes a lower bound on achievable classification error.

For binary $Y$,
[
M=2,
]
so
[
H(Y\mid O)
\leq
h_2(P_e).
]

This links information loss to unavoidable error.

The relationship is not generally an equality. Conditional entropy and Bayes error measure different aspects of unresolved target uncertainty. But Fano-type bounds show that information insufficiency cannot be made irrelevant merely by choosing a more sophisticated classifier.

\subsection{Probability Does Not Supply Warrant}

Finally, let
[
p(a)

P(\text{desired outcome}\mid a,E).
]
Suppose
[
p(a)>p(b).
]
This establishes a probabilistic ordering under the model and evidence.

It does not determine a decision without additional structure.

Let
[
U(a,y)
]
be utility. Then the relevant quantity may be
[
Q(a)

\sum_y
P(y\mid a,E)U(a,y).
]

It is possible that
[
p(a)>p(b)
]
while
[
Q(a)<Q(b).
]

For example, let the desired outcome produce utility $1$, ordinary failure produce utility $0$, but let action $a$ carry an independent catastrophic cost $C$:
[
Q(a)=p(a)-C.
]
If
[
p(a)=0.9,
\qquad
p(b)=0.8,
\qquad
C=0.2,
]
then
[
p(a)>p(b)
]
but
[
Q(a)=0.7<0.8=Q(b).
]

Even a perfectly calibrated probability therefore does not itself determine action.

The probabilistic analogue of the paper's central distinction is:
[
\boxed{
\text{probability}
\neq
\text{warrant}.
}
]

Probability is evidence entering a decision rule. The rule also depends upon loss, utility, constraints, authority, reversibility, and whatever other properties govern admissibility.

\subsection{Summary}

The deterministic condition
[
\Phi=\widehat{\Phi}\circ\pi
]
has several probabilistic counterparts.

Exact almost-sure sufficiency is
[
H(\Phi(X)\mid\pi(X))=0.
]

Residual representational uncertainty is
[
H(\Phi(X)\mid\pi(X)).
]

The best achievable zero-one error from the representation is
[
R^\star(\pi(X)).
]

Approximate sufficiency can be expressed as
[
R^\star(\pi(X))\leq\varepsilon
]
or, under general loss,
[
R_L^\star(\pi(X))\leq\varepsilon.
]

The information deficit for a deterministic target is
[
I(\Phi(X);X\mid\pi(X))

H(\Phi(X)\mid\pi(X)).
]

Post-processing obeys
[
I(Y;Z)\leq I(Y;O),
]
so information erased upstream cannot be recreated by downstream inference alone.

Repeated independent noisy measurements can improve verification, while repeated copies of the same lossy projection cannot.

Distribution shift can turn an almost-sufficient benchmark representation into an insufficient deployment representation by moving probability mass onto previously unseen collisions.

These relations produce a hierarchy:
[
\boxed{
\begin{aligned}
&\text{universal factorization}\
&\qquad\Downarrow\
&\text{almost-sure sufficiency under every distribution}\
&\qquad\Downarrow\
&\text{almost-sure sufficiency under a specified distribution}\
&\qquad\Downarrow\
&\varepsilon\text{-sufficiency under a specified loss and distribution}.
\end{aligned}
}
]

Moving downward weakens the guarantee and introduces additional dependence upon deployment assumptions.

The distinction matters because empirical success can be excellent even when universal factorization fails. A rare collision can yield near-perfect benchmark accuracy. A distribution shift can make the same collision common. A high-accuracy verifier can remain poorly calibrated. A calibrated verifier can remain informationally uncertain. An ensemble can eliminate independent noise while leaving common projection loss untouched.

The statistical question is therefore not merely how often a verifier is correct. It is how much of its remaining error belongs to the verifier and how much belongs to the representation through which the verifier is forced to see the problem.

\section{Continuation Equivalence and Online Verification}
\label{app:continuation}

The factorization results developed above are naturally retrospective. Given a completed object
[
x\in\mathcal X,
]
a representation
[
\pi(x)
]
is sufficient for a property $\Phi$ when the final value $\Phi(x)$ can be recovered from $\pi(x)$. Online verification poses a stronger problem. A system receives a prefix, retains some summary of that prefix, and must later update its judgment as new events arrive.

The relevant question is no longer merely whether two prefixes have the same property value now. It is whether they remain interchangeable under every future continuation relevant to the property.

This distinction produces a stronger equivalence relation.

\subsection{Histories, Prefixes, and Continuations}

Let $\Sigma$ be an event alphabet and let
[
\Sigma^\ast
]
denote the set of finite histories. A history is written
[
h=e_1e_2\cdots e_t.
]
Concatenation is denoted
[
hc.
]

Not every continuation need be possible after every history. Let
[
\mathcal C(h)\subseteq\Sigma^\ast
]
be the set of continuations admissible after $h$.

For two histories $h_1,h_2$, define their common continuation set
[
\mathcal C(h_1,h_2)

\mathcal C(h_1)\cap\mathcal C(h_2).
]

Let
[
\Phi:\Sigma^\ast\to\mathcal Y
]
be a property of completed histories.

Retrospective equivalence is simply
[
h_1\sim_\Phi h_2
\quad\Longleftrightarrow\quad
\Phi(h_1)=\Phi(h_2).
]

For online verification this is generally too coarse.

\begin{definition}[Continuation equivalence]
Two histories $h_1$ and $h_2$ are continuation-equivalent with respect to $\Phi$, written
[
h_1\equiv_\Phi h_2,
]
if
[
\Phi(h_1c)=\Phi(h_2c)
]
for every
[
c\in\mathcal C(h_1,h_2).
]
\end{definition}

Thus two prefixes are continuation-equivalent when no admissible common future can make their property values diverge.

Taking the empty continuation
[
c=\epsilon
]
whenever it is admissible gives
[
h_1\equiv_\Phi h_2
\Longrightarrow
\Phi(h_1)=\Phi(h_2).
]

Hence continuation equivalence is generally finer than present-value equivalence:
[
\equiv_\Phi
;\subseteq;
\sim_\Phi.
]

The converse can fail dramatically.

\subsection{Present Agreement Does Not Imply Future Interchangeability}

Consider a cumulative burden
[
B(h)

\sum_{e\in h}c(e)
]
and threshold
[
K>0.
]
Define
[
\Phi(h)

\mathbf 1{B(h)\leq K}.
]

Let
[
K=10.
]
Suppose
[
B(h_1)=2,
\qquad
B(h_2)=9.
]
Then
[
\Phi(h_1)=\Phi(h_2)=1.
]

Thus
[
h_1\sim_\Phi h_2.
]

Now choose a continuation $c$ with
[
B(c)=2.
]
Then
[
B(h_1c)=4
]
and
[
B(h_2c)=11.
]
Therefore
[
\Phi(h_1c)=1,
\qquad
\Phi(h_2c)=0.
]

Hence
[
h_1\not\equiv_\Phi h_2.
]

A one-bit current verdict
[
\mathbf 1{B(h)\leq K}
]
is sufficient to report whether the threshold has already been exceeded. It is insufficient as the complete online verification state for future updates.

The verifier must preserve at least enough information to distinguish different remaining margins.

\subsection{Prospective Sufficiency}

Let
[
\sigma:\Sigma^\ast\to\mathcal S
]
be an online summary.

\begin{definition}[Prospectively sufficient summary]
A summary $\sigma$ is prospectively sufficient for $\Phi$ if
[
\sigma(h_1)=\sigma(h_2)
\Longrightarrow
h_1\equiv_\Phi h_2.
]
\end{definition}

Thus the partition induced by $\sigma$ must refine continuation equivalence:
[
\sim_\sigma
\subseteq
\equiv_\Phi.
]

This is stronger than retrospective sufficiency:
[
\sim_\sigma
\subseteq
\sim_\Phi.
]

\begin{proposition}
Every prospectively sufficient summary is retrospectively sufficient whenever the empty continuation is admissible after every history.
\end{proposition}

\begin{proof}
If
[
\sigma(h_1)=\sigma(h_2),
]
prospective sufficiency gives
[
h_1\equiv_\Phi h_2.
]
Taking
[
c=\epsilon
]
yields
[
\Phi(h_1)=\Phi(h_2).
]
Therefore $\Phi$ is constant on the fibers of $\sigma$, and hence factors through $\sigma$.
\end{proof}

The converse is false by the cumulative-budget example.

This gives a strict hierarchy:
[
\boxed{
\text{prospective sufficiency}
\Longrightarrow
\text{retrospective sufficiency}.
}
]

\subsection{Canonical Prospective Quotient}

Continuation equivalence induces a quotient
[
q_\Phi^{+}:
\Sigma^\ast
\to
\Sigma^\ast/{\equiv_\Phi},
]
defined by
[
q_\Phi^{+}(h)

[h]{\equiv\Phi}.
]

This is the canonical abstract online verification state.

\begin{theorem}[Canonical prospective state]
Every prospectively sufficient summary $\sigma$ refines the continuation quotient:
[
\sigma\succeq q_\Phi^{+}.
]
\end{theorem}

\begin{proof}
Suppose
[
\sigma(h_1)=\sigma(h_2).
]
Prospective sufficiency implies
[
h_1\equiv_\Phi h_2.
]
Hence
[
q_\Phi^{+}(h_1)

q_\Phi^{+}(h_2).
]
Therefore the map
[
r(\sigma(h))

[h]{\equiv\Phi}
]
is well defined on $\Img(\sigma)$ and satisfies
[
q_\Phi^{+}=r\circ\sigma.
]
Thus
[
\sigma\succeq q_\Phi^{+}.
]
\end{proof}

The quotient
[
\Sigma^\ast/{\equiv_\Phi}
]
is therefore the prospective analogue of the property quotient developed in Appendix~\ref{app:factorization}.

The difference is important. The retrospective quotient asks:
[
\text{Which histories have the same answer now?}
]
The prospective quotient asks:
[
\text{Which histories will always have the same answer under every common future?}
]

The latter can require substantially more information.

\subsection{Recursive Online Summaries}

An online verifier should ideally update its state without reconstructing the complete history.

Let
[
\sigma(h)\in\mathcal S
]
be its current summary. A recursive update rule is a function
[
F:\mathcal S\times\Sigma\to\mathcal S
]
such that
[
\sigma(he)

F(\sigma(h),e).
]

If the final property is recoverable through
[
G:\mathcal S\to\mathcal Y
]
with
[
\Phi(h)=G(\sigma(h)),
]
then the verification architecture is
[
\sigma_{t+1}

F(\sigma_t,e_{t+1}),
]
[
\Phi(H_{\leq t})

G(\sigma_t).
]

This permits exact history-level verification without retaining the complete raw history in the online state.

\begin{theorem}[Recursive verification]
Suppose $\sigma$ satisfies
[
\sigma(he)=F(\sigma(h),e)
]
for every admissible extension and
[
\Phi(h)=G(\sigma(h)).
]
Then $\Phi$ can be verified online using state $\sigma(h)$ without replaying the complete history.
\end{theorem}

\begin{proof}
Initialize
[
\sigma_0=\sigma(\epsilon).
]
After receiving events
[
e_1,\ldots,e_t,
]
compute recursively
[
\sigma_j

F(\sigma_{j-1},e_j).
]
By induction,
[
\sigma_j

\sigma(e_1\cdots e_j).
]
Therefore
[
G(\sigma_t)

\Phi(e_1\cdots e_t).
]
\end{proof}

This theorem shows why "history-level verification'' need not mean "store and repeatedly inspect the entire history.''

The required object can be a sufficient state propagated through time.

\subsection{Recursive State and Continuation Equivalence}

A recursively updated summary is safe for prospective verification only if identical summary states remain behaviorally indistinguishable under future updates.

Suppose
[
\sigma(h_1)=\sigma(h_2).
]
Then for any common next event $e$,
[
\sigma(h_1e)

F(\sigma(h_1),e)

F(\sigma(h_2),e)

\sigma(h_2e).
]

Inductively, for every common continuation
[
c=e_1\cdots e_k,
]
we obtain
[
\sigma(h_1c)=\sigma(h_2c).
]

If
[
\Phi=G\circ\sigma,
]
then
[
\Phi(h_1c)=\Phi(h_2c).
]

Therefore:

\begin{theorem}[Recursive summaries induce continuation equivalence]
If $\sigma$ admits a deterministic recursive update
[
\sigma(he)=F(\sigma(h),e)
]
and
[
\Phi=G\circ\sigma,
]
then
[
\sigma(h_1)=\sigma(h_2)
\Longrightarrow
h_1\equiv_\Phi h_2.
]
\end{theorem}

Thus a deterministic recursive verification state is automatically prospectively sufficient for the property it computes.

The converse requires additional regularity concerning admissible continuation sets, but the quotient by continuation equivalence provides the abstract minimal state from which such recursive constructions can be sought.

\subsection{Right Congruence}

When every event can follow every history,
[
\mathcal C(h)=\Sigma^\ast,
]
continuation equivalence satisfies a useful compatibility law.

\begin{proposition}[Right congruence]
If
[
h_1\equiv_\Phi h_2,
]
then for every
[
e\in\Sigma,
]
[
h_1e\equiv_\Phi h_2e.
]
\end{proposition}

\begin{proof}
Let $c$ be any continuation. Since
[
h_1\equiv_\Phi h_2,
]
the concatenated continuation
[
ec
]
satisfies
[
\Phi(h_1ec)=\Phi(h_2ec).
]
Therefore
[
h_1e\equiv_\Phi h_2e.
]
\end{proof}

This makes the continuation quotient naturally state-like. Define
[
F([h]{\equiv\Phi},e)

[he]{\equiv\Phi}.
]
The right-congruence proposition guarantees that this update is well defined.

Thus the quotient itself carries a canonical transition system:
[
[h]
\xrightarrow{e}
[he].
]

The property value is
[
G([h])=\Phi(h).
]

Consequently, when the continuation quotient is finite, it directly defines a finite-state exact online verifier.

\subsection{Finite-State Verification Theorem}

Let
[
N_\Phi

\left|
\Sigma^\ast/{\equiv_\Phi}
\right|
]
be the number of continuation-equivalence classes.

\begin{theorem}[Finite-state realizability]
If
[
N_\Phi<\infty,
]
then $\Phi$ admits an exact deterministic online verifier with at most
[
N_\Phi
]
internal states.
\end{theorem}

\begin{proof}
Use the continuation-equivalence classes themselves as verifier states:
[
\mathcal S

\Sigma^\ast/{\equiv_\Phi}.
]
The initial state is
[
[\epsilon].
]
The transition rule is
[
F([h],e)=[he],
]
which is well defined by right congruence. Define
[
G([h])=\Phi(h).
]
Since continuation-equivalent histories agree under the empty continuation,
$G$ is well defined. The resulting verifier computes $\Phi$ exactly.
\end{proof}

The converse also holds in a useful form.

\begin{theorem}[State lower bound]
Any exact deterministic online verifier for $\Phi$ must have at least
[
N_\Phi
]
distinguishable states.
\end{theorem}

\begin{proof}
Suppose an exact verifier has state map
[
\sigma:\Sigma^\ast\to\mathcal S.
]
If two histories are not continuation-equivalent,
[
h_1\not\equiv_\Phi h_2,
]
there exists a continuation $c$ such that
[
\Phi(h_1c)\neq\Phi(h_2c).
]
If
[
\sigma(h_1)=\sigma(h_2),
]
deterministic recursive updates under the same continuation $c$ would leave the verifier in identical final states and therefore produce identical outputs, contradicting exactness.

Thus different continuation-equivalence classes require different verifier states. Hence
[
|\mathcal S|\geq N_\Phi.
]
\end{proof}

Combining the two results yields:

[
\boxed{
\text{minimum number of exact deterministic online states}

N_\Phi.
}
]

This is stronger than a retrospective property-class count.

\subsection{Bit Lower Bound}

If an online verifier uses $b$ bits of state, it can represent at most
[
2^b
]
distinct states.

Therefore:

\begin{corollary}[Memory lower bound]
If
[
N_\Phi<\infty,
]
every exact deterministic online verifier requires at least
[
\boxed{
\left\lceil
\log_2 N_\Phi
\right\rceil
}
]
bits of internal state.
\end{corollary}

If
[
N_\Phi=\infty,
]
no finite-state deterministic verifier can implement exact prospective verification over the unrestricted history space.

This does not necessarily imply that the complete history must be stored. An infinite but compact numerical state can sometimes suffice. A cumulative unbounded counter is the simplest example.

\subsection{Cumulative Budget}

Let event costs satisfy
[
c(e)\in\mathbb N
]
and define
[
B(h)

\sum_{e\in h}c(e).
]
For fixed budget
[
K,
]
let
[
\Phi_K(h)

\mathbf 1{B(h)\leq K}.
]

A natural recursive statistic is
[
\sigma(h)

\min{B(h),K+1}.
]

Its update rule is
[
F(b,e)

\min{b+c(e),K+1}.
]

The output is
[
G(b)

\mathbf 1{b\leq K}.
]

Thus the exact online verifier needs only the truncated cumulative burden.

There are at most
[
K+2
]
states:
[
0,1,\ldots,K,K+1,
]
where
[
K+1
]
represents all already-exceeded histories.

The memory requirement is therefore at most
[
\left\lceil
\log_2(K+2)
\right\rceil
]
bits.

Under a sufficiently rich event alphabet, these states are continuation-distinct and the bound is also necessary.

For example, suppose every integer cost
[
0,\ldots,K
]
can occur as a one-event continuation. If
[
0\leq b_1<b_2\leq K,
]
choose continuation cost
[
d=K-b_2+1.
]
Then
[
b_2+d>K,
]
while
[
b_1+d
\leq
b_2-1+K-b_2+1

K.
]
Hence the two burden states are not continuation-equivalent.

The verifier must distinguish remaining margins.

\subsection{Maximum-Risk Statistic}

Suppose each event carries risk score
[
r(e)\in\mathbb R
]
and the global property is
[
\Phi_\tau(h)

\mathbf 1
\left{
\max_{e\in h}r(e)\leq\tau
\right}.
]

Define
[
\sigma(h)

\max_{e\in h}r(e),
]
with
[
\sigma(\epsilon)=-\infty.
]

Then
[
\sigma(he)

\max{\sigma(h),r(e)}.
]

The final property is
[
\Phi_\tau(h)

\mathbf 1{\sigma(h)\leq\tau}.
]

If only the Boolean threshold property matters prospectively and exceeding the threshold is irreversible, an even smaller state suffices:
[
z(h)

\mathbf 1
\left{
\exists e\in h:r(e)>\tau
\right}.
]
Its update is
[
z(he)

z(h)
\lor
\mathbf 1{r(e)>\tau}.
]

Here the current pass/fail bit \emph{is} prospectively sufficient because once failure occurs it cannot be undone, and all currently passing histories respond identically to every future event with respect to the threshold property.

This contrasts with the cumulative-budget case.

The distinction is algebraic. Maximum-threshold verification has an absorbing failure state and does not depend upon the margin below threshold. Cumulative-budget verification does.

\subsection{Monoid Homomorphisms as Verification Summaries}

Many recursively sufficient statistics arise from algebraic aggregation.

Let
[
(M,\odot,e_M)
]
be a monoid and let
[
w:\Sigma\to M
]
map each event into the monoid.

Extend $w$ to histories by
[
\overline w(\epsilon)=e_M
]
and
[
\overline w(e_1\cdots e_t)

w(e_1)\odot\cdots\odot w(e_t).
]

Then
[
\overline w(hc)

\overline w(h)\odot\overline w(c).
]

If
[
\Phi(h)=G(\overline w(h)),
]
then
[
\overline w(h)
]
is an online sufficient statistic.

Examples include:
[
(\mathbb R,+,0)
]
for cumulative burden,
[
(\mathbb R\cup{-\infty},\max,-\infty)
]
for maximum risk, and
[
({0,1},\lor,0)
]
for irreversible violation detection.

This gives a constructive criterion:
[
\boxed{
\text{find an algebraic summary whose operation composes with history concatenation}.
}
]

The verifier need not preserve raw events when the target property factors through such a homomorphic summary.

\subsection{Product Verification States}

Suppose properties
[
\Phi_1,\ldots,\Phi_m
]
possess recursive sufficient states
[
\sigma_i(h)\in S_i.
]

Then the joint state
[
\sigma(h)

(\sigma_1(h),\ldots,\sigma_m(h))
]
is sufficient for the joint property
[
\mathbf\Phi(h)

(\Phi_1(h),\ldots,\Phi_m(h)).
]

If
[
F_i:S_i\times\Sigma\to S_i
]
are the update functions, define
[
F(
(s_1,\ldots,s_m),
e
)

(
F_1(s_1,e),
\ldots,
F_m(s_m,e)
).
]

This gives the crude state bound
[
|S|
\leq
\prod_{i=1}^{m}|S_i|.
]

Therefore
[
\log_2|S|
\leq
\sum_{i=1}^{m}\log_2|S_i|.
]

The product construction is not necessarily minimal. Different property states can share structure, and some product states may be unreachable. Minimizing the joint continuation quotient can yield a smaller verifier.

This distinction parallels the difference between composing independently designed monitors and designing one monitor for the joint property.

\subsection{Irreversible Safety Properties}

A Boolean property $\Phi$ is prefix-closed when
[
\Phi(h)=0
\Longrightarrow
\Phi(hc)=0
]
for every continuation $c$.

Failure is then irreversible.

For such a property, every failed history belongs to the same continuation-equivalence class if all failed continuations remain failed:
[
\Phi(h)=0
\Longrightarrow
[h]{\equiv\Phi}

F_{\mathrm{dead}}
]
provided no other output distinctions are required.

The passing histories may nevertheless require many distinct states because they can possess different vulnerabilities to future events.

Thus prefix closure simplifies the failed region without necessarily simplifying the admissible region.

A verifier that records only
[
\Phi(h)\in{0,1}
]
is prospectively sufficient exactly when all passing histories are continuation-equivalent:
[
\Phi(h_1)=\Phi(h_2)=1
\Longrightarrow
h_1\equiv_\Phi h_2.
]

This condition is substantially stronger than prefix closure.

\subsection{Recoverable Properties}

Now suppose failure can be repaired. Let
[
\Phi(h)=0
]
at some prefix, but there may exist continuation $c$ with
[
\Phi(hc)=1.
]

Then failed histories can themselves require multiple continuation classes.

For example, let the state be an integer balance
[
b(h)\in\mathbb Z
]
and define
[
\Phi(h)

\mathbf 1{b(h)\geq0}.
]
Events may add positive or negative increments.

Histories with balances
[
-1
\qquad\text{and}\qquad
-100
]
both currently fail, but a continuation adding
[
2
]
repairs the first and not the second.

Therefore
[
\Phi(h_1)=\Phi(h_2)=0
]
does not imply
[
h_1\equiv_\Phi h_2.
]

A failed verdict is not sufficient to determine repairability.

This gives a formal distinction among:
[
\text{current validity},
\qquad
\text{distance to validity},
\qquad
\text{recoverability}.
]

A verifier designed only for the first may be incapable of supporting the latter two.

\subsection{Repair Distance}

Let
[
\mathcal C(h)
]
be the admissible continuations after $h$, with continuation cost
[
C(c)\geq0.
]

Define repair distance
[
d_\Phi(h)

\inf
\left{
C(c):
c\in\mathcal C(h),
;
\Phi(hc)=1
\right}.
]

If no repairing continuation exists, define
[
d_\Phi(h)=\infty.
]

Two histories can share the same present verdict but differ in repair distance:
[
\Phi(h_1)=\Phi(h_2)
]
while
[
d_\Phi(h_1)\neq d_\Phi(h_2).
]

Therefore a representation sufficient for current validity need not be sufficient for repair planning.

If repair policy depends on $d_\Phi$, the appropriate verification unit must preserve the corresponding distinctions.

\subsection{Admissible Continuation Sets}

The continuation relation itself may depend upon history.

Let
[
\mathcal C(h)
]
denote the legal, physically reachable, authorized, or otherwise admissible continuations after $h$.

Then two histories can have identical current property values and identical physical state descriptions while differing in their admissible future sets:
[
\mathcal C(h_1)\neq\mathcal C(h_2).
]

If future verification depends upon what continuations remain admissible, the continuation set is property-relevant information.

For example, suppose both histories arrive at the same machine state $s$, but only one has valid authorization for action $a$. A state-only projection
[
\pi(h)=s
]
merges them. Yet
[
a\in\mathcal C(h_1),
\qquad
a\notin\mathcal C(h_2).
]

Thus equality of physical state does not imply equality of continuation structure.

This motivates an enriched online state
[
\sigma(h)

\bigl(
s(h),
\Gamma(h)
\bigr),
]
where
[
\Gamma(h)
]
represents the governing constraints relevant to continuation.

\subsection{Reachability and Admissibility}

Let
[
\mathcal R(h)
]
be the physically or computationally reachable continuations and
[
\mathcal A(h)
\subseteq
\mathcal R(h)
]
the admissible continuations.

Then
[
c\in\mathcal R(h)
]
does not imply
[
c\in\mathcal A(h).
]

An online verifier concerned with admissibility must therefore distinguish histories whenever the admissible continuation sets differ in a way relevant to the target property.

Define
[
h_1\equiv_{\mathcal A}h_2
]
when the two histories induce equivalent admissible future behavior.

A representation sufficient for physical prediction may fail to be sufficient for admissibility:
[
\pi_{\mathrm{phys}}(h_1)

\pi_{\mathrm{phys}}(h_2)
]
while
[
\mathcal A(h_1)
\neq
\mathcal A(h_2).
]

This is the prospective counterpart of the earlier distinction between prediction and warrant.

\subsection{Continuation-Sensitive Authority}

Suppose actors
[
i\in{1,\ldots,n}
]
possess actor-indexed admissible action sets
[
\mathcal A_i(h).
]

A proposed action $a$ can satisfy
[
a\in\mathcal A_i(h)
]
while
[
a\notin\mathcal A_j(h).
]

Thus the property
[
\Phi(h,i,a)

\mathbf 1{a\in\mathcal A_i(h)}
]
depends upon both history and actor.

A verifier receiving only
[
(h,a)
]
cannot determine $\Phi$ if there exist
[
i\neq j
]
with different authority over the same action in the same history.

Likewise, a verifier receiving only
[
(i,a)
]
can be insufficient when authority changes with history.

The sufficient unit must preserve the distinctions upon which standing depends:
[
\pi(h,i,a).
]

This formalizes why evidence that an action is reachable, useful, or valid in the abstract need not establish that a particular actor is authorized to perform it.

\subsection{Evidence State Versus World State}

Let
[
s_t
]
be the physical world state and
[
E_t
]
the evidence available to the verifier.

The complete online state can be represented as
[
z_t=(s_t,E_t).
]

Two histories may produce the same physical state but different evidence:
[
s(h_1)=s(h_2),
\qquad
E(h_1)\neq E(h_2).
]

If the warranted continuation depends upon evidence, then
[
s
]
alone is insufficient.

Conversely, two histories can yield the same evidence while corresponding to different physical states:
[
E(h_1)=E(h_2),
\qquad
s(h_1)\neq s(h_2).
]

The verifier may then face epistemic uncertainty.

Thus:
[
\text{world-state equivalence}
\neq
\text{evidence-state equivalence}.
]

Online verification may need to propagate both.

\subsection{Evidence Updates}

Let new observation
[
o_{t+1}
]
update evidence according to
[
E_{t+1}

U(E_t,o_{t+1}).
]

A decision rule may then be
[
a_t

D(s_t,E_t).
]

If two histories share physical state but differ in $E_t$, the same physical continuation can produce different warranted actions.

This is another example of continuation equivalence being finer than state equivalence.

The sufficient online unit is not necessarily the state of the world. It is the state required by the property and decision problem.

\subsection{Monotone Evidence and Nonmonotone Warrant}

Suppose evidence accumulates:
[
E_t\subseteq E_{t+1}.
]
Informational refinement is monotone.

It does not follow that a decision is monotone.

Let
[
D(E)\in{a,b}.
]
It is possible that
[
D(E_t)=a
]
while
[
D(E_{t+1})=b.
]

Indeed, the purpose of additional evidence is often to permit revision.

Thus
[
E_t\preceq E_{t+1}
]
in informational refinement does not imply
[
D(E_t)\preceq D(E_{t+1})
]
under any natural action ordering.

This prevents a common conflation:
[
\text{more evidence}
\not\Rightarrow
\text{stronger version of the same conclusion}.
]

More evidence can reverse the conclusion.

\subsection{Append-Only Verification Histories}

Let verification events be
[
v_1,\ldots,v_t.
]
An append-only ledger is
[
L_t=(v_1,\ldots,v_t).
]

The current verification state can be computed by a fold:
[
\sigma_t

F^\ast(L_t),
]
where
[
F^\ast
]
is the iterated update function.

If
[
F
]
is deterministic, then
[
\sigma_{t+1}

F(\sigma_t,v_{t+1}).
]

The ledger and the current state serve different purposes.

The state
[
\sigma_t
]
is an operational compression of the past.

The ledger
[
L_t
]
is a retained witness from which the state can potentially be recomputed, audited, or revised under a different update rule.

Hence
[
\text{operational state}
\neq
\text{historical witness}.
]

A system can maintain a compact online verifier while retaining a richer history for retrospective analysis.

\subsection{Replay Consistency}

Suppose the ledger update rule is
[
F.
]
Define replay from initial state $\sigma_0$:
[
\operatorname{Replay}_F(L_t)

F(
\cdots
F(
F(\sigma_0,v_1),
v_2),
\ldots,
v_t
).
]

A basic consistency invariant is
[
\sigma_t

\operatorname{Replay}_F(L_t).
]

If this equality fails, the operational state and retained witness disagree.

This creates a verification property of the verifier itself:
[
\Psi(L_t,\sigma_t)

\mathbf 1
{
\sigma_t

\operatorname{Replay}_F(L_t)
}.
]

Thus append-only history permits a second-order check that cannot be performed from $\sigma_t$ alone.

\subsection{Revision Without Erasure}

Suppose an earlier event
[
v_k
]
is later judged incorrect.

A destructive system may replace or delete it. An append-only system instead adds
[
v_{t+1}

\operatorname{Supersede}(k,v_k').
]

The current-state fold interprets the later event as changing the operative status of the earlier record while retaining the historical fact that the earlier record existed.

Let
[
\operatorname{Active}(L_t)
]
compute the currently governing event set. Then
[
v_k\in L_{t+1}
]
but perhaps
[
v_k\notin\operatorname{Active}(L_{t+1}).
]

Thus
[
\text{retained}
\not\Rightarrow
\text{currently governing}.
]

This is useful whenever verification conclusions are revisable while the history of prior conclusions remains audit-relevant.

\subsection{Exactly-Once Execution and Exactly-Once Disposition}

Let a request possess identifier
[
r.
]
Define execution count
[
N_{\mathrm{exec}}(r)
]
and disposition-witness count
[
N_{\mathrm{disp}}(r).
]

Exactly-once execution requires
[
N_{\mathrm{exec}}(r)=1.
]

Exactly-once disposition witnessing requires
[
N_{\mathrm{disp}}(r)=1,
]
where a disposition may record success, refusal, failure, compensation, cancellation, or another terminally recognized outcome.

These are different properties.

It is possible that
[
N_{\mathrm{exec}}(r)=0
]
while
[
N_{\mathrm{disp}}(r)=1
]
because the request was explicitly refused.

It is possible that
[
N_{\mathrm{exec}}(r)=1
]
while
[
N_{\mathrm{disp}}(r)=0
]
because execution occurred but no terminal witness was durably recorded.

It is also possible that
[
N_{\mathrm{exec}}(r)>1
]
while a deduplicating disposition layer emits one final result.

Thus:
[
\text{execution semantics}
\neq
\text{disposition semantics}.
]

A verifier for one property is not automatically a verifier for the other.

\subsection{Continuation Completeness}

Let
[
\mathcal O
]
be a declared set of terminal outcomes:
[
\mathcal O

{
o_1,\ldots,o_m
}.
]

Define
[
\operatorname{Terminal}(h)
]
when no further operational continuation is expected for the request.

A disposition-completeness property can be written
[
\Phi_{\mathrm{disp}}(h)

\mathbf 1
\left{
\operatorname{Terminal}(h)
\Rightarrow
\exists!o\in\mathcal O:
\operatorname{Witnessed}(o,h)
\right}.
]

The notation
[
\exists!
]
requires exactly one declared disposition witness.

This property cannot be evaluated coherently unless the outcome family
[
\mathcal O
]
has been declared.

Without a declared terminal outcome space, ``complete disposition'' lacks a reference class.

The verification property therefore depends not merely upon observed events but upon the declared semantics of terminal continuation.

\subsection{Compensating Continuation}

Suppose operation
[
a
]
changes state
[
s\to s'.
]
A later discovery may establish that $a$ should not govern the final state.

If direct reversal is impossible, a compensating action
[
a^{-1}_{\mathrm{comp}}
]
may produce
[
s'\to s'',
]
where
[
s''
]
is not necessarily identical to $s$.

Thus compensation is not mathematical inversion:
[
a^{-1}_{\mathrm{comp}}\circ a
\neq
\operatorname{id}
]
in general.

Instead, one may require restoration of selected invariants:
[
I_j(s'')=I_j(s)
]
for
[
j\in J.
]

Compensation therefore requires a declared property set
[
{I_j}_{j\in J}
]
rather than an assumption of complete reversibility.

The relevant verification unit must preserve enough information to determine whether those invariants have been restored.

\subsection{Continuation-Preserving Compression}

Suppose histories are compressed by
[
\sigma:\Sigma^\ast\to S.
]

A particularly strong requirement is
[
\sigma(h_1)=\sigma(h_2)
\Longrightarrow
\mathcal C(h_1)\cong\mathcal C(h_2)
]
with preservation of the relevant property under corresponding continuations.

This can be viewed as a bisimulation-like requirement over verification futures.

For many applications full continuation-structure equivalence is unnecessary. Only $\Phi$-relevant continuation distinctions need be preserved.

Continuation equivalence is therefore a property-relative quotient of future behavior rather than a demand for complete behavioral identity.

This is the prospective analogue of the central principle of the paper:
[
\text{retain the distinctions the property requires, not every distinction the system contains}.
]

\subsection{Prospective Joint Properties}

Let
[
\Phi_1,\ldots,\Phi_m
]
be properties of histories.

Define joint continuation equivalence by
[
h_1
\equiv_{\boldsymbol\Phi}
h_2
]
if for every common continuation $c$,
[
\Phi_j(h_1c)

\Phi_j(h_2c)
]
for all
[
j=1,\ldots,m.
]

Then
[
\equiv_{\boldsymbol\Phi}

\bigcap_{j=1}^{m}
\equiv_{\Phi_j}.
]

Thus adding prospective verification obligations refines the online state requirement exactly as adding retrospective properties refines the static property quotient.

If
[
N_j

|\Sigma^\ast/{\equiv_{\Phi_j}}|,
]
then the naive product construction gives
[
N_{\boldsymbol\Phi}
\leq
\prod_{j=1}^{m}N_j
]
when all quantities are finite.

Consequently,
[
\log_2N_{\boldsymbol\Phi}
\leq
\sum_{j=1}^{m}\log_2N_j.
]

Again, equality need not hold because continuation classes for different properties can share structure.

\subsection{Approximate Continuation Equivalence}

Exact continuation equivalence can be too strong in stochastic environments.

Let
[
P_h^c
]
denote the distribution of a future property value after prefix $h$ under continuation policy or intervention $c$.

One can define
[
h_1
\equiv_{\Phi,\varepsilon}
h_2
]
when
[
d(P_{h_1}^c,P_{h_2}^c)
\leq\varepsilon
]
for every relevant continuation $c$, where $d$ is an appropriate probability metric.

For total variation distance,
[
d_{\mathrm{TV}}
(P_{h_1}^c,P_{h_2}^c)
\leq\varepsilon.
]

An approximate online summary may then merge histories whose future property distributions differ only within tolerance.

The tolerance is necessarily application-relative. If small probability differences concern catastrophic outcomes, total variation alone may not capture decision relevance. A loss-weighted or risk-sensitive metric may be required.

Thus approximate continuation equivalence inherits the same dependence on loss and consequence structure developed in Appendix~\ref{app:probabilistic}.

\subsection{Continuation Error Propagation}

Suppose an approximate summary $\sigma$ merges histories $h_1,h_2$ with future-value discrepancy bounded by
[
\varepsilon.
]

If the system subsequently performs $n$ approximate updates with per-step discrepancy
[
\varepsilon_t,
]
a generic worst-case bound may take the additive form
[
\varepsilon_{\mathrm{tot}}
\leq
\sum_{t=1}^{n}\varepsilon_t.
]

If the update dynamics contract discrepancies by factors
[
\lambda_t<1,
]
one may instead obtain
[
\varepsilon_{\mathrm{tot}}
\leq
\sum_{t=1}^{n}
\left(
\prod_{j=t+1}^{n}\lambda_j
\right)
\varepsilon_t.
]

If the dynamics amplify error,
[
\lambda_t>1,
]
early compression errors can dominate later verification.

Thus the safety of approximate history compression depends not only upon immediate reconstruction error but upon the dynamics of future continuation.

\subsection{Retrospective Minimality Versus Prospective Minimality}

Let
[
N_{\mathrm{ret}}

|\Sigma^\ast/{\sim_\Phi}|
]
and
[
N_{\mathrm{pro}}

|\Sigma^\ast/{\equiv_\Phi}|.
]

Since
[
\equiv_\Phi\subseteq\sim_\Phi,
]
we have
[
N_{\mathrm{pro}}
\geq
N_{\mathrm{ret}}
]
whenever both are finite.

Thus prospective verification can require strictly more state than retrospective classification.

For a Boolean property,
[
N_{\mathrm{ret}}\leq2.
]
Yet
[
N_{\mathrm{pro}}
]
can be arbitrarily large or infinite.

The cumulative-budget example gives
[
N_{\mathrm{ret}}=2
]
but, for integer budget $K$ under a sufficiently rich event alphabet,
[
N_{\mathrm{pro}}=K+2.
]

Therefore the final property label can be a severe compression of the information required to update that label correctly in the future.

This is one of the central mathematical reasons that repeated checkpoint verification can fail. A checkpoint may retain only the current verdict:
[
\text{pass}.
]
But future verification may require the path by which that verdict was reached, or at least a sufficient summary of that path.

\subsection{The Prospective Verification Principle}

The entire appendix can be condensed into one implication.

Let
[
\sigma(h)
]
be the state retained after history $h$.

Exact online verification of $\Phi$ requires
[
\boxed{
\sigma(h_1)=\sigma(h_2)
\Longrightarrow
\forall c\in\mathcal C(h_1,h_2),
\quad
\Phi(h_1c)=\Phi(h_2c).
}
]

A representation that satisfies only
[
\sigma(h_1)=\sigma(h_2)
\Longrightarrow
\Phi(h_1)=\Phi(h_2)
]
is sufficient for the present but not necessarily for the future.

This yields three distinct verification questions:
[
\begin{aligned}
&\textbf{Retrospective:}
&&\text{What property does this completed history have?}\
&\textbf{Online:}
&&\text{What state must be retained to update that property correctly?}\
&\textbf{Prospective:}
&&\text{Which distinctions can affect the property under future continuation?}
\end{aligned}
]

The appropriate verification unit can differ across all three.

For a completed history, the final property value may be enough.

For an evolving history, a sufficient recursive statistic may be required.

For planning, repair, or admissibility, the verifier may additionally need distinctions concerning reachable and permitted futures.

Thus the mathematical unit of verification is not determined only by what the system has done. It can be determined by what differences in the past remain capable of making a difference to the future.

\newpage

\begin{thebibliography}{99}

\bibitem{adelman2014}
Rebecca A. Adelman.
\newblock \emph{Beyond the Checkpoint: Visual Practices in America's Global War on Terror}.
\newblock Amherst, MA: University of Massachusetts Press, 2014.

\bibitem{lahiri2026}
Shuvendu K. Lahiri.
\newblock ``Neuro-Formal Verification: Agentic Language-Agnostic Formal Program Reasoning.''
\newblock \emph{arXiv preprint arXiv:2608.21516}, 2026.
\newblock Version 3, revised September 14, 2026.
\newblock \url{https://arxiv.org/abs/2608.21516}.

\bibitem{ma2026}
Kailang Ma, Heye Huang, Inhi Kim, and Kitae Jang.
\newblock ``Rethinking World Models for Safety-Critical Embodied Systems.''
\newblock \emph{arXiv preprint arXiv:2609.03774}, 2026.
\newblock \url{https://arxiv.org/abs/2609.03774}.

\bibitem{mattei2026}
Clara E. Mattei.
\newblock \emph{Escape from Capitalism: An Intervention}.
\newblock New York: Simon & Schuster, 2026.
\newblock ISBN 978-1-6680-8514-1.

\bibitem{shi2026}
Xinrui Shi, Yanzhe Zhang, and Diyi Yang.
\newblock ``Emergent Collusion in Long-Horizon LLM Agent Interaction.''
\newblock \emph{arXiv preprint arXiv:2609.24967}, 2026.
\newblock \url{https://arxiv.org/abs/2609.24967}.

\bibitem{zhang2026}
Chaoyu Zhang, Hexuan Yu, Heng Jin, Shanghao Shi, Ning Zhang,
Yi Shi, Yulia R. Gel, Y. Thomas Hou, and Wenjing Lou.
\newblock ``Skynet: Workflow-Level Anomaly Detection for Agentic AI via
Semantic and Structural Modeling.''
\newblock \emph{arXiv preprint arXiv:2609.06835}, 2026.
\newblock \url{https://arxiv.org/abs/2609.06835}.

\end{thebibliography}

\end{document}